\documentclass[11pt,fleqn]{report}
\usepackage{graphicx}
\usepackage{amssymb}
\usepackage{float}
\usepackage{longtable}

\usepackage{cml}

\newcommand{\ModelName}{Table Lookup and Interpolation}
\newcommand{\ModelAuthor}{Gary Turner \\ Robert Phillips}
\newcommand{\ModelAffiliation}{Odyssey Space Research}
\newcommand{\ModelDate}{November, 2016}


\begin{document}
\pagestyle{empty}
\titlePage
\clearpage
\pagestyle{plain}
\pagenumbering{roman}
\tableofcontents % Print table of contents
\vfill
{\Large \textbf {Revision History}}
{\large
\begin{center}
\begin{tabular}{|l||l|l|l|}
\hline
 \textbf{Version} & \textbf{Date} & \textbf{Author} & \textbf{Purpose} \\
 \hline \hline
1 & January 2016 & Gary Turner & Initial Version \\ \hline
2 & November 2016 & Gary Turner & Addition of TableIndependent- \\
   & & & Variable\_WrapAround \\ \hline
3 & November 2017 & Gary Turner & Refactor; Major revision \\ \hline
4 & June 2018 & Robert Phillips & Merged TableIndependentVariable \\
   & & & and TableIndependentVariable\_- \\
   & & & WrapAround \\ \hline
5 & June 2018 & Robert Phillips & Code review changes \\ \hline
6 & July 2018 & Gary Turner  & Mathematical Formulation for \\
   & & & Quaternions; added User Guide for \\
   & & & Angles, Quaternions and Transpose\\
   & & & Data; final editing \\ \hline
7 & Sep 2020 & Gary Turner     & Verification and Validation; \\
   & & & Added verification and validation \\
   & & & content \\ \hline
8 & June 2022 & Brian Birmingham & Added data reload capability;  \\
   & & & Added code coverage results \\ \hline
\end{tabular}
\end{center}
}

\chapter{Introduction} % Make a "chapter" heading
\pagenumbering{arabic} % Start text with arabic 1
The Table Interpolation model was written to provide a universal and consistent
approach to the common task of interpolating tables.
It has been identified that multiple table interpolation systems have been
written -- and are being independently maintained -- including
multiple instances of interpolations being handled completely within models.
This redundancy adds to both maintenance and testing costs.
The purpose of this model is to consolidate the table-lookup and
table-interpolation capabilities into one common model that can be used in
a standalone manner or through inclusion in independent models.

The particular features that drove the design of the model are:
\begin{itemize}
 \item Simple things done simply; the simplest table-lookup is with one
 independent variable and one dependent variable with a 1-1 discrete match
 and no interpolation.
 \item Simple extensibility to more complex scenarios, such as interpolation
 on a $n$-dimensional data table covering multiple independent variables
 and multiple dependent variables.
 \item Interdependence of tables; in many cases, an independent variable may be
 used for more than one dependent variable.  Where the same independent
 data set is used for multiple dependent variables, those dependent variables
 share that common independent variable, thereby eliminating duplication
 of both memory space and processing associated with duplicated interpolations
 of identical data sets.
 \item Memory retention of previous lookup.  In many cases, a simulation will
 gradually progress through a lookup table -- e.g. when obtaining
 wind-speed from altitude during a descent.  To make the search more efficient,
 the model
 maintains a memory of its previous location so that it can restart at
 that location.
 \item Bi-directional data.  The model accepts independent data tables with
 values that are increasing or decreasing.  Note, however:
 \begin{itemize}
   \item the independent variable data must be monotonic and each value
         unique, even for a simple lookup.
   \item data representing dependent variables has no such requirement;
         dependent data may be completely random.
 \end{itemize}
 \item Binary search capability.  In cases where the independent variable is
 random or highly variable, the quickest path to find that value in the
 table is often a binary search: divide the table in two, discard one half,
 and repeat on the remaining half until the data point is identified.
 \item Sweep-mode searching.  For cases where the starting point is known,
 sweeping through the adjacent data points testing each one is typically
 the most efficient path to a solution.  This mode can be extremely slow in
 other cases.
 \item Linear and Wrap Around continuity. Independent variables can be
 configured such that their domains are either:
 \begin{itemize}
   \item limited by the specified lower and upper bounds
   \item unbounded by having their values cycle indefinitely, wrapping from the
         lower to upper bound (and vice-versa).
 \end{itemize}
 \item Extensible interpolation framework.  The interpolation framework has
 been written to support additional interpolation options.  Beyond the simple
 arithmetic linear interpolation, the model provides special algorithms for
 interpolating angles and quaternions.  Other interpolation options may be
 added fairly easily.
 \item Modifying existing data.  The model supports biasing and scaling the
 data values provided in the lookup tables. For the SimpleTableLookup
 implementation, the model also supports a complete reassignment of data tables
 mid-sim, including changing the number of elements in the independent data,
 and changing the number of dependent variables.
\end{itemize}


\chapter{Requirements}
\begin{enumerate}
\item  Provide a generic table-lookup capability to handle arbitrary $n$
 independent variables applied to arbitrary $m$ dependent variables.
\item  Provide a simple implementation where $n=1$ and $m=1$.
\item Provide wrap-around capabilities for:
 \begin{enumerate}
 \item independent-variables: when the independent variable value goes off the
 table, wrap the value back to the other end of the table.
 \item dependent-variables: when the dependent data represents angles,
 interpolate through the discontinuity inherent with angular data (e.g. $0^o
 \equiv 360^o$).
 \end{enumerate}
 \item Provide the ability to modify the data during a run.
\end{enumerate}

\chapter{Model Specification} \label{sec:Model_spec}
\section{Code Structure}
The model is sub-divided into multiple concepts and classes:
\begin{enumerate}
 \item \textbf{Table Management}.  There are 3 classes dealing with
  table-management:
 \begin{enumerate}
  \item \textbf{AbstractTableLookup}.

  This class provides the abstract basis for the
   classes \textit{SimpleTableLookup} and the more general
   \textit{TableLookupSet}.  These classes manage the relation between the
   independent and dependent variables.
   The class contains STL-lists of pointers to:
   \begin{itemize}
    \item All of the data-tables, each of which may be differently sized and
     differently dimensioned.
    \item All of the independent variables, with each table being assigned a
         subset of this list.
    \item All of the output variables, with each element on this list being
          assigned to one and only one table.
   \end{itemize}
   Note:  This class is not technically an abstract class, as it has no pure virtual function.  It can be instantiated (see \textit{RUN\_abstract\_table\_lookup} in \textit{SIM\_09\_code\_coverage}), but this is only for testing purposes.  The classes \textit{SimpleTableLookup} and/or \textit{TableLookupSet} should be used instead.
  \item \textbf{SimpleTableLookup}.  This class is a simple implementation of a
   \textit{AbstractTableLookup}.  While it can populate
   an arbitrary number of dependent variables, it has only one
   independent variable, and that independent variable has only one
   interpretation that is applied equally to all dependent variables.
   Behind the scenes, the execution of this class creates a
   \textit{TableIndependentVariable} and a
   \textit{GenericMultiInputTable} (or derivative thereof, e.g.
   \textit{GenericSingleInputTable}, \textit{SingleInputTableVarDeriv}, etc.)
   to hold the data for, respectively, the
   independent variable and the the dependent variable(s).  The behavior of the
   dependent variable is limited to that associated with the specified form of
   \textit{GenericMultiInputTable} thus created.
  \item \textbf{TableLookupSet}.  This class provides the management of
   the general situation where multiple dependent variables each
   depend on a subset of the available independent variables.  The number of
   dependent variables and independent variables is unconstrained.  This option
   supports a wider set of behaviors for dependent variables.
  \item \textbf{TableLookupTransposeDataSet}.  This class is a type of
  \textit{TableLookupSet}; it is used where data is presented in a transpose
  sense to that required for a \textit{TableLookupSet}.
  Typically, data is made
  available for each dependent-variable independently of the other dependent
  variables.  These data may be grouped into one data block, but typically such
  a block provides the data for the first variable, then the data for the
  second variable, and so on --
  ($x_1, x_2, x_3, ..., y_1, y_2, y_3, ..., z_1, z_2, z_3, ...$)

  Sometimes, data is presented differently -- the data block contains the data
  for the first calibration point for all dependent variables, then the data
  for the second calibration point, and so on --
  ($x_1, y_1, z_1, x_2, y_2, z_3, ...$)

  In this case, the \textit{TableLookupTransposeDataSet} can be used to process
  these data into a more conventional format, at which point it functions
  equivalently to a \textit{TableLookupSet}.
 \end{enumerate}

 \item \textbf{Independent Variables}.

    The \textbf{TableIndependentVariable} represents
    independent variables.  Each instance of this class includes the data
    representing the calibrated data values for one independent variable.  A
    complex model instantiation may include many of these instances.

    \textbf{Continuity}

    An independent variable can treat values that are out of bounds -- i.e.,
    less than the minimum calibration value or greater than the maximum
    calibration value -- in one of two ways.
    Note that the minimum and maximum calibration values of an independent
    variable are always the first and last calibration values:
    \begin{itemize}
    \item if the calibration values are monotonically increasing, the smallest
    value is the first value and the largest value is the last value.
    \item if the calibration values are monotonically decreasing, the smallest
    value is the last value and the largest value is the first value.
    \item if the calibration values are not monotonic, the model will fail to
    initialize and the independent variable will not be available for lookup or
    interpolation purposes.
    \end{itemize}

    \begin{itemize}
    \item \textbf{Linear Continuity} If the value of the independent variable
    is out of the domain of the table, the closest tabulated value will be used
    instead of the actual value of the independent variable.
    For example, when the independent variable has a value greater than the
    maximum calibrated value, the lookup is computed as
    though the variable had a value equal to the maximum calibrated value.
    \item \textbf{WrapAround Continuity} If the value of the independent
    variable is is out of bounds, the value used for the lookup is computed by
    wrapping the domain on itself until it contains the value of the
    independent variable (or equivalently, the value of the independent
    variable is shifted in discrete steps equal to the size of the domain until
     it lies within the domain).  To achieve this wrapping, the first and last
    calibration points are treated as equal.
    \end{itemize}

    \textbf{Interpretation of Continuous Variables}

    When the independent variable has value between two calibrated values,
    $s_i$ and $s_{i+1}$, the table can be resolved by:
    \begin{itemize}
     \item \textbf{Interp}.  The output will be interpolated based on the value
      of the independent variable relative to the two adjacent calibration
      points.
      Simple linear interpolation is the only available interpolation option
      except where the dependent variable is a quaternion-set, in which case
      the interpolation uses a \textit{slerp} (spherical linear interpolation)
      algorithm or an angle, in which case periodic constraints are applied.
     \item \textbf{Prev}. The output is discrete, based on the calibrated point
      $s_{i}$.
     \item \textbf{Next}. The output is discrete, based on the calibrated point
      $s_{i+1}$.
     \item \textbf{Floor}. The output is discrete, based on the smaller-value
      of $\{s_{i}, s_{i+1}\}$.
     \item \textbf{Ceil}. The output is discrete, based on the larger-value of
      $\{s_{i}, s_{i+1}\}$.
     \item \textbf{Round}. The output is discrete, based on the value of
      $\{s_{i}, s_{i+1}\}$ that is closest to the value of the independent
      variable.
    \end{itemize}

    Note that for monotonically increasing data, \textit{Prev} and
    \textit{Floor} are equivalent, as are \textit{Next} and \textit{Ceil}.
    However, for monotonically decreasing data, \textit{Prev} is equivalent to
    \textit{Ceil}, while \textit{Next} is equivalent to \textit{Floor}.
    Thus, two of these options may be considered redundant if the structure
    of the independent variable data-set is already known.

  \item \textbf{Lookup Tables for dependent variables}.  There are 3 classes
  that provide the data for the dependent variables:
  \begin{enumerate}
   \item \textbf{GenericMultiInputTable}.  This is the general multi-purpose
    table-lookup that is used in the most general applications.
    This class connects
    a set of independent variables with a set of dependent variables.  Each
    dependent variable in the table is a function of the same set of
    independent variables.

    Note - where dependent variables depend on different combinations of
    independent variables, each combination needs its own
    \textit{GenericMultiInputTable}.  This does not violate the requirement to
    avoid duplication; the table-manager, \textit{TableLookupSet}, can manage
    multiple such tables and multiple independent variables, connecting each
    table with its necessary set of independent variables.

    The class provides:
    \begin{itemize}
     \item A STL-vector of data.  This single-dimensioned data set contains
      all of the data for all of the dependent variables.
     \item A STL-vector of indexing information.  This provides the necessary
      indexing for where the data for each dependent variable resides in the
      overall data vector.
     \item A STL-vector of pointers to the target variables.
      The table will populate these variables directly.
     \item A STL-vector of pairs, with each pair providing a pointer to a
     \textit{TableIndependentVariable} and the instruction of how to
     interpret that \textit{TableIndependentVariable} for the specific
     dependent variable.

     For example, a continuous dependent-variable \textit{x} might
     interpolate the independent-variable \textit{s}, while discrete-variable
     \textit{p} might use a rounding function applied to the same
     independent-variable.
    \end{itemize}

   \item \textbf{GenericSingleInputTable}. This class is a specialization of
   the \textit{GenericMultiInputTable}, limiting the number of independent
   variables to 1. The background contents of this class are identical to those
   of the more-comprehensive multi-input form, the differences are:
     \begin{enumerate}
     \item At initialization, the number of input variables is verified to be
     1.
     \item With each update, the interpolation algorithm -- needing to only
     process through one dimensions -- is marginally more efficient.
     \end{enumerate}

   \item \textbf{SingleInputTableForAngles}.  This class is an extension of
    \textit{GenericSingleInputTable}.  It extends the functionality of
    \textit{GenericSingleInputTable} by providing two paths to interpolation.

    All dependent variables represented by this class are assumed to be
    angles.  Units can be either
    degrees or radians (but must be consistent).  The output of the table
    will be in the range $(-180,180]$ or $(-\pi,\pi]$ (user discretion).
    When interpolating
    between two data-table values -- which, by definition,
    divide the circle into two pieces -- the interpolated value will always
    lie within the smaller piece.  This avoids problems that would otherwise
    occur at the discontinuities -- either at the transition from
    $\pi$ to $-\pi$, or $2\pi$ to $0$.
    \begin{itemize}
     \item As an illustrative example, consider interpolating the mid-point
     between $2\pi - \epsilon$ and $\epsilon$.  This midpoint should be $0$.
     However a simple arithmetic computation associated with linear
      interpolation would
      return

      \begin{equation*}
      \frac{(2\pi - \epsilon) + \epsilon}{2} = \pi
      \end{equation*}

      This class recognizes that $2\pi - \epsilon$ and $- \epsilon$ are
      equivalent; it then restricts the interpolation to the smaller
      ``slice'' of the circle (i.e. between between $- \epsilon$ and
       $\epsilon$), and returns $0$ in this example.
    \end{itemize}

    In general, the architecture developed here could be used for any other
    cyclic variable with a period other than $2\pi$ or 360 degrees.  The value
    \textit{half\_circle} describes the mid-point; for the
    \textit{SingleInputTableForAngles} class it has a
    value of either $\pi$ or $180$, but could be assigned some other value
    by extension of the class.

   \item \textbf{SingleInputTableForQuaternions}.  This class is used
    exclusively for quaternion variables, i.e. a set of 4 values.  All 4
    values are handled together as a single entity.  As with
    \textit{SingleInputTableForAngles}, this class is limited to having a
    single independent variable.

    The distinction of this class is that the interpolation uses a
    \textit{slerp} algorithm instead of a simple linear interpolator.

  \item \textbf{SingleInputTableVarDeriv}. This class uses a pair of variables
  representing some variable and its derivative. Both of these could be
  independently interpolated using a pair of \textit{GenericSingleInputTable}
  instances, but this class allows them to be coupled. A cubic polynomial
  function is
  fit to the ``zeroth-derivative'' variable (equivalently a quadratic function
  is fit to its derivative) using the four end points instead of a pair of
  linear equations being fit to each de-coupled pair of end-points.
  See section \ref{sec:Math_var_deriv} for implementation details.
  \end{enumerate}
  \end{enumerate}



\section{Mathematical Formulation}
Six examples are presented here
\begin{enumerate}
 \item The differences between the two continuity settings in the processing
       of the raw value of an independent variable.
 \item The simple generation of multiple dependent variables where they
 share a common independent variable with common interpretation.
 \item The generation of multiple dependent variables where they
 share a common independent variable but have
 disparate interpretations of that independent variable.
\item The generation of a single dependent variable that is a function of two
independent variables.  Note the extension to multiple independent variables
is equivalent to the previous examples.
\item  The \textit{slerp} interpolation algorithm is presented, as applied to
an instance of a \textit{SingleInputTableForQuaternions}.
\item The derivation is presented for the interpolation of a coupled
variable-derivative pair of variables, being interpolated togetheer using
the \textit{SingleInputTableVarDeriv} capability.
\end{enumerate}

\subsection{Independent Variable}
\label{sec:tab_ind_var}
\subsubsection{Continuity Setting}
 Each independent variable has a set of calibration points
 $\{x_0,...,x_n\}$ that are
 either monotonically increasing or monotonically decreasing.

 Suppose we let:
 \begin{align*}
  x_{min} &= min( x_0, x_n ) \\
  x_{max} &= max( x_0, x_n )
 \end{align*}

 If the input data, $x$, lies inside the specified domain
 (i.e. $x_{min} \leq x \le x_{max}$) then there is no functional
 difference between the two types of \textit{Continuity}.
 The difference becomes apparent only when $x$ lies outside the specified
 domain, e.g., $x > x_{max}$.

 For basic \textit{Linear Continuity} (default behavior), the value $x'$ that
 is used in the lookup is determined by:

 \begin{equation*} \label{eq1}
  x' = \begin{cases}
         x_{min} ~,~ & x \leqslant x_{min} \\
         x       ~,~ & x_{min} \leqslant x \leqslant x_{max} \\
         x_{max} ~,~ & x \geqslant x_{max}
       \end{cases}
 \end{equation*}


 For \textit{WrapAround Continuity}, $x'$ is wrapped if needed so that
 $x_{min} \leq x' < x_{max}$.
 The original input value $x$ will be translated by some integer multiple
 of the specified domain size $(x_{max} - x_{min})$ such
 that it does lie within the sub-domain, and that value will be sent
 to the table look-up.  Essentially, the domain is extended to infinity with
 each point in the original domain being replicated periodically.

 \begin{equation*}
   \forall p \in \mathbb{Z} ~,~ ~~  x \equiv x + p(x_{max} - x_{min})
 \end{equation*}
 \begin{equation*}
   \exists p' \in \mathbb{Z} ~~:~~ \{x' =  x + p(x_{max} - x_{min}), x_{min}
   \leqslant x' < x_{max}\}
 \end{equation*}

 For example, suppose the calibration points are
 $x_i \in \{ -1, 0, 1, 2, 3\}$ and the
 input variable has value $x = 8.2$.
 If that input variable has \textit{Linear} continuity, $x' = 3$.
 If that input variable has \textit{WrapAround} continuity, $x$ will be
 twice-shifted by the size of the domain ($4$) resulting in $x' = 0.2$.

 \textbf{NOTE }An important requirement for correctly using
 \textit{WrapAround} continuity is that
 \begin{equation*}
  f(x_0) = f(x_n)
 \end{equation*}
 The output data tables must reflect this requirement -- the data set
 must have identical values in the first and last entries.
 \textbf{This is not checked for, it is left to the user to verify}.
 If this requirement is not satisfied, the function is mathematically
 indeterminate on the cycle boundaries because, for example,
 while
 \begin {equation*}
  x_n + (x_n - x_0) \equiv x_0 + 2(x_n - x_0) \implies
  f(x_n + (x_n - x_0)) = f(x_0 + 2 (x_n - x_0))
 \end{equation*}
 a simple shift of the independent-variable values would results in dissimilar
 outcomes:
 \begin{align*}
  f(x_n + (x_n - x_0)) &= f(x_n) \\
  f(x_0 + 2 (x_n - x_0)) &= f(x_0) \neq f(x_n)
 \end{align*}
 Algorithmically, the indeterminacy is resolved by making $f(x) = f(x_0)$ when
 $x=x_n$, but this may produce confusing results if the requirement is not met.

 \subsubsection{Index and Fraction}
 The calibration points define a contiguous domain that can be considered as
 being a collection of ``boxes''. For all in-range values, we can identify which
 ``box'' contains the value:
 \begin{equation*}
    \exists i \in ( 0 ... n - 1 ) ~:~ x_i \leqslant x < x_{i+1}
 \end{equation*}
 and a measure of the proximity of the value to $x_{i+1}$:
 \begin{equation*}
    \alpha = \frac{x - x_i}{x_{i+1} - x_i}
 \end{equation*}

 Note
 \begin{align*}
  x=x_i & \implies \alpha = 0.0 \\
  x \to x_{i+1} & \implies \alpha \to 1.0
 \end{align*}

 The index $i$ and interpolation fraction $\alpha$ are used by each table
 that uses this independent variable for interpolating the dependent values.

\subsection{Simple Interpolation of 1 Independent Variable}
In the first example, let $x$ represent the independent variable.
with $\{a, b, c\}$ the dependent variables.
Note that the size of the data table must be equal to the product of:
\begin{itemize}
 \item the number of dependent variables
 \item the size of the independent array
\end{itemize}
In this example, I have 3 dependent variables ( $\{a, b, c\}$) and I will
provide 4 calibrated values for the independent variable, so I need
12 values for the table data vector.
This is checked internally; the model will fail to initialize if this
requirement is not met.

The relationship between $x$ and $a$, $b$, and $c$ is illustrated:
\begin{align}
 \begin{bmatrix}
  x_0\\x_1\\x_2\\x_3
 \end{bmatrix}
:
\begin{bmatrix}
 a_0&b_0&c_0\\
 a_1&b_1&c_1\\
 a_2&b_2&c_2\\
 a_3&b_3&c_3
\end{bmatrix}
\end{align}

Using the index $i$ and interpolation fraction $\alpha$ for $x$ as described in
the previous section, the dependent variable $a$ is generated as:
\begin{equation}
  a ~=~ a_i + \alpha (a_{i+1} - a_i)~=~ a_i(1- \alpha) + a_{i+1} (\alpha)
\end{equation}
with similar expressions for b and c.

\subsection{Disparate Interpretations of 1 Independent Variable}
This section considers the implications of using a single independent variable,
but treating it differently for the generation of values for multiple dependent
variables. We illustrate this concept in an example, continuing
with the same concept as the previous example, but in
this example, dependent variable $b$ is considered discrete while $a$ and $c$
remain continuous.
For generating $b$ in this specific example, we want the independent variable
$x$ to be rounded to the nearest calibrated value.

For this example, it is not possible to use the \textit{SimpleTableLookup}
class because the multiple interpretations required (i.e. $x$ must be used
differently for generating $a$ and $c$ versus for generating $b$) violate the
assumptions of the \textit{SimpleTableLookup} class.

However, the more general \textit{TableLookupSet} can be used with two
instances of \textit{GenericMultiInputTable} (or derivatives thereof),
one for variables $a$ and $c$
(which share a common interpretation of the independent variable) and one for
$b$.

The result for $a$ and $c$ is unchanged, but $b$ is different:
\begin{equation*}
  a ~=~ a_i(1- \alpha) + a_{i+1} (\alpha)
\end{equation*}
\begin{equation*}
  c ~=~ c_i(1- \alpha) + c_{i+1} (\alpha)
\end{equation*}
\begin{equation*}  b = \begin{cases}
       b_i  ~~~  \alpha < 0.5 \\
       b_{i+1}  ~~~  \alpha \geqslant 0.5
      \end{cases}
\end{equation*}

\subsection{Multiple Independent Variables}
In this example, we consider the simplest form of the multi-dimensional
 interpolation: one dependent variable (a) and 2 independent variables (x, y)
\begin{align}
\begin{matrix}
   & \begin{bmatrix}
      y_0 & y_1 & y_2
     \end{bmatrix}
   \\
   &\\
   \begin{bmatrix}
     x_0\\x_1\\x_2\\x_3
   \end{bmatrix}
   &
   \begin{bmatrix}
     a_{00} & a_{01} & a_{02}\\
     a_{10} & a_{11} & a_{12}\\
     a_{20} & a_{21} & a_{22}\\
     a_{30} & a_{31} & a_{32}
   \end{bmatrix}
\end{matrix}
\end{align}

Note that the number of elements in the data vector must be equal to
 the product of:
\begin{itemize}
 \item the number of dependent variables
 \item the size of each of the independent arrays
\end{itemize}

In this example, that is $1 \times 4 \times 3 = 12$

This is checked
 internally; the model will fail to initialize if this requirement is not
 met.

For a second (and subsequent) dependent variable, a similar matrix would be
 stacked on top of the a-matrix in another dimension.  In this example,
 the addition of a second dependent variable
 would make the size of the data vector increase to 24 elements.

In this example, the values x and y are compared against the two independent
 arrays.  Suppose the values lie between the last two elements of the x-array
 and the first two elements of the y-array.  The interpolation fractions
$\alpha$ and $\beta$ are obtained from :
\begin{align*}
 &x = x_2 + \alpha (x_3 - x_2)
 \\
 &y = y_0 + \beta (y_1 - y_0)
\end{align*}

Then the dependent variable(s) is generated as:
\begin{equation}
 a ~=~ a_{20}(1-\alpha)(1-\beta) ~+~ a_{21}(1-\alpha)(\beta)
  ~+~ a_{30}(\alpha)(1-\beta) ~+~ a_{31}(\alpha)(\beta)
\end{equation}

Note – for 3 independent variables, this extends to
\begin{equation}\label{eq:3d_interp}
 a ~=~ a_{200}(1-\alpha)(1-\beta)(1-\gamma) ~+~ a_{201}(1-\alpha)(1-\beta)
(\gamma) ~+~ a_{210}(1-\alpha)(\beta)(1-\gamma) ~+~ ...
\end{equation}
and on into higher dimensions as needed.

\subsection{Interpolation of Quaternions} \label{sec:Math_slerp}
When using a \textit{SingleTableForQuaternions}, instead of using a simple
linear interpolation, we use the public domain Spherical Linear Interpolation
algorithm (\textit{slerp}). It is important to understand the limitations of
this algorithm; the term \textit{spherical linear} refers to a line on a
sphere, or a \textit{great circle}.  \textbf{This algorithm is only accurate
when the quaternions in the table all represent rotations on the same great
circle and do so in a consistent manner}.  In particular:
\begin{itemize}
 \item While each quaternion represents a definite operation, each operation
 may have multiple quaternion representations.  To properly interpolate between
 two quaternions, they must represent rotations or transformations about a
 consistent axis; this is sufficient to ensure that the operations are
 co-planar.
 \item When the quaternions represent rotations along a planar curve with
 irregular curvature (i.e. the rate at which the rotation advances is not a
 constant function of the independent variable), the errors introduced will
 be comparable to applying linear interpolation to a non-linear function.
 The resulting quaternion will still represent a rotation in the same plane,
 but to the wrong orientation.
 \item When the quaternions represent rotations that do not lie in a common
 plane, the algorithm fails completely.  The resulting quaternion may not even
 be a recognizable quaternion (aside from having 4 values).  Verifying that all
 quaternions in the table represent rotations in the same plane is a cumbersome
 operation.  \textbf{Co-planar verification IS NOT performed.  It is assumed
 that all quaternions represent co-planar rotations (or transformations)}.
 Notes:
 \begin {itemize}
  \item The simplest verification that all quaternions are co-planar
  is to evaluate the ratios of the \textbf{vector} components of the
  quaternion.  The vector component can be thought of as the product of a
  scalar with
  a unit vector, the scalar represents the angular rotation and the unit
  vector represents the Eigen-axis of that rotation.  Thus, by taking the
  ratios of values within the vector component of the quaternion, we are
  removing the scalar component and considering only information about
  the orientation of the Eigen-axis.  Because the
  Eigen-axis is shared between co-planar quaternions, the ratios between
  the vector components are also common between co-planar quaternions.
  If the ratios \textit{ x:y, x:z, y:z}, or equivalently
  $\frac{v_0}{v_1} ~,~ \frac {v_0}{v_2} ~,~ \frac{v_1}{v_2}$
  match for all quaternions, then their Eigen-axes match and the
  operations are co-planar.  Note -- it is only necessary to verify two of
  these three ratios; if any two match, the third must also match.
  \item This does not mean that $\frac{v_0}{v_1} = \frac{v_0}{v_2}$, etc.,
  but rather $\frac{v_0}{v_1}$ is the same for all quaternions.
  \item The converse is not necessarily true; failure of this test does not
  imply that quaternions are not co-planar, but is a strong suggestion that the
  algorithm will likely fail.
 \end{itemize}
\end{itemize}

To generate the expression for interpolating quaternions, consider a case in
which the independent variable has identified an index $i$ and
interpolation fraction $\alpha$.  The target
quaternion, $Q_T$ must represent a rotation/transformation that lies between
the two quaternions, $Q_A = Q_i$ and $Q_B = Q_{i+1}$.

Without loss of generality, let us work with left-rotation quaternions;
transformation quaternions are simply the conjugate of their respective
rotation quaternion, and right-quaternions are simply the conjugate of
left-quaternions. Consequently, the math is identical for all 4 types, but
rotations are easier to visualize and left-quaternion notation is easier to
follow because the quaternion then appears as an operator (because it compounds
to the left, so is usually written as being applied on the left of the target
vector).

Consider $Q_A$, $Q_B$, and $Q_T$ as representing rotations applied to a
common ``origin'' vector, $\vec v_0$.

\begin{align*}
 \vec v_A &= Q_A \cdot \vec v_0 \\
 \vec v_B &= Q_B \cdot \vec v_0 \\
 \vec v_T &= Q_T \cdot \vec v_0
\end{align*}

Specifically, we consider $Q_A$ and $Q_B$ to be the ``calibrated'' values
of the quaternion (i.e. the quaternions that appear in the interpolation
tables), and $Q_T$ as the interpolated quaternion, representing a rotation to
a vector $\vec v_T$ in the same plane and between vectors $\vec v_A$ and
$\vec v_B$.
We can think of these as rotations by angle $\theta_i$ about some common
Eigen-axis relative to some common initial vector, with
$\theta_A < \theta_T < \theta_B$. As with the linear
interpolation, we can introduce an interpolation fraction:
\begin{equation*}
 \theta_T = \theta_A + \alpha( \theta_B - \theta_A)
\end{equation*}

We can now define two delta-quaternions, representing the quaternions to
get from A to B and from A to T:

\begin{align*}
 \vec v_B &= Q_\Delta \cdot \vec v_A ~=~ Q_\Delta \cdot  (Q_A \cdot \vec v_0)
  \\
 \vec v_T &= Q_\delta \cdot \vec v_A ~=~ Q_\delta \cdot  (Q_A \cdot \vec v_0)
\end{align*}

The problem can now be broken down into manageable parts.  Recognizing that:
\begin{equation}
 Q_T = Q_\delta Q_A
\end{equation}
it becomes only necessary to identify two functions: $f$ providing the desired
quaternion in terms of the fractional rotation $\alpha$ of the whole rotation
between the two calibrated points $Q_\Delta$, and $g$ providing the rotation
between the two calibrated points in terms of the calibrated rotations
\textit{to} those calibrated points.
\begin{align*}
 Q_\delta & = f(\alpha, Q_\Delta) \\
 Q_\Delta & = g(Q_A, Q_B)
\end{align*}

Consider the quaternion representation of an Eigen-rotation through some angle
$\theta$ about some axis $\hat e$.
\begin{equation*}
 Q = \begin{bmatrix} cos \frac{\theta}{2} \\
                     \left[ sin \frac{\theta}{2} \hat e \right]
     \end{bmatrix}
\end{equation*}

Because we require that the axis $\hat e$ is identical for all quaternions in
the table, the three quaternions of interest ($Q_A ~,~ Q_B ~,~ Q_T$) can be
distinguished by their respective angles ($\theta_A ~,~ \theta_B ~,~ \theta_T$)
.

We can also consider $\theta_\Delta$ and $\theta_\delta$:
\begin{align*}
 \theta_\Delta & = \theta_B - \theta_A \\
 \theta_\delta & = \theta_T - \theta_A \\
\end{align*}

and recognize that they are related by the interpolation fraction $\alpha$:

\begin{align*}
 \theta_\delta & = \alpha \theta_\Delta \\
 \theta_T &= \theta_A + \alpha \theta_\Delta =
             \theta_A + \alpha (\theta_B - \theta_A)
\end{align*}

To avoid the clutter of fractions, let
\begin{equation*}
 \phi_i = \frac {\theta_i}{2}
\end{equation*}

Then
\begin{equation*}
 Q_i = \begin{bmatrix} cos \phi_i \\
                       \left[ sin \phi_i \hat e \right]
     \end{bmatrix}
\end{equation*}

and

\begin{equation*}
 \phi_T = \phi_A + \alpha \phi_\Delta
\end{equation*}

Solving for $\phi_T$
\begin{align*}
 cos( \phi_T) &= cos( \phi_A + \alpha \phi_\Delta)
               = cos( \phi_A) cos(\alpha \phi_\Delta) -
                 sin( \phi_A) sin(\alpha \phi_\Delta) \\
 sin( \phi_T) &= sin( \phi_A + \alpha \phi_\Delta)
               = sin( \phi_A) cos(\alpha \phi_\Delta) +
                 cos( \phi_A) sin(\alpha \phi_\Delta)
\end{align*}

We can apply the same relations to $\phi_B$ towards the goal of
identifying a linear relationship between the trigonometric
functions of $\phi_T$, $\phi_A$, and $\phi_B$:
\begin{align*}
 cos( \phi_B)  &= cos( \phi_A) cos(\phi_\Delta) -
                  sin( \phi_A) sin(\phi_\Delta) \implies
 sin( \phi_A)  =  \frac{ cos( \phi_A) cos(\phi_\Delta) - cos( \phi_B)}
                        { sin(\phi_\Delta)} \\
 sin( \phi_B)  &= sin( \phi_A) cos(\phi_\Delta) +
                  cos( \phi_A) sin(\phi_\Delta) \implies
 cos( \phi_A)  =  \frac{ sin( \phi_B) - sin( \phi_A) cos(\phi_\Delta) }
                        { sin(\phi_\Delta)}
\end{align*}

Thus:
\begin{align*}
 cos( \phi_T) &= cos( \phi_A) cos(\alpha \phi_\Delta) -
                 sin(\alpha \phi_\Delta)
                   \frac{ cos( \phi_A) cos(\phi_\Delta) - cos( \phi_B)}
                        { sin(\phi_\Delta)} \\
              &= \frac {cos( \phi_A) cos(\alpha \phi_\Delta) sin(\phi_\Delta) -
                        sin(\alpha \phi_\Delta)cos( \phi_A) cos(\phi_\Delta) +
                        sin(\alpha \phi_\Delta)cos( \phi_B)}
                       {sin(\phi_\Delta)} \\
              &= \frac {cos( \phi_A)  sin( (1-\alpha) \phi_\Delta) +
                        cos( \phi_B) sin( \alpha \phi_\Delta)}
                       {sin(\phi_\Delta)}
\end{align*}
and similarly,
\begin{align*}
 sin( \phi_T) &= \frac {sin( \phi_A)  sin( (1-\alpha) \phi_\Delta) +
                        sin( \phi_B) sin( \alpha \phi_\Delta)}
                       {sin(\phi_\Delta)}
\end{align*}

The symmetry in these two expressions (i.e.
 $cos( \phi_T) = p cos( \phi_A) + q cos( \phi_B)$ and
 $sin( \phi_T) = p sin( \phi_A) + q sin( \phi_B)$ )
means that we can write $Q_T = p Q_A + q Q_B$:

\begin{equation}\label{eqn:slerp}
 Q_T = \frac { Q_A sin( (1-\alpha) \phi_\Delta) +
               Q_B sin( \alpha \phi_\Delta)}
             { sin(\phi_\Delta)}
\end{equation}

It remains to find $\phi_\Delta$.  For this, consider the real (i.e. scalar)
component of the quaternion.  We defined $Q_\Delta$ to be the quaternion
satisfying $Q_B = Q_\Delta \cdot Q_A$,
therefore $Q_\Delta = Q_B \cdot Q_A^*$.  Hence:
\begin{align}
 cos (\phi_\Delta) &= Re(Q_\Delta) = Re(Q_B \cdot Q_A^*) =
 \sum_{j=1}^4 (Q_A)_j  (Q_B)_j \\
\phi_\Delta &= cos^{-1} \left( \sum_{j=1}^4 (Q_A)_j (Q_B)_j \right)
\end{align}

Note to settle confusion -- the scalar part of the product of 2 quaternions,
A and B, includes the negative sign on the vector dot product because $ii =
jj = kk = -1$, so is $q_{As} q_{Bs} - \vec q_Av \cdot \vec q_Bv$.  In this
case, we multiply $Q_B$ with the conjugate of $Q_A$, $Q_A^*$, with the
conjugation introducing another negative sign to the vector dot-product.
Thus $Re(Q_B \cdot Q_A^*) =  q_{As} q_{Bs} + \vec q_Av \cdot \vec q_Bv$.

\fbox{\begin{minipage}[c]{0.8\linewidth}
\textbf{Aside on More General Case}

A similar analysis can be used to generate the spherical-linear interpolation
between non-co-planar quaternions, although the resulting algorithm is not
included in the model at this time.

The important recognition is that the vectors
($ \vec v_A ~,~ \vec v_T ~,~ \vec v_B$)
are co-planar, so the internal eigen-vectors are equal, i.e.
$\hat e_\Delta = \hat e_\delta$.
With this
observation, the following algorithm outline can be used to compute $Q_T$ in
most cases.
\begin{align*}
Q_\Delta    &= Q_B \cdot Q_A^* = \begin{bmatrix} cos \phi_\Delta \\
                                  \left[ sin \phi_\Delta \hat e_\Delta \right]
                                 \end{bmatrix} \implies (\hat e_\Delta~,~
                                 \phi_\Delta)\\
\phi_\delta &= \alpha \phi_\Delta \\
Q_\delta &= \begin{bmatrix} cos \phi_\delta \\
                                  \left[ sin \phi_\delta \hat e_\Delta \right]
            \end{bmatrix} \\
Q_T &= Q_\delta \cdot Q_A
\end{align*}
\end{minipage}}

\textbf{Limit Case}

If $Q_i \approx Q_{i+1}$, then $\phi_\Delta \approx 0$.
In such proximity,
\begin{itemize}
  \item The difference between spherical and linear interpolation is
        negligible (and linear interpolation is much less expensive to compute).
  \item If $Q_B \cdot Q_A^* > 1$ (caused by a numerical rounding error in the
        summation), $\phi_\Delta = \cos^{-1} ( Q_B \cdot Q_A^*)$ is illegal.
  \item $\sin(\phi_\Delta)$ can be be close to zero, which can introduce
        large errors in the evaluation of equation \ref{eqn:slerp}.
\end{itemize}

The algorithm protects against all of these issues by using linear
interpolation when $Q_B \cdot Q_A^* > 1 - \epsilon$.


\subsection{Coupled Interpolation of a Variable and its Derivative}
\label{sec:Math_var_deriv}
The \textit{SingleInputTableVarDeriv} class provides a mechanism for
coupling the interpolation between a variable and its derivative by fitting
the variable to a cubic polynomial $x = f(t)$ and applying the edge
conditions:
\begin{align*}
x_i            &= f(t_i) \\
x_{i+1}        &= f(t_{i+1}) \\
{\dot x}_i     &= f'(t_i) \\
{\dot x}_{i+1} &= f'(t_{i+1})
\end{align*}

With $f(t) = at^3 + b t^2 + c t + d$ and $f'(t) = 3a t^2 + 2b t + c$ we can
evaluate the cubic coefficients ${a,b,c,d}$ as a function of the calibrated
values for variable and its derivative.

Here, we make an important assumption:

\textbf{ The derivative MUST BE TAKEN  with respect to the variable that is
used as the independent variable in the table interpolation algorithm}

For example, if we are coupling the interpolation between position and
velocity, the table values must use \textbf{time} as the independent
variable because velocity is the \textbf{time}-based derivative of position.

Without loss of generality, we are going to substitute $t$ with
$\tau = t - t_i$, then we have three relationships:

\begin{align*}
\tau_i &= 0 \\
f(\tau) &= a\tau ^3 + b \tau ^2 + c \tau + d \\
\frac{df}{d \tau} &= \frac{df}{dt} = f'
\end{align*}

To simplify the algebra, we also use $T = t_{i+1} - t_i = \tau_{i+1}$. Then the
4 coefficients can be expressed as:

\begin{align*}
d  &=  x_i  \\
c  &=  {\dot x}_i \\
b  &= \frac{1}{T^2} [  3(x_{i+1} - x_i) - T( {\dot x}_{i+1} + 2 {\dot x}_i)] \\
a  &= \frac{1}{T^3} [ -2(x_{i+1} - x_i) + T( {\dot x}_{i+1} +   {\dot x}_i)] \\
\end{align*}

Now set the fraction of the interval, as used in previous discussion on
interpolation: $\alpha = \frac{t-t_i}{t_{i+1} - t_i} = \frac{\tau}{T}$ so $\tau
= \alpha T$.

Now,
\begin{align*}
x(\alpha) & = d + c (\alpha T) + b (\alpha T)^2 + a (\alpha T)^3 \\
{} &  = x_i + \alpha T {\dot x}_i +
              \alpha^2[ 3(x_{i+1} - x_i) - T( {\dot x}_{i+1} + 2 {\dot x}_i)] +
              \alpha^3[-2(x_{i+1} - x_i) + T( {\dot x}_{i+1} +   {\dot x}_i)]\\
{} &=  x_i              (1 -3 \alpha ^2 + 2 \alpha^3) +
       x_{i+1}          (3 \alpha ^2 - 2 \alpha ^3) +
       {\dot x}_i     T (\alpha - 2 \alpha^2 + \alpha ^3) +
       {\dot x}_{i+1} T ( - \alpha^2 + \alpha ^3) \\
{} &=  x_i              (1 - \alpha)^2(1 + 2 \alpha) +
       x_{i+1}          \alpha ^2 (3 - 2 \alpha) +
       {\dot x}_i     T \alpha (1 - \alpha)^2 +
       {\dot x}_{i+1} T ( - \alpha^2 (1 - \alpha)) \\
\end{align*}


Conssequently, instead of the decoupled simple linear interpolation form:
\begin{align*}
x(\alpha) &= x_i (1-\alpha) + x_{i+1}(\alpha) \\
{\dot x}(\alpha) &= {\dot x}_i (1-\alpha) + {\dot x}_{i+1}(\alpha)
\end{align*}
we have the coupled forms:
\begin{align}
x(\alpha) & =  x_i  (1 - \alpha)^2(1 + 2 \alpha) +
             x_{i+1} (\alpha ^2 (3 - 2 \alpha)) +
             {\dot x}_i (    T \alpha (1 - \alpha)^2) +
             {\dot x}_{i+1} (T ( - \alpha^2 (1 - \alpha))) \\
{\dot x}(\alpha) & = \left(\frac{x_{i+1}-x_i}{T} \right) (6 \alpha (1 - \alpha)) +
                   {\dot x}_i            (1- \alpha) (1 - 3 \alpha) +
                   {\dot x}_{i+1}        (\alpha (3 \alpha - 2))
\end{align}

Note -- throughout this derivation, we have presented $x$ as a scalar quantity.
The derivation, and the the model implementation, applies equally well to a
vector variable with its vector derivative. The requirements are that the
two variables:
\begin{itemize}
\item have equal dimensions
\item have index-matching, such that the derivative of the nth element of
      the vector representing the variable is represented by the nth element of
      the vector representing the derivative.
\item each dimension can be treated independently of the others.
\end{itemize}


\chapter{User's Guide}
The User's Guide is divided into 8 parts to cover:
\begin{itemize}
 \item The simple table lookup with a single input variable:
 \begin{itemize}
  \item single output
  \item multiple outputs
  \item special case of angular outputs
 \end{itemize}
 \item A table utilizing multiple inputs:
  \begin{itemize}
  \item single output
  \item multiple outputs
 \end{itemize}
 \item A set of tables, each of which may have one or more unique or shared
       inputs.
 \item A set of multiple configurations, with each configuration having
       the potential for multiple tables.
 \item Creating new tables dynamically (a.k.a ``on-the-fly'', ``at runtime'')
\end{itemize}

\section{General Instructions for Table Management} \label{sec:UG_tab_mgr}
The \textit{AbstractTableLookup} class (and thus the classes
 \textit{SimpleTableLookup}, \textit{TableLookupSet}, and
 \textit{TableLookupTransposeDataSet} that inherit from it)
inherit activity-management from the \textit{SubscriptionBase} model.

The \textit{subscribe()} and \textit{unsubscribe()} methods can be used to
manage the activity of the table.  Note that \textit{subscribe()} will
always turn a table on, while \textit{unsubscribe()} will turn it off
\textit{if and only if}
doing so removes the last subscription to the table. Because unsubscribing
doesn't guarantee that table updates will stop, it is strongly recommended
that the user follow a call to \textit{unsubscribe} with a call to
\textit{is\_active()} to determine whether the table is still active due to
another subscriber.
See the \textit{SubscriptionBase} model for more details on this
(inherited) feature.

\textbf{Important Notes:}
\begin{itemize}
 \item the \textit{active} flag is protected.  It can \textit{only} be modified
       via the \textit{subscribe()} / \textit{unsubscribe()} methods;
       it cannot be modified by direct external assignment.
 \item the \textit{active} flag defaults to false.  To use a table, it must
       be deliberately subscribed.  This is often conducted in the
       constructor of the model using the \textit{AbstractTableLookup} (or,
       more correctly, one of its derivatives).
       Alternatively, to support cases where the table is intermittently used,
       the \textit{subscribe()} / \textit{unsubscribe()} calls may be enacted at
       some particular phase of the simulation.
\end{itemize}

\section{Simple Single-input Single-output} \label{sec:UG_single_single}
  For situations in which a single input is used as the only source,
  and the dependent variables all have identical interpretations of that
  variable, the \textit{SimpleTableLookup} class can be used.
  The \textit{SimpleTableLookup} class inherits basic capability from the
  \textit{AbstractTableLookup} class, and implements the most basic
  table-lookup capabilities.
In this section, we consider the simplest of all cases - the single-output
 version of \textit{SimpleTableLookup}.

\subsection{Instantiate}

Create an instance of SimpleTableLookup.  This takes no construction arguments.

\begin{code}
class MyModel
{
 public:
%*\begin{codecomment}
Create a single independent variable and a single dependent variable
\end{codecomment}*)
  double             my_independent_variable;
  double             my_dependent_variable;
  SimpleTableLookup  table_lookup;
  ...
%*\begin{codecomment}
Constructors:
\end{codecomment}*)
  MyModel()
    :
    table_lookup()
  {};
\end{code}


\subsection{Populate}

\subsubsection{Independent Variable Data}
\label{sec:UG_single_single_populate_indep}

Create data for the independent variable and load those into the table.
There are two options for performing this task, using an array of data, or a
 STL-vector of data:
\begin{itemize}

 \item \textbf{Array}

\begin{code}
double scratch_indep[5] = {1,2,3,5,8};
table_lookup.load_independent_data( my_independent_variable,
                                    scratch_indep,
                                    5);
\end{code}

The arguments are:
\begin{enumerate}
 \item A reference to the independent variable itself
 \item A pointer to the first element of the data array.
  For 1-dimensional arrays, the variable name itself serves as that pointer.
  This could also be specified as the address of the first element of the
  array; $\&scratch\_indep[0]$ is equivalent to $scratch\_indep$.
 \item The size of the data array.  In this example that is $5$
 because \textit{scratch\_indep} has five elements.
\end{enumerate}


 \item \textbf{STL-vector}

There are multiple methods by which a STL-vector can be constructed,
 one of the more deliberate methods is illustrated below.
 If generating the data from scratch, the array mechanism is probably the
 easier interface.

\begin{code}
std::vector<double> scratch_indep[5];
scratch_indep.push_back(1);
scratch_indep.push_back(2);
scratch_indep.push_back(3);
scratch_indep.push_back(5);
scratch_indep.push_back(8);
table_lookup.load_independent_data( my_independent_variable,
                                    scratch_indep);

\end{code}
Notes:
\begin{itemize}
 \item The vector is passed by reference and received as a \textit{const}.
 \item This version does not need the array size; the size is taken
 from the size of the STL-vector.
  That is not possible for the array-method, which passes only a pointer
 which contains no intrinsic information about the size of the array to
 which it points.
 \item Behind the scenes, this method \textit{will} be used.
 The array-method internally converts its data array into a STL-vector and
 calls this STL-vector-method.
\end{itemize}

\item In both cases, there are two optional arguments that may be applied
 at the
 end (arguments 4,5 for array-method; arguments 3,4 for vector-method).
 Note that if the second optional argument is needed, the first must also
 be specified.
\begin{enumerate}
 \item \textbf{TableIndependentVariable::LookupMethod}
 specifies how the independent
 variable's value will be interpreted for generating the dependent variable.
 Options are discussed in the Model Specifications section (chapter
 \ref{sec:Model_spec}), and include:
  \begin{itemize}
   \item TableIndependentVariable::Interp (default)
   \item TableIndependentVariable::Prev
   \item TableIndependentVariable::Next
   \item TableIndependentVariable::Floor
   \item TableIndependentVariable::Ceil
   \item TableIndependentVariable::Round
  \end{itemize}
 \item \textbf{TableIndependentVariable::Continuity}
 specifies how the domain of the function will be identified.
 Options include:
  \begin{itemize}
   \item TableIndependentVariable::Linear (default)
   \item TableIndependentVariable::WrapAround
  \end{itemize}

 For the \textit{Linear} option, the domain is
 as specified, i.e.  $\{1,2,3,5,8\}$ (discrete) or $[1,8]$ (continuous).
  For the \textit{WrapAround} option,
 the domain is extended under the assumption that the first value wraps to
 the last value and the last value wraps around to the first.  This feature is
 discussed more comprehensively in the Model Specifications section (chapter
 \ref{sec:Model_spec}),
\end{enumerate}

Note -- with either method, the values are copied from \textit{scratch\_indep}
 into the model for permanent storage.
The variable \textit{scratch\_indep} may then be overwritten or fall out of
 scope without repercussions.
\end{itemize}

\subsubsection{Dependent Variable Data}
\label{sec:UG_dependent_variable_data}
Create data for the dependent variable(s) and load those into the table.
  There are multiple options for performing this task, these are
 discussed in the next section for the general case with multiple
 dependent variables.  The simplest versions for the case of a single
 dependent variable are presented below.  You may recognize that their
 appearance is very similar to that for the independent data load;
 the advantages/disadvantages of array over vector are identical to those
 presented in the discussion of loading independent data.
\begin{itemize}
 \item \textbf{Array}
\begin{code}
 double scratch_dep[5] = {10, 7, 4, 1, -2};
 table_lookup.load_dependent_data( my_dependent_var,
                                   scratch_dep,
                                   5);
\end{code}

\item \textbf{STL-vector}
\begin{code}
std::vector<double> scratch_dep[5];
scratch_dep.push_back(10);
scratch_dep.push_back(7);
scratch_dep.push_back(4);
scratch_dep.push_back(1);
scratch_dep.push_back(-2);
table_lookup.load_dependent_data( my_dependent_variable,
                                  scratch_dep);
\end{code}

\item In both cases, there is one optional arguments that may be applied
 at the end (argument 4 for array-method; argument 3 for vector-method).
\begin{enumerate}
 \item \textbf{AbstractTableLookup::TableType}
 This argument is used to configure the dependent variable to represent one of
 6 options:
  \begin{itemize}
   \item AbstractTableLookup::Generic (\textbf{default}, linear)
   \item AbstractTableLookup::GenericSingle (linear, limited to 1 independent
   variable; this is marginally more efficient than using \textit{Generic} with
   1 independent variable)
   \item AbstractTableLookup::AngularDeg (applies a linear interpolation along
   the shortest path between two angular values; requires 1 independent
   variable and 1 dependent variable, with a value in degrees)
   \item AbstractTableLookup::AngularRad (applies a linear interpolation along
   the shortest path between two angular values; requires 1 independent
   variable and 1 dependent variable, with a value in radians)
   \item AbstractTableLookup::Quaternion (applies a slerp interpolation to a
   quaternion variable; requires 1 independent variable and a dependent
   variable representing the 4 elements of a quaternion)
   \item AbstractTableLookup::VarDeriv (applies coupled interpolation to a
   variable and its derivative; requires 1 independent variable and 2 dependent
   variables)
  \end{itemize}

 For most applications, this argument may be omitted, in which case it will
 default to produce a simple linear function, using a
 \textit{GenericMultiInputTable}.
 If one of the other options is specified, the model will switch to using
 one of the following table types:
 \begin{itemize}
 \item \textit{GenericSingleInputTable}
 \item \textit{SingleInputTableForAngles}
 \item \textit{SingleInputTableForQuaternions}
 \item \textit{SingleInputTableVarDeriv}
 \end{itemize}
 The differences between these tables are discussed in
 the Model Specifications section (chapter \ref{sec:Model_spec}).
\end{enumerate}

Note -- regardless of which \textit{load\_dependent\_data(...)} method is used,
        the values are copied from \textit{scratch\_dep} into the model for
        permanent storage.  The variable \textit{scratch\_dep} may then
        be overwritten or fall out of scope without repercussions.
\end{itemize}

\subsection{Initialize}\label{sec:UG_single_single_initialize}
The initialize method verifies internal consistency.  It is required to
 perform the data look-up.
\begin{code}
table_lookup.initialize();
\end{code}

\subsection{Configure (optional)}\label{sec:UG_single_single_config}
All instances of \textit{TableIndependentVariable} have the option to
 specify whether the lookup mechanism should be a sequential sweep
 from the last know value (default), or a complete search of the entire
 domain on every cycle.  The latter (non-default) option is faster if
 the value of the independent variable is highly random or irregular.
 It is substantially slower if the incoming value changes slowly (by a
 small fraction of the overall state-space each step) or
 monotonically.

For the case of the \textit{SimpleTableLookup}, the user is not responsible
 for creating a \textit{TableIndependentVariable}, so its presence is hidden.
 But it is still there, it is just created internally.
 A pointer is provided to the
 internally-created \textit{TableIndependentVariable} for
 access to this setting.
\begin{code}
table_lookup.independent->perform_full_search = true;
\end{code}



\subsection{Execute}\label{sec:UG_single_single_exec}
The \textit{update} call is the only executive-level call.
It takes no arguments; the independent and dependent variables
were connected when their respective data was loaded.
\begin{code}
table_lookup.update();
\end{code}


\section{Simple Single-input Multi-output} \label{sec:UG_single_multi}
In this section, we consider the slightly more complicated situation in which
one independent variable is used to perform the lookup for a set of dependent
variables.  It uses the same \textit{SimpleTableLookup} architecture as
was explored in section \ref{sec:UG_single_single}.

Discussion of aspects of this implementation that were covered in section
 \ref{sec:UG_single_single} will not be repeated.  Users should consider
an understanding of section \ref{sec:UG_single_single} as a prerequisite.

\subsection{Instantiate}
Instantiate just as for the case with a single dependent variable, only this
time there will be multiple variables to populate.
\begin{code}
class MyModel
{
 public:
%*\begin{codecomment}
Create a single independent variable, two dependent variables,
and a table to connect them.
\end{codecomment}*)
  double             my_independent_variable;
  double             dep_variable_foo;
  double             dep_variable_bar;
  SimpleTableLookup  table_lookup;
  ...
%*\begin{codecomment}
Constructors:
\end{codecomment}*)
  MyModel():
    table_lookup()
  {};
\end{code}


\subsection{Populate}
\subsubsection{Independent Variable Data}

Create data for the independent variable and load those into the table.
This is equivalent to the configuration for the single output case.
See section \ref{sec:UG_single_single_populate_indep}.

\subsubsection{Dependent Variable Data}
\label{subsec:dep_var_data}
Creating data for multiple dependent variables is similar to the case for
a single output, except now the data array must accommodate both sets of
data.

For example, with 2 dependent variables, and the 5-array already specified
 for the independent variable, I would need a 2-by-5 array (or a 10-element
 array) to represent both dependent variables.

Analogous to how we had two options for loading the data, we also have two
options for specifying the dependent variables themselves.  If the target
dependent
variables are already part of an array and contiguous within that array,
the address of the first dependent variable and the number of variables is
sufficient information.  If the dependent variables are independent of one
another, their addresses should be loaded into a STL-vector and passed into
the \textit{load\_dependent\_data} method.

Consequently, there are six methods for loading dependent data when multiple
dependent variables are being used.
Two have already been covered for the case of a single dependent
variable.  The other four are:
\begin{itemize}
 \item array of variables, array of data
 \item array of variables, vector of data
 \item vector of pointers to variables, array of data
 \item vector of pointers to variables, vector of data
\end{itemize}
Note that these four new methods can be used to load data for the
 single dependent variable case (the array is specified to have 1 element,
 or the vector has 1 element).  However, the single-variable methods
 cannot be used to load multiple variables.

Analagous to the presentation on loading data in section
\ref{sec:UG_single_single}, the use of the
array-method for specifying the dependent variables must be accompanied
by a value specifying the number of
elements to extract from that array.  The use of the vector-method does
not require that additional argument.

However, in a departure from the presentation in section
 \ref{sec:UG_single_single}, the value specifying the number of data points
 per variable must be passed in regardless of whether the \textit{data}
 is presented
 as an array or a STL-vector.

For all four new methods, the optional argument presented in section
 \ref{sec:UG_single_single} (to assign the dependents as angular or quaternion
 variables), is still available.
  If used, the setting is applied to all dependent variables in this table.
  The simple table lookup cannot handle different types of dependent data; if
  that capability is desired, it is necessary to use the more generic
  \textit{TableLookupSet}.

Consequently, we have 4 new function signatures, any one of which may be used:
\begin{code}
  bool load_dependent_data(
           double       * dependent_vars,
           const size_t   num_vars,
           const double * data,
           const size_t   num_data_points_per_variable,
           TableType      type = GenericSingle);

  bool load_dependent_data(
           double       * dependent_vars,
           const size_t   num_vars,
           const std::vector< double> & data,
           const size_t   num_data_points_per_variable,
           TableType      type = GenericSingle);

  bool load_dependent_data(
           const std::vector< double *> & dependent_vars,
           const double * data,
           const size_t   num_data_points_per_variable,
           TableType      type = GenericSingle);

  bool load_dependent_data(
           const std::vector< double *> & dependent_vars,
           const std::vector< double>   & data,
           const size_t   num_data_points_per_variable,
           TableType      type = GenericSingle);
\end{code}

Note -- as discussed previously, the \textit{TableType} argument is optional
and will
default to \textit{Generic}. because I know we are using single-input
interpolation in this example, I specified it to use the marginally more
efficient \textit{GenericSingle} form.

In our example, our two dependent variables are independent instantiations,
so I will use the 3rd of these options.
\begin{code}
std::vector<double *> deps;
deps.push_back( &dep_variable_foo);
deps.push_back( &dep_variable_bar);
double scratch_dep[2][5] = {{8,7,6,5,4},
                            {5,1,3,9,2}};
table_lookup.load_dependent_data( deps,
                                  &scratch_dep[0][0],
                                  5);
\end{code}
\textbf{Notes:}
\begin{enumerate}
 \item The variable name \textbf{does not} act as the pointer to the first
 data element for a multi-dimensional array.  When \textit{scratch\_dep}
 was created
 as a 2-dimensional array, its data type became \textit{double **}, not the
 \textit{double *} that is required by the \textit{load\_data(...)} method.
 It is necessary (and sufficient) to explicitly
 take the address of the first element with the $\& var[0][0]$ syntax.
 \item When declaring a multi-dimensional array, the curly-braces represent
 the data assigned to one particular level of one particular dimension, and
 they are ordered accordingly.
\item Multi-dimensional array data should be coded starting with the last
 dimension and working backwards.
 In the example, with a [2][5] array, the data is entered as 2 sets of 5
 values.
\item The front index ($2$ in this example) is used to separate the variables,
the second index is used to separate the calibration values of each variable.
\end{enumerate}


\subsection{Initialize}
See section \ref{sec:UG_single_single_initialize}
\subsection{Configure (optional)}
See section \ref{sec:UG_single_single_config}
\subsection{Execute}
See section \ref{sec:UG_single_single_exec}


\section{Simple Single-input, Outputs as Angles}
 The use of the \textit{SimpleTableLookup} has already been presented in
 sections \ref{sec:UG_single_single} and \ref{sec:UG_single_multi}.
 Typically where angles are being extracted, a single output is the most likely
 desired response
 (see section \ref{sec:UG_single_single}), but the model will handle
 multiple outputs (see section \ref{sec:UG_single_multi}).

 In the process of loading the data, if the optional argument is specified as
 either
 \begin{itemize}
  \item AbstractTableLookup::AngularDeg (angular, value in degrees)
  \item AbstractTableLookup::AngularRad (angular, value in radians)
 \end{itemize}
 then the outputs will be interpreted as wrapping around at either $180$
 (to $-180$) or $\pi$ (to $-\pi$) depending on the selection made.

 Example:
 \begin{code}
my_table.load_dependent_data( my_angle_output
                              my_angular_data,
                              AbstractTableLookup::AngularDeg);
 \end{code}

 Be aware that:
 \begin{itemize}
  \item the dependent variable(s) in this table-manager instance have an output
        limited to $(-180,180]$ or $(-\pi, \pi]$.  Even if the table-data is
        input in a $[0, 2 \pi)$ format, the output will still be bounded by
        $(-\pi, \pi]$.
  \item it is not possible to mix behaviors of the output variables
        with a \textit{SimpleTableLookup} instance; if the outputs are specified
        as being \textit{AngularDeg}, then \textit{all} outputs will be in
        degrees.  Where it is desired to  provide a mix of interpreted
        variables off the same independent variable (e.g. an angle in
        degrees, another angle in radians, and/or a non-angular output), it
        is necessary to use the \textit{TableLookupSet} even when there is
        only one independent variable.
 \end{itemize}
 \textbf{Example} - suppose the \textit{AngularDeg} option were chosen and the
 data table were input with values
 $[20, 160, 200, 340, 20]$


\fbox{\begin{minipage}[c]{0.8\linewidth}
\textbf{Aside} -
 For any reader questioning the requirement on uniqueness and
 monotonicity of data at this point, I will point out that those
 requirements are placed on the
 independent variables.  These are dependent variable data; neither
 requirement applies to data sets for dependent variables.
 \end{minipage}}


 The table shows how different inputs result in their corresponding outputs.
 \begin{center}
\begin{tabular}{|l|l|l|l|}
\hline
\textbf{input index} & \textbf{input fraction} & \textbf{raw output / deg} &
 \textbf{output /deg}\\
\hline
1 & 0.45 & $160 + 0.45 (200-160) = 178$ & $178$\\
1 & 0.55 & $160 + 0.55 (200-160) = 182$ & $-178$\\
3 & 0.30 & $340 + 0.30 (20-340+360) = 352$ & $-8$\\
3 & 0.60 & $340 + 0.60 (20-340+360) = 364$ & $4$\\
\hline
 \end{tabular}
 \end{center}

\textbf{Note well} -- if the \textit{AngularRad} option were specified
 with these data, the outputs would keep wrapping on multiples of $2\pi$
 and the outputs would be nonsensical $\{2.07, -2.07, -1.71, -2.28\}$

\section{Simple Single-input, Output as Quaternions}
The model provides a Spherical Linear Interpolation (\textit{slerp}) algorithm
to support interpolation between quaternion values.  In the process of loading
the data, if the optional argument is specified as
\textit{AbstractTableLookup::Quaternion}, then the manager will have access to
the slerp algorithm.

Note that this implementation comes with limitations:
\begin{itemize}
 \item The algorithm can only process one input value, so it is well-suited to
 the \textit{SimpleTableLookup}.
 \item The algorithm can only process \textbf{exactly 4} outputs.
 \begin{itemize}
  \item Unlike the angular interpolation option, the quaternion option does
        not support interpolation of multiple quaternions within the
        \textit{SimpleTableLookup} framework.  Providing multiple quaternion
        outputs, or a mix-and-match of quaternion outputs with non-quaternion
        outputs requires the use of the \textit{TableLookupSet}.
  \item While it is possible to use a quaternion-table to populate a JEOD
        Quaternion data type, this feature requires the use of the
        \textit{TableLookupSet} and is not supported in the automatic
        table-creation used by the \textit{SimpleTableLookup}.
 \end{itemize}
\end{itemize}

\section{Single-input, Variable+Derivative Output}
The model provides the capability to couple a variable and its derivative, with
these representing scalar or vector variables of user-specifiable dimension.
Each instance of \textit{SingleInputTableVarDeriv} extends the concept of
\textit{GenericSingleInputTable} (by inheritance, it is a
\textit{GenericSingleInputTable}), to represent the variable, and adds an
additional instance of \textit{GenericSingleInputTable} to represent the
derivative.

The algorithm accessed by this option will always generate an even number of
outputs, being the variable and its derivative for each dimension.

Using the \textit{SimpleTableLookup} format to configure this option is a
little more involved than that for the other simple-table options.
The complexity comes about in identifying how to unambiguously
load coupled data. If the data is put into a single vector or array, it is not
obvious whether it should be intepreted as containing data for each dimension,
separated into variable and derivative, or as being presented with data for
all dimensions of variable then all dimensions of derivative.
This ambiguity seems like a perfect opportunity for mis-configuration, so we
do neither; loading data with a single array is not supported.

Instead, the \textit{load\_dependent\_data} method loads ONLY the data for the
variable, sizing the implementation accordingly.
\begin{code}
simple_table.load_dependent_data( x,
                                  5,
                                  &x_data[0][0],
                                  6,
                                  AbstractTableLookup::VarDeriv);
\end{code}
This constructs the instance of \textit{SingleInputTableVarDeriv} and populates
the ``intrinisic instance'' of \textit{GenericSingleInputTable} with data for the
variable. The instance of \textit{GenericSingleInputTable} for the derivative
is left empty and must be loaded independently.

We start by getting a pointer to the newly created table:
\begin{code}
SingleInputTableVarDeriv * table =
  dynamic_cast<SingleInputTableVarDeriv*>(simple_table.get_table());
\end{code}
Then we add the dependent variables to the table's \textit{deriv} instance of
\textit{GenericSingleInputTable}
\begin{code}
for (int ii = 0; ii<5; ++ii) {
  table->derivs.add_dependent(xdot[ii]);
}
\end{code}
and finally, load data for this table:
\begin{code}
table->derivs.load_data(&xdot_data[0][0], {5,6});
\end{code}
\begin{code}
\end{code}



\section{General Use with Single Input}
The rest of the User's Guide will focus on using the \textit{TableLookupSet}
management process (rather than the \textit{SimpleTableLookup} which has
been covered to this point). Typically, the justification for needing to
implement the \textit{TableLookupSet} is that interpolation is required
across multiple inputs.  However, it is also required where there is only one
input and multiple outputs \textit{of different types}, such as an angle or
quaternion.

The typical implementation involving a \textit{TableLookupSet} is presented
in section \ref{sec:UG_multi_single}.  It is recommended to review that
content before proceeding, bearing in mind that it will not be necessary to
have more than one \textit{TableIndependentVariable} when there is only 1
input.

Note - it is possible to use the \textit{TableLookupSet} architecture to
support the simple case of a single input and a single output type.  In this
case, the content of section \ref{sec:UG_multi_single} should be
comprehensive.

However, where the situation demands outputs of different types, it must be
recognized that each data-table instance is limited to supporting a single
type.
When there are output variables of different types, it is necessary to
implement at least as many tables as there are types.  The instructions
for adding
multiple tables to a \textit{TableLookupSet} manager are presented in section
\ref{sec:UG_multi_multi_multi}.

Finally, the instructions and options for implementing specific instances of
\textit{SingleInputTableForAngles},
\textit{SingleInputTableForQuaternions} and \textit{SingleInputTableVarDeriv}
for addition to a table-set are found in sections
\ref{sec:UG_angles}, \ref{sec:UG_quaternions}, and \ref{sec:UG_var_deriv}
respectively.

\section{Multi-input Single-output} \label{sec:UG_multi_single}
In this introduction to the case involving multiple input variables,
we will consider the simplest multi-input scenario, that in which a single
output value is derived from multiple inputs.

Of particular note, this general case may be applied to a single input
scenario, but because the model complexity is significantly higher, the
 configuration is substantially more complex than for the implementation
 described in section \ref{sec:UG_single_single}.

One of the important concepts that is necessary to advance beyond the simple
 single-input scenarios is that of the \textit{TableLookupSet}.  The
\textit{TableLookupSet} represents a collection of lookup-tables and
 independent variables that are intended to be used concurrently.

The strength of the \textit{TableLookupSet} architecture is found in cases
where each of the dependent variables depends on different subsets of the
available  independent variables.  In those cases, there are three options:
\begin{itemize}
 \item Use the same independents for all dependents, and fill the table array
 with duplicated data to cover the redundant independents
 \item Use a different table for each combination of independent variables
 and treat them as independent entities.  Make the necessary calls
 for each table, processing the common independent variables multiple times.
 \item Use a different table for each combination of independent variables
 and collect them into one place.  Make a single call to the collection, which
 can process the independents as needed and share the common variables across
 the collection.  This is the chosen architecture.
\end{itemize}

The concept of a set of lookup-tables is not directly relevant for a
scenario resulting in a single output.  However,
there is substantial infrastructure associated with this architecture that
we can utilize for this single-output case.
Thus, even though it is completely unnecessary
to have the capability for a set of tables in this case, it is actually easier
to take advantage of the wrapping that manages the more complex capability
than it is to make those calls individually.


\subsection{Instantiate}
In the example presented below, there are instantiations for the independent
 and dependent variables,
although these would typically already exist.

Declare instances of
\begin{itemize}
 \item \textit{TableLookupSet} (1 instance can manage multiple tables)
 \item Independent Variables (1 instance per independent variable).
        Each independent variable is provided with
        a table of its own; the location of the value of the independent
        variable
        in its table determines the location of where to find the value of the
        output variable in its table.  There are two options, described in more
        detail in section \ref{sec:tab_ind_var}.
 \begin{itemize}
  \item \textit{TableIndependentVariable} with \textit{Linear Continuity}
  covers a finite domain
  \item \textit{TableIndependentVariable} with \textit{WrapAround Continuity}
  wraps its domain to produce a cyclic infinite domain
 \end{itemize}

 \item Lookup Tables.  Each lookup table represents a `block' of data.
\begin{itemize}
 \item \textit{GenericMultiInputTable} provides any number of outputs from
 any number of inputs.
\end{itemize}
\end{itemize}

\begin{code}
class MyModel
{
public:
%*\begin{codecomment}
Create four independent variables and a dependent variable
\end{codecomment}*)
  double                   my_independent_variable0;
  double                   my_independent_variable1;
  double                   my_independent_variable2;
  double                   my_independent_variable3;
  double                   my_dependent_variable;
%*\begin{codecomment}
Instantiate a TableLookupSet as a `manager' for the table and its variables. \\
 Create four TableIndependentVariable instances to represent the four \\
 independent variables. Create one table.
\end{codecomment}*)
TableLookupSet           table_set;

TableIndependentVariable table_independents0;
TableIndependentVariable table_independents1;
TableIndependentVariable table_independents2;
TableIndependentVariable table_independents3;

GenericMultiInputTable   table;
\end{code}


\subsection{Construct}
Provide the following arguments to the constructors:
\begin{itemize}
 \item \textit{TableLookupSet} takes no argument
 \item \textit{TableIndependentVariable} can take 1, 2 or 3 arguments,
 those being:
\begin{enumerate}
 \item a character-string (std::string) name for the independent variable
  (\textbf{optional}).  The name can be useful for debugging and for
  connecting independents with the tables.  It may also be assigned
  later with the \textit{set\_name} method.  When specified,
  names must be unique.
 \item a reference to the independent variable to be associated with
       this instance (\textbf{required})
 \item An epsilon criterion for testing whether the variable is
  sufficiently close to the calibrated value to call it exact.
  Note that this is typically only useful when the data is being
  discretized; it stops very small deviations from incorrectly triggering
  into the next cell.  This argument is optional, and defaults to
  $1.0 \times 10^-9$ if not specified.  (\textbf{optional}).
\end{enumerate}

 \item \textit{GenericMultiInputTable} can take 0, 1, or 2 arguments that
  are used to populate its associations with particular dependent variables.
  The options are:
 \begin{itemize}
 \item 0 arguments.  The table must be connected with its dependent
  variables later by using the \textit{populate\_output} method.
  This method takes a
  STL-vector of pointers to the intended dependent variables.
 \item 1 argument.  A single reference to the only dependent variable
  connected to this table.
 \item 1 argument. A STL-vector of pointers to all of the dependent variables
  connected to this table.
 \item 2 arguments:
  \begin{enumerate}
   \item A pointer to the first (relevant) component of an array of values.
   \item The number of (contiguous) variables to extract from said array.
  \end{enumerate}
 \end{itemize}
\end{itemize}

\begin{code}
  MyModel():
    table_set(),
%*\begin{codecomment}
The three different constructors for the TableIndependentVariable look like:
\end{codecomment}*)
    table_independents0(my_independent_variable0),
    table_independents1(my_independent_variable1),
    table_independents2(``my_var_2'',
                        my_independent_variable[2]),
    table_independents3(``my_var_3''
                        my_independent_variable[3],
                        1.0E-12),
    table(my_dependent_variable),
    ...
\end{code}

\subsection{Populate}
\subsubsection{Independent Variable Data}
Each independent must have its own data-set representing the range of values
 that the
variable can take.  The \textit{load\_data} method is used for this purpose;
 it has a structure similar to that discussed when loading the independent
 data for the \textit{SimpleTableLookup} (indeed, the
 \textit{SimpleTableLookup} just called this method behind the
 scenes).

As we saw in the \textit{SimpleTableLookup} discussion, the arguments to
 this method can be either:
\begin{itemize}
 \item a STL-vector
 \item a pointer and the number of elements to extract starting at the pointer.
\end{itemize}

\begin{code}
  double scratch_ind0[2] = {11,12};
  table_independents0.load_data(scratch_ind0, 2);
%*\begin{codecomment}
Assign these for each independent.
\end{codecomment}*)
  ...
  double scratch_ind3[5] = {1,2,3,5,8};
  table_independents3.load_data(scratch_ind3, 5);
\end{code}

\subsubsection{Dependent Variable Data}
Each table must have its own data set representing the values that the
 dependent variables can take.
The dimension of this data-set is always 1 higher than the number of
 independent variables.  The first dimension is sized equal to the number
 of dependent variables; each of the other dimensions are sized
 according to the size of the data array for each of the independent
 variables.

In the case presented below, we have four independent variables, with
 2, 3, 4 and 5 elements respectively.  They are being used to compute
 the value of
 a single output (dependent variable).
 Note that for the case of a single output, the first
 dimension, 1, is redundant, and is not used when creating the data array
 in this example.

When populating a multi-dimensional array, the data should be entered by
 the last dimension first and working backwards.  We have already seen an
 example in the previous section where a [2][5] array was entered as 2
 strings of 5 values.
 For a five-dimensional [1][2][3][4][5] array, the data would be entered as
 1 group of 2 collections of 3 sets of 4 strings of 5 values:

The \textit{load\_data} takes 2 arguments, being either:
\begin{itemize}
 \item \begin{enumerate}
        \item A pointer to the first element of the data array
        \item An STL-vector with elements equal to the number of dependent
         variables followed by the number of elements for
         each of the independent variables.
       \end{enumerate}
 \item \begin{enumerate}
        \item A STL-vector of data (passed by reference)
        \item An STL-vector with elements equal to the number of dependent
         variables followed by the number of elements for
         each of the independent variables.
       \end{enumerate}
\end{itemize}

\begin{code}
double temp_table[2][3][4][5] =
                       { { {
/*temp_table[0][0][0][0:4]*/ {   1,   2,   3,   4,   5},
/*temp_table[0][0][1][0:4]*/ {  11,  12,  13,  14,  15},
/*temp_table[0][0][2][0:4]*/ {  21,  22,  23,  24,  25},
/*temp_table[0][0][3][0:4]*/ {  31,  32,  33,  34,  35} },
/*temp_table[0][1][0:3][0:4]*/
                            {{ 101, 102, 103, 104, 105},
                             { 111, 112, 113, 114, 115},
                             { 121, 122, 123, 124, 125},
                             { 131, 132, 133, 134, 135} },
/*temp_table[0][2][0:3][0:4]*/
                            {{ 201, 202, 203, 204, 205},
                             { 211, 212, 213, 214, 215},
                             { 221, 222, 223, 224, 225},
                             { 231, 232, 233, 234, 235} } },

/*temp_table[1][0:2][0:3][0:4]*/
                           {{{1001,1002,1003,1004,1005},
                             {1011,1012,1013,1014,1015},
                             {1021,1022,1023,1024,1025},
                             {1031,1032,1033,1034,1035} },

                            {{1101,1102,1103,1104,1105},
                             {1111,1112,1113,1114,1115},
                             {1121,1122,1123,1124,1125},
                             {1131,1132,1133,1134,1135} },

                            {{1101,1102,1103,1104,1105},
                             {1111,1112,1113,1114,1115},
                             {1121,1122,1123,1124,1125},
                             {1131,1132,1133,1134,1135} } } };
%*\begin{codecomment}
Create the sizing vector.\\
Remember to specify the first dimension, which is always the number of \
dependent variables.\\
The other sizes should follow in the same order in which the array was \
constructed.\\This vector should always have at least two elements.
\end{codecomment}*)
std::vector<size_t> size_vec;
size_vec.push_back(1);
size_vec.push_back(2);
size_vec.push_back(3);
size_vec.push_back(4);
size_vec.push_back(5);
%*\begin{codecomment}
Now load these data into the table.
\end{codecomment}*)
table.load_data(&temp_table[0][0][0][0],size_vec);
\end{code}
\textbf{Important Note:  It is vitally important to get these numbers correct.}
Particularly when loading via the array-method, there is no verification
 possible to ensure that the table just created matches the size specified
 by the integers that are passed in via the sizing vector.
It is really only possible to check the first element of the sizing vector
 against the number of outputs registered.  Beyond that, the code will use the
 values in the sizing vector to compute the total number of elements to copy,
 and will grab that much data off the array.  If the array is too large, the
 model will grab only part of the array; if the array is too small, the model
 will grab data from beyond the array.

Loading using the vector-method is safer but less convenient.  The number of
elements in the vector will be compared against the expected number resulting
from the values in the sizing vector.

\subsection{Connect}
\subsubsection{Add Elements to the Manager}
Add the tables and the independents to the manager:
\begin{code}
    table_set.add_independent_variable(table_independents0);
    table_set.add_independent_variable(table_independents1);
    table_set.add_independent_variable(table_independents2);
    table_set.add_independent_variable(table_independents3);
    table_set.add_table( single_table);
\end{code}

\subsubsection{Connect Tables with Independents}

The \textit{associate\_table\_and\_independent} method is used to connect each
 table with its independent variables.
Each table maintains a STL-vector of its associated independent variables;
this method populates that vector.

\textbf{Important Aside:  It is important to understand the order in which the
 independent variables are indexed in the data array.}

For example, suppose some table, \textit{table1}, uses season, altitude, and
 time-of-day to generate atmospheric pressure.
This table would provide a 3-dimensional array, with the output value
 identified from the three respective indices:
\begin{equation*}
 pressure = table1\,[season-index]\,[altitude-index]\,[time-index]
\end{equation*}

In this case, the independents \textbf{must} be ordered in the STL-vector in
 that order -- season, then altitude, then time-of-day.
There is no automated checking possible to ensure that this has been done
 correctly.  Consider the following example to illustrate the potential problem
 if this is done incorrectly.

Suppose the pressure table is based on 2 season-values, 3 altitude-values, and
 4 time-values.  The table is generated as a 2x3x4 array:
\begin{code}
double scratch[2][3][4] = {{{ 1.0,  2.0,  3.0,  4.0},
                            {11.0, 12.0, 13.0, 14.0},
                            {          ...           }}};
\end{code}
If we need to look-up the value associated with the 1st season-value
 (index 0 represents the first element),
 2nd altitude-value (index = 1), and 3rd time-value (index = 2), the resulting
 table index would evaluate to

$index = 0 \times (3 \times 4) + 1 \times (4) + 2 = 6$.

 which directs the code to the 7th data element, for an output value of 13.0.

If, however, those variables are loaded in the wrong order of altitude, time,
 season, then the 1st season, 2nd altitude, and 3rd time indices would evaluate
 to

$index = 1 \times (3 \times 4) + 2 \times (4) + 0 = 20$.

 which will result in the wrong value being extracted from the table.

\textbf{End of Aside}

Returning now to the discussion of how to configure the table-lookup:
There are four pieces of information that are needed to correctly connect
the table and its independent variables.  Some or all of these four can be
assumed.
\begin{enumerate}
 \item Identification of which table to use.  This is specified as a
  \textit{GenericMultiInputTable} pointer (note - the alternative tables
  \textit{SingleInputTableForAngles} and
  \textit{SingleInputTableForQuaternions} both derive from
  \textit{GenericMultiInputTable}, so pointers to these table-types are
  acceptable).  If omitted, it defaults to the
  last table added to the manager.

 \item The independent variable.  This can be specified as either:
 \begin{itemize}
  \item a \textit{TableIndependentVariable} pointer, or
  \item the \textit{std::string} name of the
        \textit{TableIndependentVariable} instance.
 \end{itemize}
 If omitted, it defaults to the last \textit{TableIndependentVariable}
 added to the manager.

 \item The index in the selected \textit{GenericMultiInputTable}'s
  STL-vector of independents into which the pointer to the selected
  \textit{TableIndependentVariable} will be inserted. This is specified as a
  \textit{size\_t} index (starting with 0).   If this argument is such that
  it will overwrite an existing element, a warning message will be provided
  \textbf{and the overwrite will be executed.}  If omitted, the selected
  \textit{TableIndependentVariable} will be added to the end of the selected
  \textit{GenericMultiInputTable}'s list of independents.

 \item The interpretation to be used by the selected
 \textit{GenericMultiInputTable} when extracting index data from the selected
 \textit{TableIndependentVariable}.  The options are:
  \begin{itemize}
   \item TableIndependentVariable::Interp (default)
   \item TableIndependentVariable::Prev
   \item TableIndependentVariable::Next
   \item TableIndependentVariable::Floor
   \item TableIndependentVariable::Ceil
   \item TableIndependentVariable::Round
  \end{itemize}
  If omitted, the default (Interp) option will be selected.
\end{enumerate}

\textbf{Notes:}
\begin{itemize}
 \item That makes a total of 24 different signatures.
 \item The name option can only be used if the
  \textit{TableIndependentVariable}  instance
  has been given a name, in the constructor for example.
 \item Pushing a \textit{TableIndependentVariables} onto a large index will
  result in the creation of NULL values to fill the gap.  Those must be
  populated before the model is initialized.  The model cannot ``skip''
  indices.
 \item The following 3 examples are equivalent ways of adding 2 independent
 variables:
\begin{code}
table_set.add_independent_variable(&table_independents0);
table_set.add_independent_variable(&table_independents1);
\end{code}

\begin{code}
table_set.add_independent_variable(&table_independents0,0);
table_set.add_independent_variable(&table_independents1,1);
\end{code}

\begin{code}
table_set.add_independent_variable(&table_independents1,1);
table_set.add_independent_variable(&table_independents0,0);
\end{code}
\end{itemize}


The simplest procedure for this relatively simple example is to make the
 associations while adding the elements, add the independents in the order
 in which the table(s) use them, and to use the no-argument option.

\begin{code}
  table_set.add_table( single_table);
  table_set.add_independent_variable(table_independents0);
  table_set.associate_table_and_independent();
  table_set.add_independent_variable(table_independents1);
  table_set.associate_table_and_independent();
  table_set.add_independent_variable(table_independents2);
  table_set.associate_table_and_independent();
  table_set.add_independent_variable(table_independents3);
  table_set.associate_table_and_independent();
\end{code}

\subsection{Alternative Implementations}
\subsubsection{SingleInputTableForAngles}
A \textit{SingleInputTableForAngles} instance may be substituted
in place of a \textit{GenericMultiInputTable} instance in the
\textit{TableLookupSet} collection of tables.

The construction of this specialized type of a
 \textit{GenericMultiInputTable} follows
the pattern previously presented for the \textit{GenericMultiInputTable}
base class, with one exception:

The \textit{SingleInputTableForAngles} constructor takes an optional
additional argument to specify the units on the dependent variable.
This \textit{bool} argument defaults to true, indicating that the units
will be \textbf{radians}.  If \textbf{degrees} are required, simply
add a false to the end
of the construction arguments for the \textit{GenericMultiInputTable}.

Example:
\begin{code}
class MyClass:
  ...
  SingleInputTableForAngles table;
  ...
  MyClass()
    :
    ...
    table( my_dependent_var,
           false),
    ...
\end{code}

\textbf{Reminder }-- all dependent variables covered by this table will be
interpreted as angles, and they will all have the same units.

\subsection{Initialize}

The initialize method verifies internal consistency.  All allocation/creation,
 population, registration, and connection processes
 should be completed before this is run.
Initialization is required to perform the data
 look-up.  Initialization of the table-set performs
 initialization of all elements that have been added to the set.
\begin{code}
table_set.initialize();
\end{code}


\subsection{Configure (optional)}
Section \ref{sec:UG_single_single_config} explores the use of the
\textit{perform\_full\_search} option.
This flag can be set directly for any of the
\textit{TableIndependentVariable} instances.

\subsection{Execute}
The update call is the only executive-level call.  It takes no arguments; the
 input and output variables are pre-connected. It is important to
understand that \textit{update()} first calculates the values for all
independent variables for all tables, and then calculates all dependent
values. It is possible to use the same outside variable as the value
location for an independent variable and a dependent variable in the same
table or different tables. This could be used to, e.g., model a feedback
loop.  Because \textit{update()} has such a definite processing order, the
end results are consistent and repeatable.

\begin{code}
table_set.update();
\end{code}

\section{Multi-input Multi-output Single-Table} \label{sec:UG_multi_multi}
This more general case of the multiple-input scenario is a simple extension
of the discussion in section \ref{sec:UG_multi_single}.

There are only two substantive differences between a single-output
 and multi-output implementation
\begin{itemize}
 \item When constructing the \textit{GenericMultiInputTable}, the optional
 argument using a reference to a dependent variable is no longer an option
 because there is more than one such variable.  The other argument
 options remain valid and unchanged.
 \item When creating the data array (if using the array-method), it is
 important to specify the first dimension size to be equal to the
 number of dependent variables.
 This was discussed in the previous section, but the first dimension was
 omitted because it was just 1.
\end{itemize}

\section{Multi-input Multi-output Multi-table} \label{sec:UG_multi_multi_multi}

This more comprehensive option allows users to implement completely
 general n-dimensional table-lookups.
As described earlier in this document, The \textit{TableLookupSet}
 functionality was designed to support the completely general case where
the output data is split across multiple data tables.
So far, we have seen how to configure a table-lookup scenario in which
a single data table is used to provide the output data.  We will now
 take the very small step to extend that to the completely general case.

\subsection{Instantiate}
A single \textit{TableLookupSet} is required, but it can support any number
 of tables and independent variables.

Create as many tables for dependent variables:
\begin{itemize}
 \item \textit{GenericMultiInputTable}, and/or
 \item \textit{SingleInputTableForAngles}, and/or
 \item \textit{SinglenputTableForQuaternions}, and/or
 \item other table types provided by extension of this model
\end{itemize}
and as many independent variables:
\begin{itemize}
 \item \textit{TableIndependentVariable} or
 \item other independent variable types provided by extension of this model
\end{itemize}
 as are needed.  Remember that each table (e.g.
 \textit{GenericMultiInputTable},
 \textit{SingleInputTableForAngles} may represent any number of dependent
 variables as long as they share common independent variables and are the same
 type -- linear, angular, quaternion, etc.).

 Follow the instructions in sections
 \ref{sec:UG_single_single} - \ref{sec:UG_multi_multi}
for a guide on how to create these.


\subsection{Populate}
Populate the independents and the tables as outlined in the same previous
 sections (\ref{sec:UG_single_single} - \ref{sec:UG_multi_multi}).

\subsection{Connect}
Following the outline presented in section \ref{sec:UG_multi_single}, all
 tables and their
 independent variables must be connected using the
\textit{associate\_table\_and\_independent} method.

The added complexity in this section comes about because when
table-to-variable dependencies cross over, the no-argument version of the
\textit{associate\_table\_and\_independent} method may be insufficient.

Reminder: In section \ref{sec:UG_multi_single}, first the table was added
to the manager, then the independent variables were added.
With the addition of
each independent variable, the connection was made between variable and
table using the easy route of associating the last added variable with the
last-added table.

That may no longer be possible with multiple tables, and care must be
taken to specify connections.


Examples:

Note --- these examples start with adding the tables and
 independents to the table-set.
This step has already been covered in section \ref{sec:UG_multi_single},
and is duplicated here for completeness.
\begin{enumerate}
 \item Simple configuration with a \textit{table\_set} comprising:
\begin{itemize}
 \item 3 independent variables,
  \textit{var0}, \textit{var1}, and \textit{var2}) and
 \item 2 tables
 (\textit{table0}, and \textit{table1})
 \item table0 represents function $f(var0, var1)$
 \item table1 represents function $g(var1, var2)$.
\end{itemize}

\begin{code}
%*\begin{codecomment1}
Add table0, var0 and make var0 the first independent associated with table0.
No arguments needed.
\end{codecomment1}*)
table_set.add_table(table0);
table_set.add_independent_variable(var0);
table_set.associate_table_and_independent();
%*\begin{codecomment1}
Add var1, and make it the second independent associated with table0.
Table0 is still current so no arguments needed.
\end{codecomment1}*)
table_set.add_independent_variable(var1);
table_set.associate_table_and_independent();
%*\begin{codecomment1}
Add table1 and make var1 the first independent associated with it.
Var1 is still current so no arguments needed.
\end{codecomment1}*)
table_set.add_table(table1);
table_set.associate_table_and_independent();
%*\begin{codecomment1}
Add var2 and make it the second independent associated with table0.
Table1 is still current so no arguments needed.
\end{codecomment1}*)
table_set.add_independent_variable(var2);
table_set.associate_table_and_independent();
\end{code}
The notable limitations on this simple method are:
\begin{itemize}
 \item the variables must appear in a convenient order.  If
 \textit{table1} instead represented function $h(var0, var2)$,
 this method will not work.
\item   It is also difficult to maintain because it relies on a certain
 ordering and is not as well structured as it might be.
\end{itemize}

\item Similar configuration using a more generic configuration definition.

In this example, all of the independents get added together, then with
the addition of each table, the associated connections get made to
explicitly identified variables.
Two methods are illustrated for specifying the variables:
\begin{itemize}
 \item table0 variables are specified by name; this requires either
 \begin{itemize}
  \item they were either constructed with a name, or
  \item the \textit{set\_name} method was employed to give them a name following
  their construction.
 \end{itemize}
 \item table1 variables are specified by pointer.
\end{itemize}

\begin{code}
%*\begin{codecomment1}
Start by adding all of the independents
\end{codecomment1}*)
table_set.add_independent_variable(var0);
table_set.add_independent_variable(var1);
table_set.add_independent_variable(var2);

%*\begin{codecomment1}
Add table0 and connect the first two elements to var0 and var1 respectively.\\
Because table0 is current (following its addition), we do not need to specify\
it ... but do for completeness.\\
Because these associations are made in order, the ``0'' and ``1'' in the\
argument list are redundant and could also be removed.
\end{codecomment1}*)
table_set.add_table(table0);
table_set.associate_table_and_independent('variable0',0);
table_set.associate_table_and_independent('variable1',1);

%*\begin{codecomment1}
Add table1 and connect its first two elements to var0 and var2 respectively.
This time, the argument list uses pointers instead of names and no redundant
indexing.
\end{codecomment1}*)
table_set.add_table(table1);
table_set.associate_table_and_independent(&var0);
table_set.associate_table_and_independent(&var2);
\end{code}

\item Similar configuration again, this time with complete separation of
 registration and association:
\begin{code}
%*\begin{codecomment1}
Start by adding all of the independents and tables
\end{codecomment1}*)
table_set.add_independent_variable(var0);
table_set.add_independent_variable(var1);
table_set.add_independent_variable(var2);

table_set.add_table(table0);
table_set.add_table(table1);

%*\begin{codecomment1}
Make the connections as before, explicitly stating both the table and variable
for each connection.  To illustrate that order is not important (if indexing
is provided), the order in which table1's variables are associated has been
reversed.  This reversal is purely for illustration; because reverse-ordering
can confuse maintenance, it is preferable to keep associations ordered.
\end{codecomment1}*)
table_set.associate_table_and_independent( &table0,
                                           'variable0');
table_set.associate_table_and_independent( &table0, &var1);
table_set.associate_table_and_independent( &table1,
                                            &var2,
                                            1);
table_set.associate_table_and_independent( &table1,
                                           'variable0',
                                            0);
\end{code}
\end{enumerate}


\subsection{Initialize}
See section \ref{sec:UG_multi_single}.

\subsection{Configure (optional)}
See sections \ref{sec:UG_multi_single} and \ref{sec:UG_multi_multi}.

\subsection {Execute}
See section \ref{sec:UG_multi_single}.

\section{Single-input Table for Angles}\label{sec:UG_angles}
Thus far, we have focused on using the \textit{GenericMultiInputTable} as the
tabular side of the \textit{Table-LookupSet}.  The
\textit{SingleInputTableForAngles} class is a type of
\textit{GenericMultiInputTable} specifically designed to handle linear
interpolation of angles.  It provides an extra verification step that the
output value lies on the shortest arc between the two calibrated values. So,
for example, the midpoint of the range (355, 5) degrees is 0 degrees, while a
simple arithmetic interpolation would provide 180 degrees.

\subsection{Instantiate}
To switch from a simple \textit{GenericMultiInputTable} to a
\textit{SingleInputTableForAngles} instance, simply replace the class type,
and include the new header.

\subsection{Construct}
The constructor options for the \textit{SingleInputTableForAngles} are
identical to those for the \textit{GenericMultiInputTable} with the exception
that all options have one additional optional argument that may be provided
at the end of the argument list (or as the only argument if adapting the
no-argument \textit{GenericMultiInputTable} constructor).

The final argument is a boolean; it indicates whether the angles should be
interpreted as radians, and it defaults to \textit{true}.

If the angles are interpreted as radians (default), the output will always
lie in the range $(-\pi, \pi]$. For this option,
\textit{no additional steps are necessary}.

If, however, the angles are interpreted as degrees,
the output will always lie in the range $(-180, 180]$.
To make the model switch from radians (default) to degrees, one of two steps
should be taken:
\begin{itemize}
 \item Add a \textit{false} argument as the last argument of the constructor,
   \begin{code}
SingleInputTableForAngles table;
...
<class-constructor>
  :
  table(dep_var,
        false); // angles are degrees (not radians).
   \end{code}
       or
 \item Set \textit{table.output\_in\_radians = false} before the table is
       initialized (e.g. in the input file).
   \begin{code}
table.output_in_radians = False
   \end{code}
\end{itemize}

\subsection{Populate and Connect}

As the class name suggests, this class supports only a single independent
variable.  Because angles are non-additive in multiple dimensions,
interpolating an angle across multiple dimensions is not a well-defined
process and is not supported.

As with the \textit{GenericMultiInputTable}, the dependent variables are
typically specified at construction time, but may be added later (before
initialization) with the \textit{add\_dependent(double \&)} method.

Populating the dependent variables with data (i.e. loading the dependent
data) and making the connections to an independent
variable follow an identical process to those already
presented for instances of \textit{GenericMultiInputTable}.


\section{Single-input Table for Quaternions}\label{sec:UG_quaternions}
This class is a type of \textit{GenericMultiInputTable} that
supports interpolation of a complete quaternion as a single
entity.

While the other table-classes (\textit{GenericMultiInputTable},
\textit{SingleInputTableForAngles}) can support any number of
dependent-variables, each instance of this table will support
exactly 1 quaternion.

\subsection{Instantiate}
To switch from a simple \textit{GenericMultiInputTable} to a
\textit{SingleInputTableForQuaternions} instance, simply replace the class
type, and include the new header.

\subsection{Construct}
There are four constructors available; closely correlated with the 4
available for \textit{GenericMultiInputTable}.

For the first three options, recognize that at initialization there must be
\textit{exactly 4} dependent variables, and they must be
added in the same order as the
quaternion values appear in the data table.  This is typically:
\begin{enumerate}
 \item scalar component
 \item vector i-component
 \item vector j-component
 \item vector k-component
\end{enumerate}


\begin{itemize}
 \item \textbf{No argument}.  Effectively equivalent to the no-argument
       \textit{GenericMultiInputTable} constructor.  With this constructor,
       the 4 dependent variables must be added prior to initialization using
       the \textit{add\_dependent(double \&)} method.

 \item \textbf{(Array, size) argument (double*, size\_t)}.  Analogous to the
       \textit{GenericMultiInputTable} constructor with identical signature.
       Typically, if using this construction method, the 4 elements of the
       quaternion would be loaded at this time and the second argument would
       take the value 4.

       Recognize that because dependent variables cannot be removed, if
       more than 4 elements are added at construction time (2nd argument $>
       4$),
       the table will not initialize successfully and cannot be used.
       If fewer than 4 are added at
       construction time (2nd argument $< 4$), the others must be added prior
       to initialization with the \textit{add\_dependent(double \&)} method.

 \item \textbf{Vector argument (std::vector$<$double*$>$)}.  Analogous to the
       \textit{GenericMultiInputTable} constructor with identical signature.
       Typically, if using this construction method, the vector would contain
       the 4 elements of the quaternion and nothing else.

       Recognize that because dependent variables cannot be removed, if
       the vector contains more than 4 elements,
       the table will not initialize successfully and cannot be used.
       If it contains fewer than 4 elements, the others must be added prior
       to initialization with the \textit{add\_dependent( double \&)} method.

 \item \textbf{Quaternion argument (jeod::Quaternion \&)}  The Quaternion class is a
       JEOD construction used widely in dynamics.  This option replaces the
       single-variable \textit{GenericMultiInputTable} constructor signature.

       If constructed with a reference to a Quaternion instance, this
       class will populate the elements of that instance directly.
       If using this construction method, it is assumed that the data is
       ordered according to the sequence presented above.
\end{itemize}

\subsection{Populate and Connect}

As the class name suggests, this class supports only a single independent
variable.  Because quaternions are non-additive in multiple dimensions,
interpolating across multiple dimensions is not a well-defined
process and is not supported.

Populating the dependent variables with data (i.e. loading the dependent
data) and making the connections to an independent
variable follow an identical process to those already
presented for instances of \textit{GenericMultiInputTable}.  Remember to
confirm that the data are ordered correctly; mis-ordering them
is an easy mistake to make when they are just components of a single variable.


\section{Single-input Table for Variable and Derivative}
\label{sec:UG_var_deriv}
The \textit{SingleInputTableVarDeriv} extension of
\textit{GenericMultiInputTable} supports the coupled interpolation of a
variable (scalar or vector) and its derivative (likewise).

An important note is that the one independent variable
\textbf{MUST BE the variable with respect to which the derivative is taken}.
For example, position and velocity may be coupled when the independent variable
is time, but not when it is some other variable, such as the position of
another vehicle.

\subsection{Instantiate}
To switch from a simple \textit{GenericMultiInputTable} to a
\textit{SingleInputTableVarDeriv} instance, simply replace the class
type, and include the new header.
\begin{code}
#include "../../include/single_input_table_var_with_deriv.hh"
SingleInputTableVarDeriv  interpolation_table_x_xdot;
\end{code}
Each instance includes a pair of \textit{GenericSingleInputTable} tables, one
intrinsically included within the class by inheritance/extension, and one as a
class member; the latter is named \textit{deriv}.

\subsection{Construct}
There are multiple construction options for this type:
\begin{itemize}
\item All of the constructors for \textit{GenericMultiInputTable} are
available, the options (described previously) will be applied to the intrinsic
instance of \textit{GenericSingleInputTable}, leaving the \textit{derivs}
instance requiring independent configuration.
\item There are 3 class-specific constructors that will connect both the
intrinisic table and the derivs table to the dependent variables:
  \begin{itemize}
  \item 3-argument form, addressing 2 C-style arrays of variables and the number
  of members of these arrays:
  \begin{code}
  SingleInputTableVarDeriv(
            double *dependent_variables,
            double *dependent_variables_derivs,
            size_t num_vars);
  \end{code}

  \item 2-argument form, using references to the variable and its derivative:
  \begin{code}
  SingleInputTableVarDeriv( double & dependent_var,
                            double & dependent_var_deriv);
  \end{code}

  \item 2-argument form using STL-vectors of pointers to the dependent variables:
  \begin{code}
  SingleInputTableVarDeriv(
            const DoublePtrVec & dependent_variables,
            const DoublePtrVec & dependent_variables_deriv);
  \end{code}
  \end{itemize}
\end{itemize}

\subsection{Connecting to Variables}
As the class name suggests, this class supports only a single independent
variable; this is assigned to the table-set to which the tables will be added:
\begin{code}
table_set.add_independent_variable( tiv);
\end{code}

The class may support any (even) number of dependent variables:
any number of variables and a derivative for each variable.
If using the standard \textit{GenericMultiInputTable} constructors, the
dependent variables for the derivatives will have to be provided separately:
\begin{code}
for (int ii = 0; ii < 5; ++ii) {
   coupled_table.derivs.add_dependent(xdot[ii]);
\end{code}
If using one of the class-specific constructors, these already populate the
dependent variables for both the zeroth- and first- derivative variables.


\subsection{Populate Table Data}
Populating the dependent variables with data (i.e. loading the dependent
data) follows an identical process to those already
presented for instances of \textit{GenericMultiInputTable}. Loading data onto
the table populates the data for the zeroth-derivative variable(s):
\begin{code}
coupled_table.load_data( &data_x[0][0], {5,6});
\end{code}

Loading data onto the table's \textit{derivs} element populates the data
for the first-derivative variable(s):
\begin{code}
coupled_table.derivs.load_data( &data_xdot[0][0], {5,6});
\end{code}

The two data arrays should have identical dimensions:
\begin{itemize}
\item the first number (5 in the above examples) is the number of variables
      to be populated; because each variable has a derivative, the number of
      variables to be populated by each table must match.
\item the second numer (6 in the above example) is the number of reference
      points values in the lookup data. Because each variable uses the same
      independent variable, these must also be the same number for both tables.
\end{itemize}

\subsection{Optional Configurations}
There are two optional configuration flags associated with an instance of
\textit{SingleInputTableVarDeriv}:
\begin{itemize}
  \item \textbf{use\_linear\_interpolation} allows the user to switch between
  using decoupled linear interpolation and the coupled cubic interpolation on
  both the primary and derivative variables. This allows easy comparison
  between using the two available interpolation algorithms, and for switching
  between the two algorithms on the fly.
  \item \textbf{omit\_derivative\_value} allows the use of this class for
  obtaining only the interpolated values of the primary variable. This supports
  the loading of the derivative variable reference points, and the use of those
  reference points in interpolating the primary variable, without needing to
  also interpolate the derivative values.
\end{itemize}
Both flags default to false.


\section{TableLookupSet with Transpose Data}\label{sec:UG_transpose}

To illustrate the meaning of transpose data sets, consider two examples
where we need to extract data for multiple variables as a function of a
single input variable:
\begin{enumerate}
  \item The temperature, density, and pressure are calibrated at different
  locations within an atmosphere.
  \item the 3-element position vector is calibrated as a function of time
  for some trajectory.
\end{enumerate}

Consider the first example, with variables that are conceptually
independent.  Data for such a situation might be expected to be presented as
${(T_1, T_2, T_3, ...); (\rho_1, \rho_2, \rho_3, ...); (P_1, P_2, P_3, ...)}$
with each variable given its complete set of values before moving onto the
next variable.

However, in the second example, we expect to see data presented as:
${(x_1, y_1, z_1); (x_2, y_2, z_2); ...}$, where each calibration point is
fully described before moving on to the next calibration point.

When a table-lookup extracts data from its internal array, it must have some
concept of how the data are structured.  When the data are loaded into a
table, the table only knows the number of entries and the value of each
entry.  At some point, it must break those data down into one of the two
presentation formats.  The default format is the first one -- \textbf{that
data is organized first by dependent variable, and then by calibration point}

Thus, if we pass ${x_1, y_1, z_1, x_2, y_2, z_2, x_3, y_3, z_3}$ into a
table, default behavior will result in it being incorrectly interpreted.
For the second calibrated value
of $x$, the model will return $y_1$, not $x_2$.

The \textit{TableLookupTransposeDataSet} class is a type of
\textit{TableLookupSet}; it is used where data is presented in a transpose
sense to that assumed for a \textit{TableLookupSet}.

When we provide transpose data, such as
${x_1, y_1, z_1, x_2, y_2, z_2, x_3, y_3, z_3}$,
the \textit{TableLookupTransposeDataSet} can be used to process
these data into the default format,
${x_1, x_2, x_3, y_1, y_2, y_3, z_1, z_2, z_3}$,
at which point it functions
equivalently to a \textit{TableLookupSet}.

\subsection{Instantiate}
There are four necessary components to make the
\textit{TableLookupTransposeDataSet} function correctly:
\begin{enumerate}
 \item The table manager; in this case, this would be an instance of
       \textit{TableLookupTransposeDataSet}.
 \item One and only one \textit{TableIndependentVariable}. This may be a
       standalone instantiation or, like the \textit{SimpleTableLookup}, the
       \textit{TableLookupTransposeDataSet} will attempt to create one if one
       has not been provided.
 \item At least one data table (e.g. a \textit{GenericMultiInputTable} or
       derivative thereof).
       Often, the model uses only one table.  Remember that each table can
       handle multiple dependent variables when:
       \begin{itemize}
        \item the dependent variables share a common set of independent
        variables (in this case a single common independent variable), AND
        \item the dependent variables share the interpretation of the
        independent variable(s)
        (i.e. all use interpolation, or all use round-to-the-nearest, etc.),
        AND
        \item the dependent variables have the same type (i.e. linear,
        angular, quaternion variables)
       \end{itemize}
       Note - Because the model can only support
       one independent variable, the requirement on matching independent
       variables is met automatically; the other two requirements may
       impose the presence of more than one table.

       There is an additional requirement
       with this sub-model -- that the data used to populate each table must be
       from a contiguous block of the original data array.  So if the desired
       variables are found in the 1st and 3rd columns, then either:
       \begin{itemize}
        \item and a redundant variable to take the 2nd column, or
        \item assign the two columns to two separate tables (which can both be
        managed by the \textit{TableLookupTransposeDataSet} manager instance).
       \end{itemize}

       It is recommended -- but not required -- that the instances of the data
       tables are pre-instantiated and addressed via the table-configuration
       specification (see next item). If a table-configuration does not point
       to a pre-instantiated table instance, the
       \textit{TableLookupTransposeDataSet} instance will
       attempt to create the necessary table(s) when it processes the data (see
       section \ref{UG:transpose_data_set:populate}).

       \textbf{Caution} -- creating an instance of
       \textit{SingleInputTableVarDeriv} will also require populating the
       data for the first-derivative variable(s). Creating and populating an
       instance of \textit{GenericSingleInputTable} for the first-derivative
       will not automatically connect it to the derivs member in the instance
       of \textit{SingleInputTableVarDeriv}. It is strongly recommended that,
       when planning to use a \textit{SingleInputTableVarDeriv}, a
       pre-existing instance should be used; both the inherent and embedded
       instances of \textit{GenericSingleInputTable} can then be configured
       independently without the difficulty of trying to connect them
       afterwards.

.
 \item At least one instance of the configuration instructions,
       \textit{TableLookupTransposeDataSet\_\\TableConfig}. This table
       configuration specifies how to break out the transpose data into a
       set that can be read by the \textit{GenericMultiInputTable}.

       Each table requires its own configuration.  However, the architecture
       allows one instance of \textit{TableLookupTransposeDataSet\_TableConfig}
       to provide configurations to multiple tables.  For details, see section
       \ref{sec:UG_transpose_dependent}.
\end{enumerate}


\subsection{Construct}
The \textit{TableLookupTransposeDataSet} and
\textit{TableLookupTransposeDataSet\_TableConfig} constructors \\take no
arguments.  The other classes have already been covered.

\subsection{Configure Independent Variable}
Like the \textit{SimpleTableLookup}, this class supports only a single
independent variable.  This can be provided as an instance of a
\textit{TableIndependentVariable}, or the model will create such an instance
and connect it to the specified \textit{independent\_var} variable.

The configuration of a stand-alone instance of
\textit{TableIndependentVariable} has already been explored extensively.

To use the latter option, populate two fields:
\begin {enumerate}
\item \textit{TableLookupTransposeDataSet::independent\_var} pointer to
       access the intended variable.
\item \textit{TableLookupTransposeDataSet::indep\_continuity} to specify
       whether the independent variable domain is bounded (default) or wraps
       around indefinitely.
\end{enumerate}

Example:
Have the model create a new \textit{TableIndependentVariable} with wrap-around
domain.

\begin{code}
TableLookupTransposeDataSet manager;
double                      independent;
...
manager.independent_var = &independent;
manager.indep_continuity = TableIndependentVariable::WrapAround
\end{code}

Note - if a \textit{TableIndependentVariable} instance is added to the manager
before initialization, that instance will take precendence and the model will
not create a new instance.

\subsection{Configure Dependent Variables}\label{sec:UG_transpose_dependent}
Each column of data in the transpose-data-set \textit{potentially} represents a
dependent variable.  This section describes how to configure the model to get
the data out of the transpose-data-set and into a
\textit{GenericMultiInputTable}.

\textbf{Important note} - it is \textbf{NOT} required that all columns of a
transpose-data-set be used.  Users have the flexibility to ignore any columns
of irrelevant data.

The configuration of how to extract the data for the dependent variables is
contained in the \textit{TableLookupTransposeDataSet\_TableConfig} class.
Each table (\textit{GenericMultiInputTable}) requires its own configuration
(\textit{TableLookupTransposeDataSet\_TableConfig}).

However, the architecture allows one instance of
\textit{TableLookupTransposeDataSet\_TableConfig} to provide configurations
to multiple tables because when these configurations are loaded onto the
manager, their data is replicated.  This allows the instance to be
re-populated and re-pushed onto the manager indefinitely.

Alternatively, it is also acceptable to create one configuration per table.

Each configuration specifies:
\begin{itemize}
 \item The lookup-method that the table will use to interpret the
       independent variable.
       \begin{itemize}
        \item \textbf{TableIndependentVariable::LookupMethod lookup\_method}
              This defaults to interpolation (
              \textit{TableIndependentVariable::Interp}).
       \end{itemize}
 \item A selection of columns of data where each column
       represents a dependent variable:
       \begin{itemize}
       \item If the data columns are contiguous, it is only necessary to
             specify the first and last column number:
             \begin{itemize}
              \item \textbf{size\_t index\_front} provides the low column index
                    (note the first column has index 0, the second index 1, etc.)
              \item \textbf{size\_t index\_back} provides the high column index.
             \end{itemize}
             The selected columns then are from \textit{index\_front} to
             \textit{index\_back} inclusive.
       \item If the data columns are not contiguous, the column numbers can be
             specified in a STL-vector:
             \begin{itemize}
              \item \textbf{std::vector<size\_t> indices} provides the list of
                    column indices.
                    (note the first column has index 0, the second index 1, etc.)
             \end{itemize}
             This option may also be used for contiguous data columns, but each
             column must then be specified. Because of the limitations of SWIG,
             this option may not be available to specify in a Python input
             file.
       \end{itemize}
 \item Which \textit{GenericMultiInputTable} to populate with this
       range.  If unspecified, a new table will be created.
       \begin{itemize}
        \item \textbf{GenericMultiInputTable * table\_ptr} (optional)
       \end{itemize}
 \item If creating a new table, the type of table to be created and
       a list of the dependent variables associated with each column.
       \begin{itemize}
        \item \textbf{AbstractTableLookup::TableType table\_type} (optional)
              This value defaults to AbstractTableLookup::Generic which results
              in the creation of a \textit{GenericMultiInputTable}.
        \item \textbf{std:vector$<$double *$>$ dependent\_variables} (optional)
              is a vector of pointers to the dependent variables.
       \end{itemize}
 \item An optional name; this is most useful to obtain access to an instance
       after it has been loaded into the manager.  See below for details.
\end{itemize}


Example:
This code snippet will result in the creation of a new table because none is
specified; the new table will default to a \textit{GenericMultiInputTable}
because no type is specified.  The lookup-method is set to round-to-nearest.
The new table will be populated with data taken from columns 4-6 inclusive.
\begin{code}
TableLookupTransposeDataSet_TableConfig table_config;
...
table_config.lookup_method = TableIndependentVariable::Round;
table_config.index_low = 3;
table_config.index_high = 5;
\end{code}

\subsection{ Pushing and Accessing Configurations}
Once a \textit{TableLookupTransposeDataSet\_TableConfig}  instance has been
configured, it can be pushed onto the manager with the
\textit{add\_config(...)} method.
\begin{code}
TableLookupTransposeDataSet             manager;
TableLookupTransposeDataSet_TableConfig table_config;
...
table_config.lookup_method = TableIndependentVariable::Round;
table_config.index_low = 3;
table_config.index_high = 5;
manager.add_config( table_config);
\end{code}

Note that changing the values of a configuration after it has been loaded
will have no effect because the values are \textit{copied} into the manager.
To change values after loading, it is necessary to find the loaded instance
and change those values.
There are two access methods provided to support ``undoing'' the pushing of
\textit{TableLookupTransposeDataSet\_TableConfig} data onto the manager.
\begin{itemize}
\item \textbf{get\_config} returns a pointer to the desired instance of the
      \textit{TableLookupTransposeDataSet\\\_TableConfig} as it resides in the
      manager.  It has two signatures, both using the name of the
      configuration:
      \begin{itemize}
      \item \textit{get\_config(const char * name)}
      \item \textit{get\_config(const std::string \& name)}
      \end{itemize}
\begin{code}
table_config.index_low = 3;
table_config.index_high = 5;
table_config.name = "my_config";
manager.add_config( table_config);
...
// Change index_high to 4
TableLookupTransposeDataSet_TableConfig * table_correction
     = manager.get_config( "my_config");
table_correction->index_high = 4;
\end{code}
\item \textbf{remove\_config} can be used to completely remove a configuration
      from the manager.  This takes a reference to an instance of the original
      \textit{TableLookupTransposeDataSet\_TableConfig} and searches the
      loaded copies to find the one with matching parameters, then removes it
      compeletely.  The original instance can now be corrected and re-pushed if
      desired, or left absent.
      \begin{itemize}
      \item \textit{remove\_config(TableLookupTransposeDataSet\_TableConfig
             \& config)}
      \end{itemize}
\begin{code}
table_config.index_low = 3;
table_config.index_high = 5;
manager.add_config( table_config);
...
// Change index_high to 4
manager.remove_config(table_config);
table_config.index_high = 4;
manager.add_config( table_config);
\end{code}
\end{itemize}

\subsection{Populate}\label{UG:transpose_data_set:populate}
\textbf{Independent Variable Table}

The independent-variable data may be populated independently (default), or it
may be extracted from the transpose-data-set.  For example, suppose the
transpose-data-set comprises 2 columns - time and distance.  We want to get
distance as a function of time.  We can extract the time values directly from
the data-set and use these to populate the independent variable.

Stand-alone population of the independent variable has already been covered
in this User's Guide.

Using the data-set tp populate the independent variable requires the setting
of 2 variables:
\begin{itemize}
 \item \textbf{populate\_independent\_from\_file} defaults to false; set this
       to true to extract data from the data-set.
 \item \textbf{indep\_index} is the index of the column that contains the
       independent-variable data.  The first column has index 0.
\end{itemize}

These data will then be populated with the execution of the
\textit{process\_data(...)} method (see below).

\textbf{Dependent Variables Table}

Beneath the surface, the operation of the \textit{TableLookupTransposeDataSet}
is largely equivalent to that of the \textit{TableLookupSet}.  The difference
comes in the ability to extract the data from the transpose-data-set and
populate the regular tables with those data.  So far, we have looked at the
configuration of how to do that; the remaining step is the execution.

The model provides two methods to process a transpose-data-set:
\begin{itemize}
\item Takes a filename and an (optional) line-limit.  The line-limit allows the
user to truncate the transpose-data-set at a specified line.
\begin{code}
  process_data( const std::string & filename,
                size_t terminate_on_line = 100000);
\end{code}

\item Takes a reference to an instance of
      \textit{std::list$<$ std::vector $<$ double $> ~ >$} (which is a list of
      vectors of type-double).  Each line of the transpose-data-set is
      assigned to a STL-vector, and each such vector added to a STL-list.

      If you are interested, the architecture here is deliberate.  The lines
      are always processed sequentially, so a STL-list of these is sufficient.
      However, it may be necessary to pick out individual values from each
      line, so each line must be a vector to facilitate indexing.
\begin{code}
  process_data( std::list< std::vector < double >  >)
\end{code}

Note that this option always runs.  If the first option is called; those data
will be read in from the specified file, converted to a list-vector data
structure, and this option will then be executed.

\end{itemize}



\subsection{Connect and Initialize}
The tables and the independent variables are connected automatically within the
\textit{process\_data()} method.

The \textit{TableLookupTransposeDataSet} must be initialized; the
initialization method is inherited from \textit{AbstractTableLookup} so matches
that for \textit{TableLookupSet}, see section \ref{sec:UG_multi_single}.
.

\subsection{Execute}
As for \textit{TableLookupSet}, see section \ref{sec:UG_multi_single}.

\section {Multi-set Configurations}
The usage of \textit{TableLookupSet} is intended to be an all-encompassing
manager for all variable lookups at any given time.  It can handle multiple
tables simultaneously, with cross-over usage of shared independent variables.

However, there are situations in which the selection of which variables to use,
or which variables to populate, or the domain of the independent variables
is dependent on the model status.  In these unusual cases, it may be
necessary to implement multiple instances of \textit{TableLookupSet}.
While it is possible to have
these multiple instances running concurrently, the intended
architecture is to switch between instances, each of which provides the
comprehensive set of tables for the current model status.

When it is time to switch from one \textit{TableLookupSet} to another, the
\textit{subscribe()} and \textit{unsubscribe()} methods can be used to
manage the activity of the set.  See section \ref{sec:UG_tab_mgr} for
details.

Note that this activity-management capability applies only to the
manager classes.  Individual tables within an active manager will always be
active.

Any single instance of any table -- \textit{GenericMultiInputTable}, or
\textit{SingleInputTableFor*} --
can be loaded into multiple instances of
\textit{TableLookupSet}.  The \textit{TableLookupSet} simply maintains a list
of pointers to these tables, which exist as independent entities.

\textbf{Example}: consider the simple case in which a single variable
 depends on a
pair of independent variables.  For one of those independents, the dependency
is constant in time; for the other, there is some threshold event that triggers
a different response as the vehicle shifts from one phase to another.

In this case, we would need 3 instances of TableIndependentVariable - one for
the consistent independent and two for the independent with two distinct
phases.
\begin{code}
double                     independent_a;
double                     independent_b;
double                     dependent;

TableIndependentVariable   ind_var_a;
TableIndependentVariable   ind_var_b_phaseA;
TableIndependentVariable   ind_var_b_phaseB;
\end{code}

We also need 2 instances of TableLookupSet, and 2 tables, one for each phase.
\begin{code}
TableLookupSet             phaseA_set;
TableLookupSet             phaseB_set;
GenericMultiInputTable     phaseA_table;
GenericMultiInputTable     phaseB_table;
\end{code}

Construction is as outlined in section \ref{sec:UG_multi_single}.
\begin{code}
phaseA_set();
phaseB_set();
ind_var_a(independent_a);
ind_var_b_phaseA(independent_b);
ind_var_b_phaseB(independent_b);
phaseA_table(dependent);
phaseB_table(dependent);
\end{code}

Each table gets populated with its own data following the instructions
in section \ref{sec:UG_multi_single}.  The following conceptual-code
illustrates the process.
\begin{code}
double scratch_ind_a[ Num_a ] = {<data>};
ind_var_a.load_data( scratch_ind_a, Num_a);

double scratch_ind_bA[ Num_bA ] = {<data>};
ind_var_b_phaseA.load_data( scratch_ind_bA, Num_bA);

double scratch_ind_bB[ Num_bB ] = {<data>};
ind_var_b_phaseB.load_data( scratch_ind_bB, Num_bB);

std::vector<size_t> size_vec();
size_vec.push_back(1);
size_vec.push_back(Num_a);
size_vec.push_back(Num_bA);
double scratch_table_A[ Num_a ][ Num_bA ] = ;
phaseA_table.load_data( &scratch_table_A[0][0],size_vec);

size_vec[2] = Num_bB;
double scratch_table_B[ Num_a ][ Num_bB ] = ;
phaseB_table.load_data( &scratch_table_B[0][0],size_vec);
\end{code}

The tables and independent variables are added and connected within
the respective
\textit{TableLookupSet} as outlined in section \ref{sec:UG_multi_single}.
Note that in this example, the common \textit{ind\_var\_a} gets added to both
sets, whereas each version of the second variable only gets added to the set
associated with that phase.
\begin{code}
phaseA_set.add_table( phaseA_table);
phaseA_set.add_independent_variable( ind_var_a);
phaseA_set.associate_table_and_independent();
phaseA_set.add_independent_variable( ind_var_bA);
phaseA_set.associate_table_and_independent();

phaseB_set.add_table( phaseB_table);
phaseB_set.add_independent_variable( ind_var_a);
phaseB_set.associate_table_and_independent();
phaseB_set.add_independent_variable( ind_var_bB);
phaseB_set.associate_table_and_independent();
\end{code}

Both sets must be initialized.
\begin{code}
phaseA_set.initialize();
phaseB_set.initialize();
\end{code}

Both sets may be configured independently - see section
\ref{sec:UG_multi_single}.

Both sets may be called as direct scheduled updates:
\begin{code}
phaseA_set.update();
phaseB_set.update();
\end{code}

The final piece of the puzzle is in activating and deactivating the two sets.
Start by switching on only the the phase-A table.
Without a subscription to activate it, the \textit{phaseB\_set.update()}
 execution will return immediately without processing any data.
\begin{code}
phaseA_set.subscribe();
\end{code}
At some time later, the two tables can be switched:
\begin{code}
phaseA_set.unsubscribe();
phaseB_set.subscribe();
\end{code}

\textbf{Note:}
A viable alternative to activating/deactivating the table-sets is to wrap the
table-set instances in another class, and have that class's \textit{update()}
method call the \textit{update()} method on only the intended table-set.
But that requires more code-writing.




\section{Creating Tables On-the-fly}
Thus far, we have been exclusively concerned with pre-configuring a scenario
at the top-level.  All of that functionality is also available on-the-fly, with
the flexibility of creating tables at run-time.

The methods presented in this section are intended for user-access to
support \textit{TableLookupSet} configurations. However, they are found in
the \textit{AbstractTableLookup} class.  The reason for this is that they are
used behind the scenes when creating a \textit{SimpleTableLookup}
implementation (\textit{SimpleTableLookup} is a type of
\textit{AbstractTableLookup}).  When using a \textit{SimpleTableLookup}, it
should not be necessary for a user to access these methods directly.

\subsection{Instantiate and Populate}
To create instances of the components of \textit{AbstractTableLookup}
on-the-fly, use the methods:
\begin{itemize}
 \item \textit{AbstractTableLookup::create\_independent\_variable(...)},
 \item \textit{AbstractTableLookup::create\_table(...)}
\end{itemize}

\subsubsection{Independent Variable}
The \textit{create\_independent\_variable(...)} takes two or three arguments:
\begin{enumerate}
 \item std::string \textit{name} to identify the instance.
  This must be unique.
 \item Reference to the actual variable being used as the independent variable.
 \item (optional) Continuity type of the independent variable.
  If unspecified, defaults to \textit{Table-IndependentVariable::Linear}.
  May alternatively be
  specified as \textit{TableIndependentVariable-::WrapAround}
  to create a wrap-around independent variable with infinite domain.
\end{enumerate}

 It creates the
 instance of \textit{TableIndependentVariable} and returns a pointer to that
 instance.  That pointer should be used to register the new instance with the
 \textit{TableLookupSet} instance via the
 \textit{TableLookupSet::add\_independent\_variable} method.

 This architecture is somewhat fragile.  Unless the instance is registered
 (or the address of the new instance stored elsewhere), access to it can
 easily be lost if the returned pointer is overwritten or falls out of scope.
 In those circumstances, the instance will continue to exist, but there will
 be no user-interface to change its contents or register it.

It is recommended that data for the new instance also be populated at the
same time
that it is created while the pointer to it is still easily available.
Note that creating regular instances of \textit{TableIndependentVariable}
(rather than on-the-fly instances) altogether avoids this potential problem of
losing handles because the instance will always be available at a known
location.

An example is shown following the discussion on creating a new table.

\subsubsection{Dependent Variable / Table}
The \textit{create\_table(...)} has three distinct forms:
\begin{itemize}
 \item \textbf{Single Variable}.  This option takes 1 or 2 arguments:
   \begin{enumerate}
    \item A reference to the actual dependent variable being populated
     by this table.
    \item (optional) The type of table.  If unspecified, defaults to
     \textit{AbstractTableLookup::Generic}.  Alternatively, may be
     specified as \textit{AbstractTableLookup::AngularDeg} or
     \textit{AbstractTable-Lookup::AngularRad} to create an angular output.
  \end{enumerate}

 \item \textbf{Multi Variable, multi-argument}.  This option takes a number of
   arguments equal to the number of variables plus 1.  They are:
  \begin{enumerate}
   \item An integer specifying the number of variables.
   \item A series of pointers, one to each of the variables.
  \end{enumerate}
  \textbf{WARNING}:  Be careful with this option.  If the integer
   specifying the number of variables does not match with the number of
   variables, bad things can happen.  If the integer is too small, the latter
   variables will simply not get included.
  \textbf{If the integer is too large},
   erroneous addresses will be introduced, which can lead to a
   \textbf{Segmentation Fault}.

  Note - The dimensions of the loaded data will be tested against the
  specified integer, not the number of arguments.

 \item \textbf{Multi Variable, STL-vector}.  This option takes 1 or 2
   arguments:
   \begin{enumerate}
    \item A reference to an STL-vector containing pointers to all of the
    dependent variables being populated by this table.
    \item (optional) The type of table.  If unspecified, defaults to
    \textit{AbstractTableLookup::Generic}.  May alternatively be specified as
    \textit{AbstractTableLookup::AngularDeg} or
    \textit{AbstractTable-Lookup::AngularRad} to create an angular output.
  \end{enumerate}
\end{itemize}

\subsubsection{Example}


\begin{code}
%*\begin{codecomment}
Create a new \textit{TableIndependentVariable}
\end{codecomment}*)
TableIndependentVariable * new_var =
    table_set.create_independent_variable(
                                'indy_0',
                                my_independent_variable0);
%*\begin{codecomment}
Create the data and load it onto the newly-created
\textit{TableIndependentVariable} while we still have easy access to the new
instance.

\end{codecomment}*)
  double scratch[5] = {1,2,3,5,8};
  new_var->load_data(scratch, 5);
%*\begin{codecomment}
Create additional instances of  \textit{TableIndependentVariable}.
The same local pointer may be re-used to temporarily store the access
to the new instance, but realize that this overwrites the address of the first,
making it more difficult to access that one.  Alternatively, a new local
pointer could be created for each new instance.
\end{codecomment}*)
  new_var = table_set.create_independent_variable(
                               'indy_1',
                               my_independent_variable1);
%*\begin{codecomment}
Continue to populate this instance and create others
\end{codecomment}*)
...
%*\begin{codecomment}
In the same way, create and populate the tables one at a time:
\end{codecomment}*)
  GenericMultiInputTable * new_table =
     table_set.create_table( 2,
                            &my_dependent_variables0,
                            &my_dependent_variables1);
  double temp_table[2][5][5][5] = {{{{...}}}};
%*\begin{codecomment}
Load the data as described in the section \ref{sec:UG_multi_single}.
\end{codecomment}*)
  std::vector<size_t> size_vec;
  size_vec.push_back(2);
  size_vec.push_back(5);
  size_vec.push_back(5);
  size_vec.push_back(5);
  new_table->load_data(&temp_table[0][0][0][0],size_vec);
%*\begin{codecomment}
Create and populate a single-variable table:
\end{codecomment}*)
  new_table = table_set.create_table(
                           my_dependent_variables[2]);
  ...
\end{code}


\subsection{Connect}
The tables and independent variables must be connected with the
\textit{associate\_table\_and\_independent} method whether they are created
 statically or dynamically (on-the-fly).
 See section \ref{sec:UG_multi_single}.

Important Limitation:  The specific recognition of the table
 (e.g. \&\textit{table0}) and independent variable (e.g. \&\textit{var1})
 presented in the examples in section \ref{sec:UG_multi_single} requires
 that these addresses are known.  There are 3 options still available for
 providing these identifying arguments:
\begin{itemize}
 \item Use the pointers returned by the \textit{create\_*} methods instead
       of taking the addresses of statically allocated elements.
 \item Use names, as also described in section \ref{sec:UG_multi_single}.
 \item Use the indices of the location of the target in the table-set's lists
       This is a `last resort'; it is not recommended because it results in
       an implementation that is very difficult to maintain.
\end{itemize}


\subsection{Initialize}
See section \ref{sec:UG_multi_single}.

\subsection{Configure (optional)}
See section \ref{sec:UG_multi_multi}.

\subsection{Execute}
See section \ref{sec:UG_multi_single}.


\section{Modifying Default Settings}\label{sec:UG_bias_scale}

In most cases, the data used for table lookups and interpolations is fixed
from scenario to scenario,  Throughout the User's Guide, the examples we
have seen of how to establish and populate these tables have primarily
used compiled code.  While it is more efficient to organize large data
sets through compiled code, doing so does come with a significant downside
-- what happens if the data needs to be modified for a particular scenario.

An example of such a scenario might be using a table to generate a signal
from a particular piece of hardware; a particular requirement may call for
an analysis of the effect of introducing a bias in the hardware.
This could be resolved by writing a pre-processing or post-processing
algorithm to bias the nominal inputs or outputs (respectively), or by
biasing the table data itself.  But if the table data is compiled and
locked-in, modifying it means editing the code and recompiling -- not
a task well suited to a fast turnaround of a simple analysis.

To overcome this problem, the model provides some methods for modifying
values in the data tables for both independent variables (inputs) and
dependent variables (outputs).

\subsection {Complete Overwrite of Data}

If using the \textit{SimpleTableLookup}, data tables for independent and
dependent data may be replaced completely in mid-run using the
\textit{load\_dependent\_data(...)} and \textit{load\_independent\_data(...)}
methods. See Section \ref{sec:UG_single_multi} for an introduction to these
methods, including instructions on calling them.
See Section  \ref{sec:UG_reloading_data} for verification of this data-reload
capability.

If using a \textit{TableLookupSet}, overwriting data arrays is not recommended.
In this case, the \textit{GenericMultiInputTable} and
\textit{TableIndependentVariable} instances contain the dependent and
independent data, and the independence of these two entities makes the process
of safely overwriting data substantially more challenging.
When the dependent data is loaded, the number of data elements is checked
against the specified dimensions. Replacing the
dependent data in the \textit{GenericMultiInputTable} can be achieved using the
\textit{GenericMultiInputTable::load\_data(...)} method \textbf{if} the new data
`array' has the same dimensionality as the old. However, if the dimensionality
of the new dependent data array changes, then the independent arrays on which
it depends must be changed as well to maintain consistency.

Similarly, re-loading independent data is feasible \textbf{if} the new data has
the same number of elements as the old data, and \textbf{if} the new data is
monotonic. Changing the independent data has the potential for additional
consequences in that
each dependent data \textit{GenericMultiInputTable} instance in the
\textit{TableLookupSet} uses some subset of the available
\textit{TableIndependentVariable} instances in the same \textit{TableLookupSet}.
If it is desired to change the data in an independent variable with the
intention of changing the way a particular dependent data table responds, it
must also be considered that every other dependent data table using that
independent variable will be similarly affected. An additional step is
required -- a call to \textit{clear\_data()} -- before overriding any data
elements in the independent data array, then the array can be reloaded.

In general though, unless the user is confident that the sizes of the arrays
are not being changed, making complete overwrites of table data when using a
\textit{TableLookupSet} \textbf{IS NOT RECOMMENDED}.
The consequences of getting it
wrong are seen in \textit{SIM\_01a\_table\_lookup\_set\_nominal/}
\textit{ERROR\_reload}.
In this example, the dependent data and independent data are succesfully
replaced with values using the same sized arrays. Then the independent data is
replaced with a longer array than the dependent data is expecting. The result
is that in the process of evaluating the independent variable against
its calibrated values, an index is generated that exceeds the dimensionality of
the dependent array. Data is still extracted from the assigned index, but it is
nonsensical data because the array does not extend into that memory space.

Instead of relying on the risky process of replacing the data arrays, when
using \textit{TableLookupSet} the model provides the ability to modify existing
data in both the dependent and independent arrays with methods to bias and
scale the existing values. These two processes may be used in conjunction, and
cannot result in changing the size of the arrays. They also include safety
checks on monotonicity that would be bypassed by a complete replacement of
existing data.

Note -- because these methods are embedded in the
\textit{GenericMultiInputTable} and \textit{TableIndependentVariable} classes,
they are conceptually available to users of \textit{SimpleTableLookup} as well
as users of \textit{TableLookupSet} because \textit{SimpleTableLookup}
internally uses these classes. However, the user interface to access these
methods in the \textit{SimpleTableLookup} is non-trivial for the casual user so
it is not explored in this document.

\subsection {Biasing the Values of the Independent Variable}

The set of precalibrated data points used when defining the independent
variable can be biased across its complete domain, or across some contiguous
subset by adding (or subtracting) a fixed value to the chosen data values.

The \textit{bias\_data} method can be called at any time to modify the data
held by an instance of \textit{TableIndependentVariable}.  It generally takes
3 arguments:
\begin{itemize}
\item the bias to apply
\item the lowest index of calibrated values to which the bias should be
applied
\item the highest index of calibrated values to which the bias should be
applied
\end{itemize}
The bias provided in the first argument will then be applied to all values at
indices between the provided limits, \textbf{including at both specified
indices}.

In the case where the bias is to be applied to only 1 value, an equivalent
two-argument method can be called:
\begin{itemize}
\item the bias to apply
\item the index of calibrated values to which the bias should be applied
\end{itemize}
This two-argument version simply calls the three-argument version behind the
scenes using the provided second argument as both the second and third argument.

In the case where the bias is to be applied uniformly across all data points, an
equivalent one-argument method can be called:
\begin{itemize}
\item the bias to apply
\end{itemize}
This one-argument version simply calls the three-argument version behind the
scenes, using 0 for the second argument and the array-size (minus 1) for the
third.


An important requirement to consider when planning application of biases to a
subset of data is that after each application, the resulting data
\textbf{must be monotonic}.  If a process results in non-monotonic data,
\textbf{it will be rejected}.  For example, biasing the first two elements of
{0,1,2} by +2 would result in {2,3,2}.  This application would be rejected
because the resulting data is not monotonic.

An example of these methods can be seen in
\textit{SIM\_01a\_table\_lookup\_set\_nominal/RUN\_biases}
\begin{code}
# apply bias of +1.0 to indices 0,1,2:
testing_model.test_so_defined.indep1.bias_data( 1.0,0,2)

# apply a bias of -2.0 to index 0 only:
testing_model.test_so_defined.indep1.bias_data(-2.0,0)
\end{code}

Further examples of application bias are found in
\textit{SIM\_05\_independent\_wraparound/RUN\_bias}.  This run tests the
detection of certain error cases when applying an invalid bias.


\subsection {Scaling the Values of the Independent Variable}

Just as the set of precalibrated data points may be biased across some
subset or all of the domain, a similar mechanism allows for the application
of a scaling factor to a subset of data point.

The \textit{scale\_data} methods have the same signatures as the
\textit{bias\_data} methods (outlined above).


Examples of a scaling factor are found in
\textit{SIM\_05\_independent\_wraparound/RUN\_bias}.  This run also tests
certain error cases when applying an invalid bias or scaling factor.

\subsection {Compounding Bias and Scale}
The biases and scaling factors are applied in the order in which they are
commanded.  Each application is limited to a contiguous block of data, but
there is no limit to how many times either process may be applied.
Because of the monotonicity requirement, achievement of some modifications may
require compound operations.

As a simple example, biasing {0,1,2} into {2,1,0} by applying a bias of 2
to the first element and -2 to the third element cannot be accomplished
by biasing alone.  It does not matter whether the +2 or -2 bias is applied
first, or whether the bias is broken down into smaller pieces, at some
point one of the two modified points will have to move across the value 1,
at which point both ends will be on the same side of 1, and the data will
not be monotonic.  It will be rejected.  However, this outcome can be
achieved with the application of a scaling by -1 followed by a bias of
+2 to all three points.


\subsection {Applying Bias or Scale to Table Data}
Just as the independent variable data can be modified as presented above,
so can the
dependent data.  The same \textit{bias\_data} and \textit{scale\_data}
signatures are
also available to the \textit{GenericMultiInputTable} class.  Note that
this capability
is not available to instances of \textit{SimpleTableLookup}.


While the methods for modifying independent and dependent data are identical in
appearance, users should be aware of a complicating factor when modifying
a subset of dependent data, that of the dimensionality of the data being
modified.  The independent variables use a one-dimensional array,
so identifying the index or
indices for application of the intended modifications is a simple undertaking.
Conversely, the dependent data is often multi-dimensional, at least in
concept.  For example, the following table illustrates the desired output
for a lookup table as a function of two variables:

\begin{table}[h!]
\caption{A Simple Lookup Table}
\begin{center}
\begin{tabular}{||c||c|c|c||}
\hline \hline
 & \multicolumn{3}{c}{independent A values} \\
 & \textbf{0} &\textbf{1} &\textbf{2}\\
independent B values & & & \\
\hline
\textbf{0} &  0 &  1 &  2 \\
\textbf{1} & 10 & 11 & 12 \\
\textbf{2} & 20 & 21 & 22 \\
\hline \hline
\end{tabular}
\end{center}
\end{table}

Applying a modification to the whole table is straightforward with use of the
1-argument method.  For example,
\begin{code}
table.bias_data(0.5)
\end{code}
creates:
\begin{table}[h!]
\begin{center}
\begin{tabular}{||c||c|c|c||}
\hline \hline
 & \multicolumn{3}{c}{independent A values} \\
 & \textbf{0} &\textbf{1} &\textbf{2}\\
independent B values & & & \\
\hline \hline
\textbf{0} &  0.5 &  1.5 &  2.5 \\
\textbf{1} & 10.5 & 11.5 & 12.5 \\
\textbf{2} & 20.5 & 21.5 & 22.5 \\
\hline \hline
\end{tabular}
\end{center}
\end{table}

However, identifying a single index can be more problematic.
Suppose we wanted to modify only the value $10$ to be $10.5$.
The first point to realize is that while this
data-set is conceptually two-dimensional, it is still stored in one
dimension. It is essential to know not only which data point to modify,
but how to locate that data point.  Continuing with the previous example,
these data could have been entered as
\begin{code}
double data[3][3] = {{0,1,2},{10,11,12},{20,21,22}};
table.load_data(&data[0][0], ...);
\end{code}
or as the transpose:
\begin{code}
double data[3][3] = {{0,10,20},{1,11,21},{2,12,22}};
table.load_data(&data[0][0], ...);
\end{code}
Neither is any more correct than the other.  Knowing
whether the value $10$ is at index $1$ or $3$ requires knowledge of how
the data were
originally loaded into the table.  It should not be assumed
that the way in which the
data are presented visually matches the way in
which they were entered in code.

Applying a modification to a subset of data can be even more challenging.
As we presented when discussing modifications to the independent variables,
a contiguous block of data may be modified using the three-argument method,
with the indices of the two ends of the contiguous block being passed as the
second and third arguments.  However, with multi-dimensional data,
data values that \textit{appear} contiguous may not always be so.
For example, in the table above, the
values ${10,11,12}$ appear contiguous, but whether they can be treated
as such depends on how the data was loaded.
\begin{itemize}
\item In the first example, the three values are contiguous, at indices
$3-5$ inclusive and
can be modified in a single three-argument call.
\item In the second example, the three values are not contiguous and must
be modified
with three separate two-argument calls, for indices 1, 4, and 7.
\end{itemize}

An example of the application of a single point bias to dependent data can
be found in\\
\textit{SIM\_01a\_table\_lookup\_set\_nominal/RUN\_biases}:
\begin{code}
testing_model.test_so_defined.single_table.bias_data(1000.0,14)
\end{code}



\chapter {Verification and Validation}

The model has a comprehensive set of unit-simulations used for verification of
the configurations and algorithms described previously in this document.
There are two distinct types of verification test:
\begin{itemize}
 \item Detection of errors and misconfigurations
 \item Confirmation of expected results.
\end{itemize}

In the first case, these
tests address configuration checks and the likes -- testing whether the model
identifies misconfigurations, provides suitable error messages in response,
and takes the appropriate action, up to and including termination of the
simulation.  These tests tend to be self-explanatory and are not discussed
further in this document.

In the second case, the model is expected to produce meaningful output data.
These data are then compared against data that have been manually inspected,
verified, and committed to the repository.  These test cases are now
considered, including:
\begin{itemize}
 \item Presentation of the purpose and construction of each test
 \item Presentation of the controlling data, and where appropriate
 discussion of the rationale for selection of the specific data values
 \item Derivation of the expected results from the tests
 \item Confirmation that the committed data matches the expected results.
\end{itemize}

The verification test suite is divided into ten simulations, each of which
has between 1 and 15 runs, as listed below.

\begin{itemize}
  \item SIM\_01a\_table\_lookup\_set\_nominal:
  \begin{itemize}
    \item RUN\_biases
    \item RUN\_ceil\_int\_int\_int
    \item RUN\_floor\_int\_int\_int
    \item RUN\_int\_ceil\_int\_int
    \item RUN\_int\_floor\_int\_int
    \item RUN\_int\_next\_int\_int
    \item RUN\_int\_prev\_int\_int
    \item RUN\_int\_round\_int\_int
    \item RUN\_interp\_all
    \item RUN\_interp\_all\_off\_on
    \item RUN\_next\_int\_int\_int
    \item RUN\_prev\_int\_int\_int
    \item RUN\_prev\_next\_floor\_ceil
    \item RUN\_queries
  \end{itemize}

  \item SIM\_01b\_table\_lookup\_set\_faults:
  \begin{itemize}
    \item RUN\_10\_test\_table\_manager\_init
    \item RUN\_11\_test\_table\_manager\_runtime
    \item RUN\_15\_duplicate\_name\_FAIL
    \item RUN\_16\_NULL\_dependent\_FAIL
    \item RUN\_17\_duplicate\_dependent\_FAIL
    \item RUN\_20\_test\_table\_configuration
    \item RUN\_21\_test\_table\_init
    \item RUN\_22\_test\_table\_trivial
    \item RUN\_23\_scale\_bias
    \item RUN\_25\_construct\_table\_with\_null\_dependent\_FAIL
    \item RUN\_26\_add\_NULL\_dependent\_FAIL
    \item RUN\_27\_assign\_output\_to\_empty\_FAIL
    \item RUN\_28\_assign\_output\_to\_null\_FAIL
    \item RUN\_30\_test\_independent
    \item RUN\_31\_tiv\_error\_tests
  \end{itemize}

  \item SIM\_02\_simple\_table\_lookup:
  \begin{itemize}
    \item RUN\_error
    \item RUN\_verif
  \end{itemize}

  \item SIM\_03\_simple\_table\_lookup\_for\_angles:
  \begin{itemize}
    \item RUN\_verif
  \end{itemize}

  \item SIM\_04a\_table\_lookup\_for\_quaternions:
  \begin{itemize}
    \item RUN\_error\_checks
    \item RUN\_verif
    \item RUN\_verif\_with\_negatives
    \item RUN\_with\_zeroes
  \end{itemize}

  \item SIM\_04b\_table\_lookup\_var\_deriv:
  \begin{itemize}
    \item RUN\_01a\_default
    \item RUN\_01b\_1array
    \item RUN\_01c\_2array
    \item RUN\_01d\_1vec
    \item RUN\_01e\_2vec
    \item RUN\_01f\_1elem
    \item RUN\_01g\_2elem
    \item RUN\_01h\_simple
    \item RUN\_02a\_default\_no\_xdot
    \item RUN\_02b\_default\_linear
  \end{itemize}

  \item SIM\_05\_independent\_wraparound:
  \begin{itemize}
    \item RUN\_bias
    \item RUN\_error
    \item RUN\_verif
  \end{itemize}

  \item SIM\_06\_transpose\_data\_manager:
  \begin{itemize}
    \item ERROR\_01\_reload
    \item ERROR\_02\_invalid\_lookups
    \item ERROR\_03\_irregular\_data\_file
    \item ERROR\_04\_extra\_table
    \item FAIL\_01\_short\_line\_corruption
    \item FAIL\_02\_no\_independent
    \item FAIL\_03\_double\_independent
    \item FAIL\_04\_inconsistent\_table\_configuration
    \item RUN\_01\_internal\_data
    \item RUN\_02\_external\_data
    \item RUN\_03\_external\_data\_input\_reset
    \item RUN\_04\_internal\_data\_use\_TIV
  \end{itemize}

  \item SIM\_07\_reload\_data:
  \begin{itemize}
    \item RUN\_verif
  \end{itemize}

  \item SIM\_08\_transpose\_config:
  \begin{itemize}
    \item RUN\_01\_get\_config
    \item RUN\_02\_remove\_processed
    \item RUN\_03\_remove\_config
    \item RUN\_04\_remove\_config\_2
  \end{itemize}

  \item SIM\_09\_code\_coverage:
  \begin{itemize}
    \item RUN\_abstract\_table\_lookup
    \item RUN\_generic\_multi\_input\_table
    \item RUN\_quaternion\_spherical\_interpolator
    \item RUN\_single\_input\_table\_for\_angles
    \item RUN\_single\_input\_table\_for\_quaternions
    \item RUN\_table\_independent\_variable
    \item RUN\_table\_lookup\_transpose\_data
  \end{itemize}

  \item SIM\_10\_special\_cases:
  \begin{itemize}
    \item RUN\_1\_2
    \item RUN\_high\_fraction
    \item RUN\_high\_index\_and\_new\_name
    \item RUN\_tiv\_full\_search
  \end{itemize}
\end{itemize}

\section {Configuring the Independent Variable}

This section investigates some of the ways in which the independent variable
can be configured, including:
\begin{itemize}
\item Extending its domain to infinity, using the specified domain
repetitively so that the infinte domain is periodic.
This is referred to as ``wrapping'' the domain.
\item Applying a bias to the specified calibrated values.
\item Applying a scaling factor to the specified calibrated values.
\end{itemize}

\subsection {Wrapping the Independent Variable}

Typically, an independent variable is provided with a set of
precalibrated data points.  This set must be monotonic (could be increasing
or decreasing), which means that it has two ends -- a lowest value and a
highest value.  With the default configuration for a
\textit{TableIndependentVariable} instance, when the independent variable
has a value outside this specified domain, the dependent variable takes its
value based on the end point (either left or right, depending on which side
of the domain the independent variable is beyond) of its own data array.

With the wrap-around configuration, the domain wraps back on itself
indefinitely; when the independent variable leaves its
specified domain, it reenters from the other side.

This capability is investigated in
\textit{SIM\_05\_independent\_wraparound/RUN\_verif} (tests error cases)
and in \textit{SIM\_02\_simple\_table\_lookup/RUN\_verif} (tests nominally
operating wrapped tables).  See section \ref{sec:VV_Simple_Lookup}
for details of the nominal operations of the wrapped table.


\subsection {Biasing the Values of the Independent Variable}

The set of precalibrated data points used when defining the independent
variable can be biased across some subset or all of the domain. See section
\ref{sec:UG_bias_scale} for a discussion of these processes.

This capability is investigated in
\textit{SIM\_05\_independent\_wraparound/RUN\_biases}.
In this scenario, a sequence of biases is applied.
\begin{itemize}
\item a bias is applied to the dependent data, this is presented in section
\ref{sec:VV_bias_dependent}
\item a block-application to bias the first 3 elements of {2 4 6 8} by +1.
\item A single-index application to bias the first element by -2.
\end{itemize}

\begin{table}[ht]
\begin{center}
\begin{tabular}{||l||c|c|c|c||}
\hline \hline
\textbf{operation}& \multicolumn{4}{c||}{Data Values} \\
  & 0 & 1 & 2 & 3  \\
\hline \hline
\textbf{initial}           &2&4&6&8 \\ \hline
\textbf{+1 to indices 0-2} &3&5&7&8 \\ \hline
\textbf{-2 to index 0}     &1&5&7&8 \\
\hline \hline
\end{tabular}
\end{center}
\end{table}

The biasing functions as expected.

\subsection {Scaling the Values of the Independent Variable}
The set of precalibrated data points used when defining the independent
variable can be scaled across some subset or all of the domain.
See section \ref{sec:UG_bias_scale} for a discussion of these processes.

The application of a scaling factor is investigated in
\textit{SIM\_05\_independent\_wraparound/RUN\_bias}.  This run also tests
certain error cases when applying a bias or scaling to the calibrated
values.

In this run, the central 3 data points of calibrated data set [2, 4, 6, 8,
9] are scaled by 1.05:
\begin{table}[!ht]
\begin{center}
\begin{tabular}{||l||l|l|l|l|l||}
\hline \hline
\textbf{operation}&
\multicolumn{5}{c||}{Data Values} \\
&0&1&2&3&5\\
\hline \hline
\textbf{initial}              &2&4&6&8&9 \\
\textbf{*1.05 to indices 1-3} &2&4.2&6.3&8.4&9 \\
\hline \hline
\end{tabular}
\end{center}
\end{table}

The scaling functions as expected.



\section {Applying Bias to Table Data} \label{sec:VV_bias_dependent}
Just as the independent variable data can be biased, so can the dependent
data.
See section \ref{sec:UG_bias_scale} for a discussion of these processes.
This is investigated in
\textit{SIM\_01a\_table\_lookup\_set\_nominal-/RUN\_biases}.  The bias is
applied as instructed to data at index 14.

This table has been populated as
\begin{code}
double temp_table1[3][4][3] =
        {{{3000,3001,3002},{3010,3011,3012},{3020,3021,3022},
          {3030,3031,3032}},
          {{3100,3101,3102},{3110,3111,3112},{3120,3121,3122},
          {3130,3131,3132}},
          {{3200,3201,3202},{3210,3211,3212},{3220,3221,3222},
          {3230,3231,3232}}};
\end{code}
(see \textit{SIM\_01a\_table\_lookup\_set\_nominal/include/test\_model.hh:83})

and can be represented visually as:
\begin{table}[H]
\begin{center}
\begin{tabular}{|l|l|l|}
\hline
3000 & 3001 & 3002\\3010 & 3011 & 3012\\3020 & 3021 & 3022\\3030 & 3031 & 3032\\
\hline
3100 & 3101 & 3102\\3110 & 3111 & 3112\\3120 & 3121 & 3122\\3130 & 3131 & 3132\\
\hline
3200 & 3201 & 3202\\3210 & 3211 & 3212\\3220 & 3221 & 3222\\3230 & 3231 & 3232\\
\hline
\end{tabular}
\end{center}
\end{table}

The data element at index 14 has the value $3102$; this is incremented
by $1000$ to $4102$ as expected.


\section{ Simple Lookup / Linear Interpolation with One Independent Variable}
\label{sec:VV_Simple_Lookup}
For this section, we consider \textit{SIM\_02\_simple\_table\_lookup},
specifically the outputs of \textit{RUN\_verif}.

This simulation is created with several instances of the
\textit{SimpleTableLookup} class, each with a different purpose:
\begin{itemize}
  \item \textbf{table\_single\_array} is created with 5 tabulated values for the
        independent variable and 5 for the dependent variable, \textit
        {\{2,4,6,8,10\}} and \textit{\{0,1,2,3,4\}} respectively.  The
        values for both variables are populated from an array.
  \item \textbf{table\_single\_vec} is identical to
        \textit{table\_single\_array} except
        that the values for both the independent and dependent variables are
        provided with a STL-vector rather than a classic array.
        For the purposes of data verification, these two runs are equivalent.
  \item \textbf{table\_interp} is similar to \textit{table\_single\_array}; it
        uses the same independent variable and tabulated values, but has 2
        dependent variables, and 2 sets of dependent values, one for each
        variable: \textit{\{\{0,1,2,3,4\},\{10,11,12,13,14\}\}}.
        The dependent data is loaded from an array.
        The table is configured to use interpolation to generate
        the value of the independent variable.
  \item \textbf{table\_prev} is similar to \textit{table\_interp}, except
        the dependent data is loaded from a 10-element STL-vector, and the
        table is configured to output the dependent values corresponding to
        the value at the index previous to the identified location in the
        independent table.
  \item \textbf{table\_next} is similar to \textit{table\_prev}, except
        the table is configured to output the dependent values corresponding to
        the value at the index following the identified location in the
        independent table.
  \item \textbf{table\_floor} is similar to \textit{table\_prev}, except
        the table is configured to output the dependent values corresponding to
        the value at the index identified as the index of the tabulated
        independent values whose value is
        immediately smaller than the current value of the independent variable.
        Because the independent variable is monotonically increasing, this
        is the same as the previous value; the output of
        \textit{table\_floor} and \textit{table\_prev} should match.
  \item \textbf{table\_ceil} is similar to \textit{table\_next}, except
        the table is configured to output the dependent values corresponding to
        the value at the index identified as the index of the tabulated
        independent values whose value is
        immediately larger than the value of the independent variable.
        Because the independent variable is monotonically increasing, this
        is the same as the ``next'' value; the output of
        \textit{table\_ceil} and \textit{table\_next} should match.
  \item \textbf{table\_round} is similar to \textit{table\_prev}, except
        the table is configured to output the dependent values corresponding to
        the value at the index identified as the index of the tabulated
        independent values whose value is closest to
        the value of the independent variable.
  \item \textbf{table\_wrap} is loaded similarly to \textit{table\_prev} but
        uses interpolation, more like \textit{table\_interp}.  For each
        of the other tables, when the independent variable takes a value
        outside the domain of its tabulated value, the dependent variable
        will take the tabulated value at the appropriate end of its set of
        values.  This table uses the wrap-around configuration, providing
        the independent variable with an infinite domain.
  \item \textbf{table\_reverse} This table works the same as
        \textit{table\_interp}, except the independent data is reversed to
        test the ability to interpolate between monotonically decreasing
        elements.  The independent variable now uses \textit{\{10,8,6,4,2\}}
        with the dependent values unchanged.
  \item \textbf{table\_error} is used for error checking; it is not used in
        \textit{RUN\_verif}
  \item \textbf{table\_1\_indep\_val} is used for error checking; it is not
        used in \textit{RUN\_verif}
\end{itemize}

The independent variable value is incremented from 0 to 12.5 in steps of 0.5,
then decremented back to 0 with steps of -0.5.

Consider a few examples:

\begin{itemize}
  \item independent value = 1.0.  This is outside the domain (off-table-left
  for most tables) of
  the independent variable.
  \begin{itemize}
    \item The single, interp, prev, next, floor, ceil and round
      configurations should all return the first element of their dependent
      values, i.e. 0 or 10.
    \item The wrap-around configuration should recognize
      that this is 1 less than the smallest value, which is therefore equivalent
      to 1 less
      than the largest value; this table should return the same output as it
      would if the independent variable had a value of 9 -- i.e. 3.5 and 13.5.
    \item For the reverse table, the input value is off-table-right,
      so the output values should be taken from the right-end of the
      tabulated values for the dependent variables, i.e. 4 and 14.
  \end{itemize}
  \item independent value = 5.0.  This is inside the domain of all tables,
  half-way between 4 and 6.  For all but the reverse table, the input value
  is halfway between the 2nd and 3rd tabulated values \{\{1,2\},\{11,12\}\}.
  \begin{itemize}
    \item The single tables should return the single interpolated value from
    the first set of data, 1.5.
    \item The interp table should return a pair of interpolated values, one
    from each set of data,
    1.5 and 11.5.
    \item The prev table should return the pair of values taken from the 2nd
    tabulated values of each set of data, 1 and 11.
    \item The next table should return the pair of values taken from the 3rd
    tabulated values, 2 and 12.
    \item The floor table should return the pair of values taken from the
    index associated with the next lower tabulated value -- i.e. 4.0, the 2nd
    tabulated value -- and return 1 and 11.
    \item The ceil table should return the pair of values taken from the index
    associated with the next higher tabulated value -- i.e. 6.0, the 3rd
    tabulated value -- and return 2 and 12.
    \item The round table should return the pair of values taken from the
    index
    associated with the closest tabulated value.  Depending on the exact
    numerical representation of 5.0, this could either round up or down,
    to the 2nd or 3rd tabulted value, resulting in either 1 and 11, or 2 and
    12.
    \item The wrap table will behave like the interp table because the
    independent value is within its ``natural'' domain.  It will return
    values 1.5 and 11.5.
    \item The reverse table has different indices, the value of 5.0 is
    halfway between the 3rd and 4th values.  It will return values 2.5
    and 12.5.
  \end{itemize}
\end {itemize}


The table shows the expected value of the first dependent variable.  The
value of the second dependent variable should always be 10 higher than the
first.  The output for the single table is also not shown in this table; its
single data output should match the first element of the interp table.

\clearpage




\begin{table}[ht]
\caption{Indices and Output Values for 8 Interpretations of the Input}
\begin{center}
\footnotesize  % small, footnotesize, scriptsize, tiny
\begin{tabular}{|l||l|l|c|c|c|c|c||l|l||l|l||}
\hline
\textbf{input} &
\multicolumn{7}{c||}{\textbf{Monotonic Incr}} &
\multicolumn{2}{c||}{\textbf{Wrap-around}} &
\multicolumn{2}{c||}{\textbf{Monotonic Decr}}\\
\hline
 & \textbf{index} &
\multicolumn{6}{c||}{\textbf{values}} &
\textbf{index} &
\textbf{value} &
\textbf{index} &
\textbf{value} \\
\hline
 & &
\textbf{interp} &
\textbf{prev} &
\textbf{next} &
\textbf{floor} &
\textbf{ceil} &
\textbf{round} & & & & \\
\hline \hline
0.5  & 0-   & 0    & 0 & 0 & 0 & 0 & 0       & 3.25 & 3.25     & 4+   & 4 \\
1.0  & 0-   & 0    & 0 & 0 & 0 & 0 & 0       & 3.5  & 3.5      & 4+   & 4 \\
1.5  & 0-   & 0    & 0 & 0 & 0 & 0 & 0       & 3.75 & 3.75     & 4+   & 4 \\
2.0  & 0    & 0    & 0 & 0 & 0 & 0 & 0       & 0    & 0 / 4 ** & 4    & 4 \\
2.5  & 0.25 & 0.25 & 0 & 1 & 0 & 1 & 0       & 0.25 & 0.25     & 3.75 & 3.75 \\
3.0  & 0.50 & 0.50 & 0 & 1 & 0 & 1 & 0 / 1 * & 0.50 & 0.50     & 3.50 & 3.50 \\
3.5  & 0.75 & 0.75 & 0 & 1 & 0 & 1 & 1       & 0.75 & 0.75     & 3.25 & 3.25 \\
4.0  & 1.00 & 1.00 & 1 & 1 & 1 & 1 & 1       & 1.00 & 1.00     & 3.00 & 3.00 \\
4.5  & 1.25 & 1.25 & 1 & 2 & 1 & 2 & 1       & 1.25 & 1.25     & 2.75 & 2.75 \\
5.0  & 1.50 & 1.50 & 1 & 2 & 1 & 2 & 1 / 2 * & 1.50 & 1.50     & 2.50 & 2.50 \\
5.5  & 1.75 & 1.75 & 1 & 2 & 1 & 2 & 2       & 1.75 & 1.75     & 2.25 & 2.25 \\
6.0  & 2.00 & 2.00 & 2 & 2 & 2 & 2 & 2       & 2.00 & 2.00     & 2.00 & 2.00 \\
6.5  & 2.25 & 2.25 & 2 & 3 & 2 & 3 & 2       & 2.25 & 2.25     & 1.75 & 1.75 \\
7.0  & 2.50 & 2.50 & 2 & 3 & 2 & 3 & 2 / 3 * & 2.50 & 2.50     & 1.50 & 1.50 \\
7.5  & 2.75 & 2.75 & 2 & 3 & 2 & 3 & 3       & 2.75 & 2.75     & 1.25 & 1.25 \\
8.0  & 3.00 & 3.00 & 3 & 3 & 3 & 3 & 3       & 3.00 & 3.00     & 1.00 & 1.00 \\
8.5  & 3.25 & 3.25 & 3 & 4 & 3 & 4 & 3       & 3.25 & 3.25     & 0.75 & 0.75 \\
9.0  & 3.50 & 3.50 & 3 & 4 & 3 & 4 & 3 / 4 * & 3.50 & 3.50     & 0.50 & 0.50 \\
9.5  & 3.75 & 3.75 & 3 & 4 & 3 & 4 & 4       & 3.75 & 3.75     & 0.25 & 0.25 \\
10.0 & 4.00 & 4.00 & 4 & 4 & 4 & 4 & 4       & 0.00 & 0 / 4 ** & 0.00 & 0.00 \\
10.5 & 4+   & 4.00 & 4 & 4 & 4 & 4 & 4       & 0.25 & 0.25     & 0-   & 0    \\
11.0 & 4+   & 4.00 & 4 & 4 & 4 & 4 & 4       & 0.50 & 0.50     & 0-   & 0    \\
11.5 & 4+   & 4.00 & 4 & 4 & 4 & 4 & 4       & 0.75 & 0.75     & 0-   & 0    \\
\hline
\end{tabular}
\medskip
\end{center}
\end{table}


* These values are a result of rounding from the midpoint.  Depending on the
numerical representation, they could round down or up.

** At the wrap-around, the outcome could go to either end of the output
array, resulting in 0.0 or 4.0.  Typically, a value that hits the maximum
value precisely will be wrapped back to the minimum value.

The recorded data includes the values for the \textit{single\_array} and
\textit{single\_vec} tables that have been omitted from this table.  The
recorded data do not include the working values of the indices for the
tables which were included in the table for clarification.

The recorded data values record the variables identified as
text.dependent[a][b].  The following table identifies the indices to which each
of the lookup tables writes their output data:

\begin{table}[ht]
\caption{Indices of the test.dependent array attributed to each table.}
\begin{center}
\begin{tabular}{|l|c|c|}
\hline
\textbf{table} &
\textbf{first variable} &
\textbf{second variable} \\
\hline


single\_array & [0][0]  & -      \\
single\_vec   & [1][0]  & -      \\
interp        & [2][0]  & [2][1] \\
prev          & [3][0]  & [3][1] \\
next          & [4][0]  & [4][1] \\
floor         & [5][0]  & [5][1] \\
ceil          & [6][0]  & [6][1] \\
round         & [7][0]  & [7][1] \\
wrap          & [8][0]  & [8][1] \\
reverse       & [9][0]  & [9][1] \\
\hline
\end{tabular}
\end{center}
\end{table}

Upon inspection, all data outputs match expectations.

The results are committed to

\textit{SIM\_02\_simple\_table\_lookup/verif\_data/RUN\_verif/log\_test\_data.csv}



\section{ Linear Interpolation with Multiple Independent Variables}
This capability is tested by \textit{RUN\_interp\_all}, found in
\textit{SIM\_01a\_table\_lookup\_set\_nominal}. See the next section for a
description of this simulation and its results.

\section{ Mixed Lookup and Linear Interpolation with Multiple Independent
Variables}

There are multiple ways in which a set of independent variables may be
combined.  The simplest implementation can be generated by applying a single
rule to all variables, such as a \textit{n-} dimensional interpolation.
However, for complete generalization, we also need to test the
ability to apply different interpretaions to each of the independent
variables, using some for continuous interpolation while
using others for any combination of the available discrete data lookups.

The ``interpolate using all variables''
case is a subset of this more general analysis, so is included in this
simulation.  For this section, we use the
\textit{SIM\_01a\_table\_lookup\_set\_nominal} simulation.

This test is obviously limited to the available implementations for
converting the value of the input variable into a table index.  However,
successful demonstration of the ability to use multiple implementations
provides support for the concept that any algorithm could be implemented and
integrated into the table architecture.  This would include, for example,
non-linear interpolation algorithms.


This simulation uses the full \textit{TableLookupSet} as the table manager
and the \textit{GenericMultiInputTable} for the actual table
lookup/interpolation.  The simulation includes two table managers; the
distinction between them is in the way the tables are created:
\begin{itemize}
  \item \textbf{test\_so\_defined} tests the results from a table-manager
  whose managed contents are defined within the simulation object itself and
  pre-compiled.
  \item \textbf{test\_on\_the\_fly} tests the same numerical capabilities
  within an environment in which the managed contents are created at
  run-time inside another model.
\end{itemize}
The results from both managers should match.

Both managers manage two data tables; combined, they populate three dependent
variables
as functions of four independent variables. The first three independent
variables use the default Linear configuration (finite domain), the fourth
has the Wrap-around capability (infinite domain).

\begin{itemize}
  \item \textbf{multi\_table} uses the first three independent variables to
  generate values for the first two dependent variables
  \item \textbf{single\_table} uses the first two and the fourth independent
  variables to generate values for the third dependent variable.
\end{itemize}

There are multiple tests, examining the different interpretation configurations
(between ceiling, floor, interpolate, next, previous, and round):
of the four independent variables
\begin{itemize}
  \item RUN\_ceil\_int\_int\_int
  \item RUN\_floor\_int\_int\_int
  \item RUN\_int\_ceil\_int\_int
  \item RUN\_int\_floor\_int\_int
  \item RUN\_int\_next\_int\_int
  \item RUN\_int\_prev\_int\_int
  \item RUN\_int\_round\_int\_int
  \item RUN\_interp\_all
  \item RUN\_next\_int\_int\_int
  \item RUN\_prev\_int\_int\_int
  \item RUN\_prev\_next\_floor\_ceil
\end{itemize}

The python script \textit{verif.py} reads in the resulting data files,
extracts the input values and recomputes the expected outputs.  The computed
outputs are then compared against the logged outputs from the simulation to
identify any differences. The script also compares the outputs of the two
table managers to verify that the output is not a function of how the tables
are initially allocated. No differences are found in either test.


Two sample points are presented here to illustrate the difference between the
interpretations of the independent variable.  Both sample points are taken
from the third dependent variable.
\begin{itemize}
  \item The first independent variable has 3 calibration points: $\{9,8,-2\}$
  \item The second independent variable has 4 calibration points: $\{2,4,6,8\}$
  \item The third independent variable has 3 calibration points: $\{2,4,6\}$ but
  wraps around indefinitely
  \item the data for the third dependent variable is drawn from the
  3-dimensional array constructed as
  $3000 + 100* index_1 + 10 * index_2 + index_3$:

  \begin{table}[ht]
  \begin{center}
  \begin{tabular}{|l|l|l|}
  \hline
  3000 & 3001 & 3002\\3010 & 3011 & 3012\\3020 & 3021 & 3022\\3030 & 3031 & 3032\\
  \hline
  3100 & 3101 & 3102\\3110 & 3111 & 3112\\3120 & 3121 & 3122\\3130 & 3131 & 3132\\
  \hline
  3200 & 3201 & 3202\\3210 & 3211 & 3212\\3220 & 3221 & 3222\\3230 & 3231 & 3232\\
  \hline
  \end{tabular}
  \end{center}
  \end{table}
\end{itemize}

Note -- because these data points lie on a plane (a 3-dimensional plane in
4-dimensional space), the linear interpolation between points can be
factored and simplified.  For interpolation in 3-dimensions, the model
algorithm weights the 8 points at the vertices of the cube containing the
point of interest.  Suppose the point of interest, $A$ is at a fractional
distance $(p,q,r)$  on each respective dimension from the base vertex
$S_{ijk}$.  Then the interpolated point is computed within the model as:
\begin{equation*}
A = S_{ijk}(1-p)(1-q)(1-r) + S_{ij(k+1)}(1-p)(-q)r + ...
\end{equation*}
as described in equation \ref{eq:3d_interp}.

Because the points are co-planar in this test case, we recognize that
\begin{equation*}
S_{(i+1)jk} - S_{ijk} =
S_{(i+1)(j+1)k} - S_{i(j+1)k} =
S_{(i+1)j(k+1)} - S_{ij(k+1)} =
S_{(i+1)(j+1)(k+1)} - S_{i(j+1)(k+1)} =
\Delta x
\end{equation*}
and similar for the y- and z-axes.  With this simplification, the evaluation
of the interpolated value simplifies to:
\begin{equation*}
A = S_{ijk} + p \Delta x + q \Delta y + r \Delta z
\end{equation*}
Because the value $S_{ijk}$ changes linearly, it can be similarly expressed as:
\begin{equation*}
S_{ijk} = S_{000} + i \Delta x + j \Delta y + k \Delta z
\end{equation*}
Thus:
\begin{equation*}
A = S_{000} + (i+p) \Delta x + (j+q) \Delta y + (k+r) \Delta z
\end{equation*}

Specifically in this example, this evaluates to
\begin{equation*}
A = 3000 + 100 (i+p) + 10 (j+q) + (k+r)
\end{equation*}

The python script \textit{verif.py} uses this specific algorithm to verify
the values generated by the vertex-weighting method used in the model.


Breakdowns are illustrated for input values of $\{5.6, 5.6, 5.0\}$ and
$\{7.6, 7.6, 6.25\}$ (entries 13 and 18 respectively in the log files)
in the following tables.  These tables also use the simplified algorithm to
compute the predicted values; these values match those generated by the
model's vertex-weighting algorithm.

Note -- the letters \{c,f,i,n,p,r\} set as the run ID in the tables map to
the independent variable behaviors
\{ceiling, floor, interpolate, next, previous, round\} respectively.

Manual spot-checks have also been applied to interpolations between non-planar
points.

\begin{table}[ht]
\caption{Results for input \{5.6, 5.6, 5.0\}}
\begin{center}
\begin{tabular}{|l||l|l|l||llll|l|}
\hline
run ID & index1 & index2 & index 3 & \multicolumn{5}{c|}{output} \\
       &        &        &         & base &  +  & +  &   + &   =      \\ \hline \hline
cii    & 1      & 1.8    & 1.5     & 3000 & 100 & 18 & 1.5 & 3119.5 \\ \hline
fii    & 2      & 1.8    & 1.5     & 3000 & 200 & 18 & 1.5 & 3219.5 \\ \hline
ici    & 1.24   & 2      & 1.5     & 3000 & 124 & 20 & 1.5 & 3145.5 \\ \hline
iii    & 1.24   & 1.8    & 1.5     & 3000 & 124 & 18 & 1.5 & 3143.5 \\ \hline
ifi    & 1.24   & 1      & 1.5     & 3000 & 124 & 10 & 1.5 & 3135.5 \\ \hline
ini    & 1.24   & 2      & 1.5     & 3000 & 124 & 20 & 1.5 & 3145.5 \\ \hline
ipi    & 1.24   & 1      & 1.5     & 3000 & 124 & 10 & 1.5 & 3135.5 \\ \hline
iri    & 1.24   & 2      & 1.5     & 3000 & 124 & 18 & 1.5 & 3145.5 \\ \hline
nii    & 2      & 1.8    & 1.5     & 3000 & 200 & 18 & 1.5 & 3219.5 \\ \hline
pii    & 1      & 1.8    & 1.5     & 3000 & 100 & 18 & 1.5 & 3119.5 \\ \hline
pnc    & 1      & 2      & 2       & 3000 & 100 & 20 & 2   & 3122   \\ \hline
\hline
\end{tabular}
\end{center}
\end{table}

\begin{table}[H]
\caption{Results for input \{7.6, 7.6, 6.25\}}
\begin{center}
\begin{tabular}{|l||l|l|l||llll|l|}
\hline
run ID & index1 & index2 & index 3 & \multicolumn{5}{c|}{output} \\
       &        &        &         & base &  +  & +  &   +   &   =      \\ \hline \hline
cii    & 1      & 2.8    & 0.125   & 3000 & 100 & 28 & 0.125 & 3128.125 \\ \hline
fii    & 2      & 2.8    & 0.125   & 3000 & 200 & 28 & 0.125 & 3228.125 \\ \hline
ici    & 1.04   & 3      & 0.125   & 3000 & 104 & 30 & 0.125 & 3134.125 \\ \hline
iii    & 1.04   & 2.8    & 0.125   & 3000 & 104 & 28 & 0.125 & 3132.125 \\ \hline
ifi    & 1.04   & 2      & 0.125   & 3000 & 104 & 20 & 0.125 & 3124.125 \\ \hline
ini    & 1.04   & 3      & 0.125   & 3000 & 104 & 30 & 0.125 & 3134.125 \\ \hline
ipi    & 1.04   & 2      & 0.125   & 3000 & 104 & 20 & 0.125 & 3124.125 \\ \hline
iri    & 1.04   & 3      & 0.125   & 3000 & 104 & 28 & 0.125 & 3134.125 \\ \hline
nii    & 2      & 2.8    & 0.125   & 3000 & 200 & 28 & 0.125 & 3228.125 \\ \hline
pii    & 1      & 2.8    & 0.125   & 3000 & 100 & 28 & 0.125 & 3128.125 \\ \hline
pnc    & 1      & 3      & 1       & 3000 & 100 & 30 & 1     & 3131     \\ \hline
\hline
\end{tabular}
\end{center}
\end{table}

%\clearpage

\section {Interpolation Between Angles}
Because angles are often intended to be interpreted as cyclically repetitive,
they can cause problems for simple interpolation algorithms.
\begin{itemize}
\item In some cases,
an angular value is intended to be unbounded, such as computing the total
angular displacement as a function of time. In cases where 370
degrees has a meaning distinct from that of 10 degrees, the treatment of this
type of
angular representation is no different than that of a regular non-cyclical
variable. This has already been verified and presented.
\item In cases where the dependent variable is a function of an angular (or
other cyclical) variable, it is necessary to wrap the independent variable
so that when it has a value of 370 degrees, the model returns the same
output as when the independent variable had a value of 10 degrees.  This
ability to wrap the independent variable has already been examined and
verified.
\item Cases where the dependent variable is an angular variable and a
function of one independent variable are examined in this section.  The
complication with this situation is the ambiguity associated with the
wrap-around of the dependent variable.  For example, the value obtained when
interpolating between 350 degrees and 10 degrees could be somewhere close to
0 degrees, or somewhere close to 180 degrees.  This model assumes the shortest
path, which should result in a value close to 0 degrees in this example.
\item Cases where the dependent variable is an angular variable and a
function of multiple variables introduces multiple levels of ambiguity.
This type of interpolation is usually very specific and not covered by this
general model.
\end{itemize}

\textbf{Important note on range} The output from this type of
lookup/interpolation will
always be in the range $(-180,180]$ degrees, or $(-\pi, \pi]$ radians.

For this section, we will be working with
\textit{SIM\_03\_simple\_table\_lookup\_for\_angles} which has only 1 run,
\textit{RUN\_verif}.


The simulation provides:
\begin{itemize}
  \item 2 table managers of type \textit{SimpleTableLookup}, one for degrees
  and one for radians.
  \item 4 tables of type \textit{SimpleInputTableForAngles}.  These instances
  are present just to test the different mechanisms by which a table might be
  constructed, they are not used for generating output in this simulation.
\end{itemize}

As with any instances of \textit{SimpleTableLookup}, the 2 multi-purpose
table-managers create new instances of specific tables according to the
table-types they are given when the dependent data is loaded. Previously, we
have seen these use only the default \textit{Generic} table-type, but in this
case, one is given the specification to construct an \textit{AngularDeg}
table-type, and the other an \textit{AngularRad} table-type.  Note that both
of these result in the creation of a \textit{SimpleInputTableForAngles}
instance for use by the \textit{SimpleTableLookup} manager; the specification
passed to the manager is forwarded to the \textit{SimpleInputTableForAngles}
instance, resulting in the setting of an internal flag that affects the
periodicity of
the wrapping -- for a Radians-based table, the dependent values are wrapped
at $ 2 \pi$,
while for the Degrees-based table the wrap is effected at $360$.

Both tables are supplied with the same independent variable, calibrated with
values of
${0,2,4,6,8,10}$ and modified in the test to increment from $0.5$ to $12.5$ in
steps of
$0.5$ and then to decrement to $-0.5$ in steps of $-0.5$

Both tables generate values for two dependent variables -- both in degrees, or
both in
radians according to the spcification of the table.  These dependent variables
are
calibrated with values of:

\begin{table}[ht]
\caption{Degree Table}
\begin{center}
\begin{tabular}{|c|c|c|c|c|c|}
\hline
30&90&150&-150&-90&-30\\
\hline
30&140&250&340&430&560\\
\hline
\end{tabular}
\end{center}
\end{table}

and
\begin{table}[ht]
\caption{Radians Table}
\begin{center}
\begin{tabular}{|c|c|c|c|c|c|}
\hline
0.5& 1.0& 1.5&-1.5&-1.0& -0.5\\
\hline
1.5& 3.0& 4.5& 6.0& 7.5& 10.0\\
\hline
\end{tabular}
\end{center}
\end{table}

Consider a couple of examples:
\begin{itemize}
  \item When the independent variable has a value of 5, the index is 2.5,
  and the
  outputs should be midway between:
  \begin{itemize}
    \item $[150,-150]$.  The simple difference between these values is $300$,
    which is more than the half-circle $180$.  Hence, the shortest path
    between these values is
    \textbf{not} the one passing through $0$.  $-150$ can be expressed as its
    equivalent
    $210$; now the difference is less than the half-circle $180$ and the
    values can be
    interpolated to yield $180$.  This is in the default range and is the
    expected output.
    \item $[250,340]$.  The simple difference between these values is $90$,
    which is less
    than the half-circle $180$.  The result is the midpoint of these two
    values, which
    computes to $295$.  However, $295$ degrees is outside the default range,
    so this value is further modified to yield the in-range equivalent value
    of $-65$ degrees.
    \item $[1.5,-1.5]$.  The simple difference between these values is $3$,
    which is less
    than the half-circle $\pi$, so the shortest distance between these values
    is the one
    passing through $0$, the result is $0$ radians.
    \item $[4.5, 6.0]$.   The simple difference between these values is $1.5$,
    which is less
    than the half-circle $\pi$.  The result is the midpoint of these two
    values, which
    computes to $5.25$.  However, $5.25$ radians is outside the default
    range, so this value
    is further modified to yield the in-range equivalent value of $-1.033$
    radians.
  \end{itemize}
  \item When the independent variable has a value of $11$, the independent
  variable is
  off-table-right, and the outputs should be the last calibrated values:
  \begin{itemize}
    \item $-30$.  This value is within the default range.
    \item $560$.  This value is outside the default range and will be
    modified to its
    equivalent $-160$ degrees.
    \item $-0.5$.  This value is within the default range.
    \item $10.0$.  This value is outside the default range and will be
    modified to its
    equivalent $-2.56$ radians.
  \end{itemize}
\end{itemize}


A partially abbreviated results table is presented below:
\begin{longtable}{|c|c|c|c|c|}
\caption{Results Table for Angular Interpolation}\\
\hline
\textbf{independent} & \textbf{degrees 1} & \textbf{degrees 2} & \
\textbf{radians 1} & \textbf{radians 2}\\
\hline
\endfirsthead
\caption{\textit{(Continued)} Results Table for Angular Interpolation}\\
\hline
\textbf{independent} & \textbf{degrees 1} & \textbf{degrees 2} & \
\textbf{radians 1} & \textbf{radians 2}\\
\hline
\endhead
\hline
\multicolumn{5}{l}{\textit{Continued on next page}} \\
\endfoot
\hline
\endlastfoot
 0.5&      45&    57.5&   0.625&   1.875 \\ \hline
   1&      60&      85&    0.75&    2.25 \\ \hline
 1.5&      75&   112.5&   0.875&   2.625 \\ \hline
   2&      90&     140&       1&       3 \\ \hline
 2.5&     105&   167.5&   1.125&  -2.908185\\ \hline
   3&     120&    -165&    1.25&  -2.533185\\ \hline
 3.5&     135&  -137.5&   1.375&  -2.158185\\ \hline
   4&     150&    -110&     1.5&  -1.783185\\ \hline
 4.5&     165&   -87.5&    0.75&  -1.408185\\ \hline
   5&     180&     -65&       0&  -1.033185\\ \hline
 5.5&    -165&   -42.5&   -0.75&  -0.658185 \\ \hline
   6&    -150&     -20&    -1.5&  -0.283185 \\ \hline
 6.5&    -135&     2.5&  -1.375&   0.091814\\ \hline
   7&    -120&      25&   -1.25&   0.466814 \\ \hline
 7.5&    -105&    47.5&  -1.125&   0.841814 \\ \hline
   8&     -90&      70&      -1&   1.216814\\ \hline
 8.5&     -75&   102.5&  -0.875&   1.841814\\ \hline
   9&     -60&     135&   -0.75&   2.466814\\ \hline
 9.5&     -45&   167.5&  -0.625&   3.091814\\ \hline
  10&     -30&    -160&    -0.5&  -2.566370\\ \hline
10.5&     -30&    -160&    -0.5&  -2.566370\\ \hline
  10&     -30&    -160&    -0.5&  -2.566370\\ \hline
 9.5&     -45&   167.5&  -0.625&   3.091814\\ \hline
   9&     -60&     135&   -0.75&   2.466814\\ \hline
 8.5&     -75&   102.5&  -0.875&   1.841814\\ \hline
   8&     -90&      70&      -1&   1.216814\\ \hline
 7.5&    -105&    47.5&  -1.125&   0.841814 \\ \hline
   7&    -120&      25&   -1.25&   0.466814 \\ \hline
 6.5&    -135&     2.5&  -1.375&   0.091814 \\ \hline
   6&    -150&     -20&    -1.5&  -0.283185 \\ \hline
 5.5&    -165&   -42.5&   -0.75&  -0.658185 \\ \hline
   5&     180&     -65&       0&  -1.033185\\ \hline
 4.5&     165&   -87.5&    0.75&  -1.408185\\ \hline
   4&     150&    -110&     1.5&  -1.783185\\ \hline
 3.5&     135&  -137.5&   1.375&  -2.158185\\ \hline
   3&     120&    -165&    1.25&  -2.533185\\ \hline
 2.5&     105&   167.5&   1.125&  -2.908185\\ \hline
   2&      90&     140&       1&       3 \\ \hline
 1.5&      75&   112.5&   0.875&   2.625 \\ \hline
   1&      60&      85&    0.75&    2.25 \\ \hline
 0.5&      45&    57.5&   0.625&   1.875 \\ \hline
   0&      30&      30&     0.5&     1.5 \\ \hline
-0.5&      30&      30&     0.5&     1.5 \\ \hline

\end{longtable}

\clearpage

\section {Interpolation Between Quaternions}
Interpolation between two calibrated quaternion values is provided by the
public domain Spherical Linear Interpolation algorithm (\textit{slerp}).
The details
and formulation of this alogrithm are presented in section \ref{sec:Math_slerp}.

Verification of this algorithm is provided by
\textit{SIM\_04a\_table\_lookup\_for\_quaternions}, especially run
\textit{RUN\_verif}.

This simple test uses a 2-element independent variable array ($\{0,1\}$), and
two quaternions \linebreak
($\{[1, 0, 0, 0] , [0.5, 0.866, 0, 0]\}$) for calibrated data points.
The two quaternions represent, respectively, the identity quaternion,
and a rotation of 120 degrees about the x-axis.

The different ways of configuring this option are investigated, including:
\begin{enumerate}
  \item Using a \textit{SimpleTableLookup} manager and the
  \textit{AbstractTableLookup::Quaternion} option when loading the
  dependent data.
  \item using a \textit{TableLookupSet} manager and individual instances of
  the \textit{SingleInputTableForQuaternions} tables:
  \begin{enumerate}
    \item Output is a double 4-array, using the default constructor.
    The elements of the
    output array are manually added to this table using
    \textit{add\_dependent}.
    \item Output is a double 4-array, using the array constructor.
    The elements of the output array are automatically added to this table
    by the constructor.
    \item Output is a JEOD Quaternion type; table is constructed using the
    Quaternion-based constructor.
    \item Output is a double 4-array, and the slerp algorithm is disabled.
    In this case, the interpolation will be nearly linear -- the model will
    use linear interpolation to generate values for the quaternion, but
    these values will then be normalized, moving
    the interpolated values off the straight line.
  \end{enumerate}
\end{enumerate}

The output from tables $1$ and $2(a-c)$ should match
numerically and provide a quaternion that respresents the rotation about the
x-axis by the interpolated angle.  The output from table $2.4$ should be
different, not necessarily in an easily predictable way.

We also test, in \textit{RUN\_verif\_with\_negatives}, the effect of
negating the second calibration quaternion:
${[1, 0, 0, 0] , [-0.5, -0.866, 0, 0]}$

Recognize that completely negating a quaternion does not change its meaning
(in this case, one is a rotation of 120 degrees and the other a rotation of
-240 degrees) and that consequently, the slerp-based outputs should be
identical, but the linear-interpolation outputs will differ.

For all tests, the value of the independent variable is incremented from $0.2$
to $1$ and back to $-0.2$ in steps of $\pm0.2$.

In the results table, we only show 2 elements of each quaternion: the scalar
component and first element of the vector component.  The other 2 elements
of the vector component are zero for all tests.

\begin{table}
\caption{Results Table for Quaternion Interpolation}
\begin{center}
\begin{tabular}{||l||r||l|l||l|l||l|l|l||}
\hline \hline

indep& angle &
\multicolumn{2}{c||}{Analytic} &
\multicolumn{2}{c||}{SLERP} &
\multicolumn{3}{c||}{Linear} \\

value  & / deg &
scalar & vec[0] &
scalar & vec[0] &
scalar & vec[0] & angle \\


\hline \hline
 0.2  &  24 &
 0.978148 & 0.207912 &
 0.978148 & 0.207909 &
 0.981982 & 0.188977 & 21.8 \\ \hline

 0.4  &  48 &
 0.913545 & 0.406737 &
 0.913548 & 0.406732 &
 0.917667 & 0.397350 & 46.8 \\ \hline

 0.6  &  72 &
 0.809017 & 0.587785 &
 0.809021 & 0.587779 &
 0.802963 & 0.596028 & 73.2 \\ \hline

 0.8  & 106 &
 0.669131 & 0.743145 &
 0.669138 & 0.743138 &
 0.654665 & 0.755919 & 93.2 \\ \hline

 1    & 120 &
 0.5                & 0.866025  &
 0.5                & 0.86      &
 -0.5               & -0.866    &120   \\ \hline

 0.8  & 106 &
 0.669131  & 0.743145  &
 0.669138  & 0.743138  &
 0.654665  & 0.755919  & 93.2 \\ \hline

 0.6  &  72 &
 0.809017 & 0.587785 &
 0.809021 & 0.587779 &
 0.802963 & 0.596028 & 73.2 \\ \hline

 0.4  &  48 &
 0.913545 & 0.406737 &
 0.913548 & 0.406732 &
 0.917667 & 0.397350 & 46.8 \\ \hline

 0.2  &  24 &
 0.978148 & 0.207912 &
 0.978148 & 0.207909 &
 0.981982 & 0.188977 & 21.8 \\ \hline

 0    &   0 &
 1        &  0            &
 1        &  4.807e-17 &
 1        & -4.807e-17 &  0   \\ \hline

-0.2  &   0 &
 1        & 0       &
 1        & 0       &
 1        & 0       &  0   \\ \hline

\end{tabular}
\end{center}
\end{table}

Use of the slerp algorithm agrees very closely with the predicted
quaternions generated from the interpolated angle values.


\section{Verification of Interpolation of Coupled Variable+Derivative}
Review section \ref{sec:UG_var_deriv} for an understanding of what is meant
by coupled data. In summary, we are looking at simultaneously using values
for a variable \textit{and its derivative} in computing the interpolated
values of both.

This capability is examined in \textit{SIM\_04b\_table\_lookup\_var\_deriv}.
This verification test has 10 runs, but many of these produce intentionally
identical results, using different formats for specifying the same
configurations.

The tests are set up to investigate the interpolation within 5 different
functions:
\begin{itemize}
  \item A constant value, derivative is zero.
  \item a linear function, derivative is constant.
  \item a quadratic function, derivative is linear.
  \item a cubic function, derivative is quadratic.
  \item a sine function, derivative is a cosine function.
\end{itemize}
For each function, reference points are generated at the integer points, and
the interpolation is conducted at the quarter-integer points.

The following tests are established to investigate these functions.
\begin{itemize}
  \item Testing a set of 5 variables and their derivatives across the same 6
  reference points using a generic TableLookupSet is conducted in the following
  runs:
  \begin{itemize}
    \item \textbf{RUN\_01a\_default} constructs the table using the default
      constructor and independently adds both sets of dependent variables.
    \item \textbf{RUN\_01b\_1array}  constructs the table using the the C-style
      array specification for the primary variable; the derivative variable
      is added independently.
    \item \textbf{RUN\_01c\_2array} constructs the table using the the C-style
      array specification for both the primary and derivative variables.
    \item \textbf{RUN\_01d\_1vec} constructs the table using the the STL-vector
      specification for the primary variable; the derivative variable is added
      independently.
    \item \textbf{RUN\_01e\_2vec} constructs the table using the the STL-vector
      specification for both the primary and derivative variables.
  \end{itemize}

  \item To test the single-output constructors, we use only the input/output
    for the sine function, and the following runs:
  \begin{itemize}
    \item \textbf{RUN\_01f\_1elem} constructs the table using a
      single reference specification for the primary variable; the derivative
      variable is added independently.
    \item \textbf{RUN\_01g\_2elem} constructs the table using a
      single reference specification for both the primary and derivative
      variables.
  \end{itemize}

  \item To test the compatibility with the SimpleTableLookup instead of a
      TableLookupSet, we use the same configuration of 5 function tests
      in \textbf{RUN\_01h\_simple}.
  \item The optional \textit{omit\_derivative\_vals} flag is tested in
    \textbf{RUN\_02a\_default\_no\_xdot}. This run uses the same configuration
    as seen in \textit{RUN\_01a\_default}, but with the flag set the
    implementation interpolates only the primary variable (still using the
    coupled interpolation algorithm).
  \item The optional \textit{use\_linear\_interpolation} flag is tested in
    \textbf{RUN\_02b\_default\_linear}. This run uses the same configuration
    as seen in \textit{RUN\_01a\_default}, but with the flag set the
    implementation reverts back to using standard linear interpolation.
\end{itemize}

As expected, the results for runs 01a-01e, and 01h are identical. For runs 01f
and 01g, the outputs for the first four functions are all zeroes because these
functions are not interpolated for these runs. The outputs for the 5th function
match those from run 01a. The equivalence of the different constructors is
confirmed.

The results from run 02a are as expected: the output values for the priamry
variable match those from run 01a, and the outputs for the derivative variable
are all zeroes (not computed). The ability to compute interpolated values for
only the primary variable while still utilizing derivative values is confirmed.

The results from run 02b confirm independent linear interpolation on both the
primary and derivative values. The abliity to switch between independent-linear
interpolation and coupled-cubic interpolation is confirmed.

For comparison, the following plots illustrate the difference between the
results from run 01a and run 02b, i.e. the difference between using
independent-linear interpolation for the primary and derivative variables
versus using the coupled-cubic interpolation for both. These plots are
generated using the script at \textit{Octave/analysis.m}.

Figure \ref{fig:var_deriv_simple} provides the coupled-interpolation outputs
for the first two simple functions (constant and linear).

For Figures \ref{fig:var_deriv_quad}, \ref{fig:var_deriv_cubic}, and
\ref{fig:var_deriv_sine} we co-plot, for both the primary and derivative
variables, the:
\begin{itemize}
  \item output of the coupled-cubic interpolation
  \item output of the independent-linear interpolation
  \item the analytical values from the functions used to generate the reference
  points.
\end{itemize}

For the sine and cubic functions, comparison plots are provided showing the
deviation of the interpolation output from the analytical result in Figures
\ref{fig:var_deriv_cubic_delta} and \ref{fig:var_deriv_sine_delta}.
As might be expected, applying cubic-interpolation to a cubic function
results in zero error. Applying cubic-interpolation to a sinusoid function does
introduce error, but substantially smaller than that seen using linear
interpolation. The errors for both algorithms go to zero at the reference
points (integers).

In all cases, the test is allowed to run past the last reference point. In all
cases, the output beyond the last reference point is held at the value of that
point. This is expected, these algorithms are intended for, and deliberately
restricted to \textbf{interpolation} and do not support \textbf{extrapolation}.



\begin{figure}[!htb]
  \centering
  \includegraphics[scale=0.7]{Images/simple.png}
  \caption{Results of coupled-cubic interpolation for
           constant and linear functions}
        \label{fig:var_deriv_simple}
\end{figure}
\begin{figure}[!htb]
  \centering
  \includegraphics[scale=0.7]{Images/quadratic.png}
  \caption{Results of coupled-cubic and independent-linear interpolation
        applied to the primary and derivative values of a quadratic function.
        The analytically-obtained values are also plotted.}
        \label{fig:var_deriv_quad}
\end{figure}
\begin{figure}[!htb]
  \centering
  \includegraphics[scale=0.7]{Images/cubic.png}
  \caption{Results of coupled-cubic and independent-linear interpolation
        applied to the primary and derivative values of a cubic function.
        The analytically-obtained values are also plotted.}
        \label{fig:var_deriv_cubic}
\end{figure}
\begin{figure}[!htb]
  \centering
  \includegraphics[scale=0.7]{Images/sine.png}
  \caption{Results of coupled-cubic and independent-linear interpolation
        applied to the primary and derivative values of a sine function.
        The analytically-obtained values are also plotted.}
        \label{fig:var_deriv_sine}
\end{figure}
\begin{figure}[!htb]
  \centering
  \includegraphics[scale=0.7]{Images/cubic_delta.png}
  \caption{The difference between the analytical values and the
        interpolation values obtained using coupled-cubic and linear
        interpolation for a cubic function.
        The coupled-cubic interpolation tracks the analytical cubic function
        perfectly while the linear interpolation deviates between the reference
        points.}
        \label{fig:var_deriv_cubic_delta}
\end{figure}
\begin{figure}[!htb]
  \centering
  \includegraphics[scale=0.7]{Images/sine_delta.png}
  \caption{The difference between the analytical values and the
        interpolation values obtained using coupled-cubic and linear
        interpolation for a sine function.
        The coupled-cubic interpolation tracks the analytical sine function
        with significantly less error than is seen with linear interpolation.}
        \label{fig:var_deriv_sine_delta}
\end{figure}

\clearpage





\section{Using Data Sets Presented in Transpose Form}
Review section \ref{sec:UG_transpose} for an understanding of what is meant
by transpose data.

This capability is examined in \textit{SIM\_06\_transpose\_data\_manager},
with specific configuration options investigated in
\textit{SIM\_08\_transpose\_config}.

In \textit{SIM\_06\_transpose\_data\_manager}, an instance of
\textit{TableLookupTransposeDataSet} is created and
loaded with data from two separate mechanisms:
\begin{enumerate}
  \item \textit{RUN\_01\_internal\_data} examines data being compiled with
  the simulation and added to the manager
  \item \textit{RUN\_02\_external\_data} examines data being loaded onto an
  empty manager from a text file.
\end{enumerate}

In both cases, the data loaded has values shown in Table \ref{tab:td_in}:

Because this data-set is going into a \textit{TableLookupTransposeDataSet},
the data will be interpreted as though each row represents a coincident data
set, while each column represents the possible values for one specific
variable.

The manager is informed that the data for the independent variable is
contained within these data, located at column with index 1, i.e $\{1,11,21\}$
(S\_define:168,169)

The manager is provided with two configurations, resulting in the management
of two tables (S\_define:171-184)
\begin{enumerate}
  \item \textbf{var1} This configuration is given two dependent variables
  (\textit{ dependent[0], dependent[1]}) and instructed to take their data
  from columns 3 and 4 respectively (i.e. $\{3,13,23\}$
  and $\{4,14,24\}$.
  \item \textbf{var2}  This configuration is assigned to a pre-allocated
  table so the manager is not responsible for creating one.  This table has
  been pre-configured to populate \textit{dependent[2]}, and is now
  instructed to take its data from column 0, i.e. $\{0,10,20\}$.
\end{enumerate}

The independent variable starts with value 5 and is incremented to $30$
in steps of $5$.  Note -- above 21, the independent variable is
off-table-high so the dependent variables will take the last calibrated
value.

These values are all observed in the log files.  The
\textit{TableLookupTransposeDataSet} accurately transposes the input data.
The results are shown in Table \ref{tab:td_out}.


\begin{table}
\caption{Transpose Data Input}
\label{tab:td_in}
\begin{center}
\begin{tabular}{|rrrrr|}
\hline
0  &  1 &  2 &  3 &  4\\
10 & 11 & 12 & 13 & 14\\
20 & 21 & 22 & 23 & 24\\
\hline
\end{tabular}
\end{center}
\end{table}

\begin{table}
\caption{Results Table for Transpose Data}
\label{tab:td_out}
\begin{center}
\begin{tabular}{|r|r|r|r|r|} \hline
indep & indep index & dependent[0] & dependent[1] & dependent[2] \\
      & (col 1)     & (col 3)      & (col 4)      & (col 0) \\
\hline\hline
 5    & 0.4          & 3 + 4 = 7   & 4 + 4 = 8    & 0 + 4 = 4 \\ \hline
10    & 0.9          & 3 + 9 = 12  & 4 + 9 = 13   & 0 + 9 = 9 \\ \hline
15    & 1.4          & 13 + 4 = 17 & 14 + 4 = 18  & 10 + 4 = 14 \\ \hline
20    & 1.9          & 13 + 9 = 22 & 14 + 9 = 23  & 10 + 9 = 19 \\ \hline
25    & 2.0          & 23 + 0 = 23 & 24 + 0 = 24  & 20 + 0 = 20 \\ \hline
30    & 2.0          & 23 + 0 = 23 & 24 + 0 = 24  & 20 + 0 = 20 \\ \hline
\end{tabular}
\end{center}
\end{table}

Other tests in \textit{SIM\_06\_transpose\_data\_manager} and
\textit{SIM\_08\_transpose\_config} probe the detection and handling of
other configurations, such as:
\begin{itemize}
  \item reassigning which columns to use for which variables,
  \item using a pre-allocated \textit{TableIndependentVariable} instance,
  relieving the manager of the responsibility for creating one
  \item using inconsistent / irregular data files,
  \item initializing with incomplete configuration
  \item etc.
\end{itemize}

\section{Reloading the independent and dependent data}
\label{sec:UG_reloading_data}
Reloading the dependent data and independent data in mid-run is examined
in \textit{SIM\_07\_reload\_data}.  The Python input file
(\textit{SIM\_07\_reload\_data/RUN\_verif/input.py}) contains the initial
and subsequent data loads.  The initial part of the run consists of an
independent variable with seven data points.  There are two output
variables (dependent data), each with seven data points.  The model
performs as expected (see
\textit{SIM\_07\_reload\_data/verif\_data/RUN\_verif/log\_test\_data.csv}).
After several sets of output are generated, the dependent data tables are
re-loaded with new data using the \textit{load\_dependent\_variables()}
method.  See section \ref{sec:UG_dependent_variable_data}, Dependent
Variable Data.  At this point, two important points should be mentioned:
\begin{enumerate}
  \item Can use any of the four method signatures to re-load the dependent
    data (see \ref{subsec:dep_var_data}, Dependent Variable Data).
  \item Can change number of outputs but not the dimension.  Since the
    independent data is not being changed, the size of each dependent variable
    collection must not change.  For example, the test case increases the
    outputs from two to three, and reloads the data accordingly.  Each
    dependent variable still has seven values.
\end{enumerate}
After several more outputs are generated, the independent AND dependent data
are reloaded.  Now it is possible to change the dimension of the data sets,
since the independent variable is being redefined.  As far is the model is
concerned, the interpolation process has started over with new data.  The user
must be mindful of the value of the independent variable, and can change it as
desired.

The verification run above is now illustrated with data.  As stated, we have
one independent variable and two dependent variables with seven values for each.
\begin{table}[H]
\caption{Initial data load}
\begin{center}
\begin{tabular}{|r|r|r|r|r|} \hline
independent\_variable & dependent\_ variable[0] & dependent\_variable[1] \\
\hline\hline
0 & 10 & 100 \\ \hline
1 & 11 & 101 \\ \hline
2 & 12 & 102 \\ \hline
3 & 13 & 103 \\ \hline
4 & 14 & 104 \\ \hline
5 & 15 & 105 \\ \hline
6 & 16 & 106 \\ \hline
\end{tabular}
\end{center}
\end{table}

The simulation assigns a value to the independent variable,
and the table interpolation process produces the following expected results:
\begin{table}[H]
\caption{Interpolation results}
\begin{center}
\begin{tabular}{|r|r|r|r|r|} \hline
assigned & computed & computed \\
independent\_variable & dependent\_ variable[0] & dependent\_variable[1] \\
\hline\hline
1.2 & 10.2 & 100.2 \\ \hline
2.2 & 11.2 & 101.2 \\ \hline
3.2 & 12.2 & 102.2 \\ \hline
\end{tabular}
\end{center}
\end{table}

Next, the dependent data is reloaded  and the number of dependent variables
increased to three.  Note that the number of data points per variable
remains constant.

\begin{table}[H]
\caption{Dependent data reload}
\begin{center}
\begin{tabular}{|r|r|r|r|r|} \hline
reloaded & reloaded & reloaded \\
dependent\_variable[0] & dependent\_ variable[1] & dependent\_variable[2] \\
\hline\hline
10000 & -10000 & 5 \\ \hline
11000 & -10100 & 4 \\ \hline
12000 & -10200 & 3 \\ \hline
13000 & -10300 & 2 \\ \hline
14000 & -10400 & 1 \\ \hline
15000 & -10500 & 0 \\ \hline
16000 & -10600 & -1 \\ \hline
\end{tabular}
\end{center}
\end{table}

The simulation continues with new independent variable assignments,
and a modified value for the independent variable.
The table interpolation process produces the following expected results:
\begin{table}[H]
\caption{Interpolation results after dependent data reload}
\begin{center}
\begin{tabular}{|r|r|r|r|r|} \hline
assigned & computed & computed & computed  \\
independent\_variable & dependent\_ variable[0] & dependent\_variable[1]
 & dependent\_variable[2]\\
\hline\hline
2.7 & 12700 & -10270 & 2.3 \\ \hline
3.7 & 13700 & -10370 & 1.3 \\ \hline
4.7 & 14700 & -10470 & 0.3 \\ \hline
\end{tabular}
\end{center}
\end{table}

Finally, both the independent and dependent data are reloaded,
and the data size for each variable is changed.
\begin{table}[H]
\caption{Complete data reload}
\begin{center}
\begin{tabular}{|r|r|r|r|r|} \hline
independent & dependent[0] & dependent[1] & dependent[2] & dependent[3]\\
\hline\hline
10 & 10000 & -10000 & 5 & 40 \\ \hline
11 & 11000 & -10100 & 4 & 30\\ \hline
12 & 12000 & -10200 & 3 & 20\\ \hline
13 & 13000 & -10300 & 2 & 10\\ \hline
14 & 14000 & -10400 & 1 & 0\\ \hline
\end{tabular}
\end{center}
\end{table}

The simulation continues with new independent variable assignments,
and the table interpolation process produces the following expected results:
\begin{table}[H]
\caption{Interpolation results after complete data reload}
\begin{center}
\begin{tabular}{|r|r|r|r|r|} \hline
assigned & computed & computed & computed & computed \\
independent & dependent[0] & dependent[1] & dependent[2] & dependent[3]\\
\hline\hline
10.4 & 10400 & -10040 & 4.6 & 36 \\ \hline
11.4 & 11400 & -10140 & 3.6 & 26 \\ \hline
12.4 & 12400 & -10240 & 2.6 & 16 \\ \hline
13.4 & 13400 & -10340 & 1.6 &  6 \\ \hline
14.4 & 14000 & -10400 & 1   &  0 \\ \hline
\end{tabular}
\end{center}
\end{table}
Note that the independent value 14.4 is off-table so the dependent variables
take the value of the last entry in their respective tables.

The logged values (see \textit{RUN\_verif/log\_multi\_data.csv}) agree with the
predictions in the tables above. A screenshot of this log file is included
below.
\begin{figure}[!htb]
  \centering
  \includegraphics[width=\linewidth]{Images/Output.png}
  \caption{Reload data output}
\end{figure}



\section{Code coverage}
Code coverage for the entire nine modules is shown below:
\begin{figure}[!htb]
  \centering
  \includegraphics[width=\linewidth]{Images/CodeCoverage.png}
  \caption{Table Interpolation Code Coverage}
\end{figure}

Not all of the modules have 100\% coverage (code coverage for the header files
is not shown, as the results can be misleading).  This is to be expected for
this model.  The Table Interpolation model is designed to be expanded and as
such, potential error sources must be anticipated.  As a result, some error
reporting is not reachable with the current baseline configuration.  However,
by using a debugger and setting breakpoints at the appropriate locations,
certain variables may be modified so that the target code will be executed.
In all cases, the uncovered code was executed.  The steps needed to execute
these lines of "unreachable" code, along with the necessary conditions, are
listed below:

The non-covered lines are as follows:
\begin{enumerate}
 \item \textit{generic\_multi\_input\_table.cc}
   \begin{enumerate}
      \item \textit{initialize()}
      \begin{code}
if (independents[ii].first->get_size() <= 1) {
   if (num_data_elements_per_increment_of_index[ii] !=
      num_data_elements_per_increment_of_index[ii+1]) {
        MessageHandler::fail(
      \end{code}
      Set
      \textit{num\_data\_elements\_per\_increment\_of\_index}[0] = 0 to cover code.
      Requires that no data has been loaded, which is impossible in current
      framework.
      \item \textit{generate\_trivial\_output()}
      \begin{code}
if (!output_ptrs_set) {
    // NOTE this is a safety check,
    //     it may not be reachable.
    // This method is called from update,
    //  which requires initialization, which...
      \end{code}
      Set \textit{output\_ptrs\_set} = false to cover code.
Requires no dependent variables being set, which requires no initialization,
which is impossible in current framework.
      \item \textit{generate\_base\_values()}
      \begin{code}
if ((*indep_iter).first->get_size() <= 1) {
      \end{code}
      Set \textit{(*indep\_iter).first.size} = 1 to cover code.  Requires an
      independent variable of one point or less...which would have been
      stripped at this point.
      \item \textit{generate\_base\_values()}
      \begin{code}
if (dwell >= num_data_points_interp) {
      \end{code}
      Set \textit{dwell} = 0 to cover code.  Requires that \textit{dwell} be
      bigger than the number of interpolation points.  Will only be,
      at most, half the value.
      \item \textit{generate\_base\_values()}
      \begin{code}
default:
      // NOTE - this should be unreachable.  ...
      // value is protected.  This is a fail-safe.
      MessageHandler::fail(
      \end{code}
      Set interpolation method, \textit{((*indep\_iter).second)}
      \begin{code}
switch ((*indep_iter).second) {
      \end{code}
      to some integer greater than 5 so that the default case will be executed.
Requires independent variable to have a non-standard LookupMethod.  This is not
 possible in the current implementation, as the only way to specify the lookup
 method is with enumerated values.
      \item \label{item:code1} \textit{precheck\_output()}
      \begin{code}
if (!output_ptrs_set) {
      \end{code}
      Set \textit{output\_ptrs\_set} = false to cover code.  This routine is called
 from \textit{update()}, which can only be called after \textit{initialization()}.
  \textit{initialization()} can only succeed if the output pointers are defined,
  which is to say that the dependent variables have been identified.  No path
  exists for this condition.
      \item \label{item:code2} \textit{precheck\_output()}
      \begin{code}
if (data_point_index.size()!=data_point_weight.size()){
    MessageHandler::error(
      \end{code}
      Clear \textit{data\_point\_index} to cover code.  \textit{data\_point\_index}
       and \textit{data\_point\_weight} are vectors and are designed to be equal in
        size.  Only a deliberate change (which is not recommended) could cause this
         section to be executed.
      \item \label{item:code3} \textit{generate\_output()}
      \begin{code}
if (!precheck_output()) {
    return false;
  }
      \end{code}
      Will be executed only if \ref{item:code1}) or \ref{item:code2}) above return
       false...which is not possible in the current implementation.
   \end{enumerate}

 \item \textit{single\_input\_table\_for\_angles.cc}
   \begin{itemize}
      \item \textit{generate\_output()}
      \begin{code}
  if (!precheck_output()) {
    return false;
  }
      \end{code}
      This is the same check as in \ref{item:code3}) above.
      \textit{SingleInputTableForAngles} inherits from
       \textit{GenericMultiInputTable} (and has no \textit{precheck\_output()} of
       its own), so the same reasoning applies.
   \end{itemize}

 \item \textit{single\_input\_table\_for\_quaternions.cc}
   \begin{itemize}
      \item \textit{generate\_output()}
      \begin{code}
if (!precheck_output()) {
  return false;
}
      \end{code}
      This is the same check as in \ref{item:code3}) above.
       \textit{SingleInputTableForQuaternions} inherits from
        \textit{GenericMultiInputTable} (has no \textit{precheck\_output()}
         of its own), so the same reasoning applies.
   \end{itemize}

  \item \textit{table\_independent\_variable.cc}
   \begin{itemize}
      \item \textit{generate\_fraction()}
      \begin{code}
if (fraction < 0.0) {
  fraction = 0.0;
}
else if (fraction > 1.0) {
  fraction = 1.0;
}
      \end{code}
      Set \textit{fraction} $<$ 0 and \textit{fraction} $>$ 1.  No path (in the
      current form) exists for these two checks.  They would require an external
      modification of the independent data, or a data corruption.
   \end{itemize}

   \item \textit{table\_lookup\_\_transpose\_data.cc}
    \begin{itemize}
       \item \textit{populate\_table()}
       \begin{code}
if (table_ptr == NULL) {
  MessageHandler::fail(
      \end{code}
      Set \textit{table\_ptr} = NULL.  This check will never be reached.  The
      pointer to the table has been set in \textit{process\_data()}, which always
      occurs before \textit{populate\_table()}.
    \end{itemize}
\end{enumerate}

\end{document}
